\documentclass[10pt,a4paper]{article}
\usepackage[utf8]{inputenc}
\usepackage[T1]{fontenc}
\usepackage[spanish,es-nodecimaldot,es-noquoting,es-noshorthands]{babel}
\usepackage{lmodern}
\usepackage{amsmath,amssymb}
\usepackage[margin=2cm,headheight=14pt]{geometry}
\usepackage{booktabs,longtable,array}
\usepackage{adjustbox}
\usepackage[table]{xcolor}
\usepackage{fancyhdr}
\usepackage[hidelinks]{hyperref}
\definecolor{azul}{HTML}{1F3864}
\definecolor{verde}{HTML}{006100}
\definecolor{rojo}{HTML}{9C0006}
\definecolor{fondodato}{HTML}{FFFFCC}
\pagestyle{fancy}
\fancyhf{}
\fancyhead[L]{\small\textsc{Diseño de Puentes - 2P}}
\fancyhead[R]{\small\nouppercase{\leftmark}}
\fancyfoot[C]{\small\thepage}
\setlength{\parindent}{0pt}
\setcounter{secnumdepth}{1}
\allowdisplaybreaks
\newcommand{\ok}[1]{\textcolor{verde}{\textbf{#1}}}
\newcommand{\nook}[1]{\textcolor{rojo}{\textbf{#1}}}
\newcommand{\info}[1]{\textcolor{azul}{\textbf{#1}}}
\newcommand{\eqitem}[3]{\par\noindent\begin{minipage}{\linewidth}{\small #1}\hfill{\scriptsize\textcolor{gray}{#2}}\\[2pt]
  \hspace*{1.5em}\adjustbox{max width=\dimexpr\linewidth-1.5em\relax}{$\displaystyle #3$}\end{minipage}\par\smallskip}
\newcommand{\nota}[1]{\par\noindent{\small\itshape\textcolor{black!70}{#1}}\par\smallskip}

\title{\color{azul}\textbf{Diseño de Puentes - 2P}\\[4pt]\large Puente tipo losa-vigas preforzadas: superestructura, pila, viga cabezal y estribo}
\author{Memoria de cálculo automatizada}
\date{30/09/2026}
\begin{document}
\maketitle
\tableofcontents
\clearpage
\section*{Introducción}
\addcontentsline{toc}{section}{Introducción}
\markboth{Introducción}{Introducción}
Esta memoria reproduce, de forma automatizada, los cálculos del documento \emph{Puentes\_2P\_calculos.pdf} (puente tipo losa-vigas preforzadas). Cada ecuación se define una sola vez en Python; a partir de esa definición se generan (1) los valores de este reporte, (2) las fórmulas vivas del libro Excel \texttt{Puentes\_2P.xlsx} y (3) las ecuaciones que se muestran aquí.

Los cálculos internos se realizan en tonf y m; las unidades (cm, kip, ksi, psi, MPa, ...) actúan como factores de conversión, igual que en Mathcad. Los datos de entrada se muestran con fondo amarillo. Las verificaciones se indican en \ok{verde} (cumple) o \nook{rojo} (no cumple); las selecciones informativas en \info{azul}.

Las cargas vivas móviles (camión HS20, camión de fatiga) se calculan con líneas de influencia de la viga simplemente apoyada, colocando cada eje sobre la sección de análisis en ambos sentidos de circulación (envolvente exacta).
\clearpage

\section*{Resumen de resultados}
\addcontentsline{toc}{section}{Resumen de resultados}
\markboth{Resumen de resultados}{Resumen de resultados}
\begin{longtable}{@{}l p{0.55\linewidth} l r@{}}
\toprule
\textbf{Mód.} & \textbf{Resultado} & \textbf{Símbolo} & \textbf{Valor}\\
\midrule
\endhead
STD & Número de torones seleccionado & $N_{\mathrm{t}}$ & $44$ \\
STD & Porcentaje de pérdidas totales & $p_{\mathrm{T}}$ & $0.2510$ \\
STD & Fuerza de pretensado efectiva & $P_{se,c}$ & $463.52$$\,\mathrm{tonf}$ \\
STD & Fuerza de sujeción total & $F_{\mathrm{sujtot}}$ & $14.651$$\,\mathrm{tonf}$ \\
STD & Momento último & $M_{\mathrm{u}}$ & $1\,205.19$$\,\mathrm{tonf\cdot m}$ \\
STD & Resistencia a flexión mayorada & $M_{\mathrm{r}}$ & $1\,473.80$$\,\mathrm{tonf\cdot m}$ \\
STD & Cortante mayorado & $V_{\mathrm{u}}$ & $129.46$$\,\mathrm{tonf}$ \\
STD & Total en el izaje, factores PCI ($\uparrow$) & $\Delta_{pg,izaje}$ & $10.188$$\,\mathrm{cm}$ \\
STD & Deflexión por CV + impacto & $\Delta_{LT}$ & $2.081$$\,\mathrm{cm}$ \\
LRFD & Factor de distribución para momento (controla) & $\mathit{FDM}_{\mathrm{imax}}$ & $0.7329$ \\
LRFD & Factor de distribución para cortante (controla) & $\mathit{FDV}_{\mathrm{imax}}$ & $0.8854$ \\
LRFD & Número de torones adoptado (2a iteración) & $N_{\mathrm{t}}$ & $48$ \\
LRFD & Fuerza total de pretensado después de pérdidas & $F_T$ & $495.56$$\,\mathrm{tonf}$ \\
LRFD & \% de pérdida final & $p_{\mathrm{t}}$ & $26.595$ \\
LRFD & Fuerza de sujeción total & $F_{\mathrm{sujtot}}$ & $16.316$$\,\mathrm{tonf}$ \\
LRFD & Momento de diseño (Resistencia I) & $M_{\mathrm{u}}$ & $1\,283.11$$\,\mathrm{tonf\cdot m}$ \\
LRFD & Resistencia a flexión mayorada & $M_{\mathrm{r}}$ & $1\,571.77$$\,\mathrm{tonf\cdot m}$ \\
LRFD & Cortante mayorado (sección 0.06L) & $V_{\mathrm{u1}}$ & $142.95$$\,\mathrm{tonf}$ \\
LRFD & Total en el izaje, factores PCI ($\uparrow$) & $\Delta_{pg,izaje}$ & $10.614$$\,\mathrm{cm}$ \\
LRFD & Deflexión por carga viva & $\Delta_{mayor}$ & $2.017$$\,\mathrm{cm}$ \\
PAC & Peso total & $W_{\mathrm{total}}$ & $246.86$$\,\mathrm{tonf}$ \\
PAC & Peso aferente a la pila & $W$ & $0.8229$$\,\mathrm{tonf/m^{2}}$ \\
PPR & Peso total & $W_{\mathrm{total}}$ & $382.05$$\,\mathrm{tonf}$ \\
PPR & Peso aferente a la pila & $W$ & $1.273$$\,\mathrm{tonf/m^{2}}$ \\
VC & Separación escogida del refuerzo transversal & $S_{\mathrm{esc}}$ & $30$$\,\mathrm{cm}$ \\
VC & Momento total & $M_{\mathrm{t}}$ & $169\,437.67$$\,\mathrm{kgf\cdot m}$ \\
VC & Número de varillas a emplear & $N_{\mathrm{var,c}}$ & $15$ \\
FR & Fuerza de frenado & $\mathit{BR}$ & $16\,329.33$$\,\mathrm{kgf}$ \\
RE & Reacción total sobre el estribo & $\mathit{RA}_{\mathrm{2carril}}$ & $107\,201.57$$\,\mathrm{kgf}$ \\
EST & Fuerza sísmica de diseño & $P_{seis}$ & $15.436$$\,\mathrm{tonf/m}$ \\
EST & Momento último de la pantalla & $M_U$ & $88.488$$\,\mathrm{tonf\cdot m/m}$ \\
EST & Separación adoptada (emplear $\varnothing$25 mm @ 14 cm) & $S_{asumido}$ & $14$$\,\mathrm{cm}$ \\
EST & Separación adoptada (emplear $\varnothing$20 mm @ 35 cm) & $S_{asumido,p,z}$ & $35$$\,\mathrm{cm}$ \\
EST & Separación adoptada (emplear $\varnothing$20 mm @ 20 cm) & $S_{asumido,t,z}$ & $20$$\,\mathrm{cm}$ \\
\bottomrule
\end{longtable}
\subsection*{Verificaciones}
\nota{Total de verificaciones: 116; no cumplen: 2.}
\begin{longtable}{@{}l p{0.7\linewidth} c@{}}
\toprule
\textbf{Mód.} & \textbf{Verificación} & \textbf{Resultado}\\
\midrule
\endhead
STD & Impacto no mayor a 0.3 & \ok{OK} \\
STD & Torones seleccionados $\ge$ requeridos & \ok{OK} \\
STD & La tabla de distribución contiene Nt torones & \ok{OK} \\
STD & Torones desviados coinciden con la tabla del extremo & \ok{OK} \\
STD & Fibra superior (con torones desviados) & \ok{OK} \\
STD & Fibra inferior (con torones desviados) & \ok{OK} \\
STD & Fibra superior en punto harp & \ok{OK} \\
STD & Fibra inferior en punto harp & \ok{OK} \\
STD & Fibra superior en centro del tramo & \ok{OK} \\
STD & Fibra inferior en centro del tramo & \ok{OK} \\
STD & Fibra superior viga - Caso I & \ok{OK} \\
STD & Fibra superior viga - Caso II & \ok{OK} \\
STD & Fibra superior viga - Caso III & \ok{OK} \\
STD & Fibra inferior viga - compresión & \ok{OK} \\
STD & Fibra inferior viga - tensión ($\ge$ -6$\surd$f'c) & \ok{OK} \\
STD & Losa - Caso I & \ok{OK} \\
STD & Losa - Caso II & \ok{OK} \\
STD & Losa - Caso III & \ok{OK} \\
STD & Bloque dentro del ala (sección rectangular) & \ok{OK} \\
STD & Resistencia a flexión Mr > Mu & \ok{OK} \\
STD & Refuerzo máximo & \ok{OK} \\
STD & Refuerzo mínimo: $\phi$Mn $\ge$ 1.2 Mcr & \ok{OK} \\
STD & Vci $\ge$ Vci mín & \ok{OK} \\
STD & Vs $\le$ Vs máx & \ok{OK} \\
STD & Separación $\le$ máxima & \ok{OK} \\
STD & Estribo proporcionado $\ge$ requerido & \ok{OK} \\
STD & Cortante horizontal - Caso (c) & \ok{OK} \\
STD & Refuerzo proporcionado $\ge$ Avh mínimo & \nook{NO OK} \\
STD & Separación $\le$ máxima (corte horizontal) & \ok{OK} \\
STD & Refuerzo en zona de anclaje & \ok{OK} \\
STD & Deflexión por carga viva $\le$ L/800 & \ok{OK} \\
LRFD & Número de vigas Nb > 4 & \ok{OK} \\
LRFD & Calzada en el voladizo $\le$ 3.0 ft & \ok{OK} \\
LRFD & (a) 3.5 $\le$ S $\le$ 16.0 ft & \ok{OK} \\
LRFD & (b) 4.5 $\le$ ts $\le$ 12.0 in & \ok{OK} \\
LRFD & (c) 20 $\le$ L $\le$ 240 ft & \ok{OK} \\
LRFD & (d) Nb $\ge$ 4 & \ok{OK} \\
LRFD & (e) 10 000 $\le$ Kg $\le$ 7 000 000 in$^{4}$ & \ok{OK} \\
LRFD & Torones adoptados $\ge$ requeridos & \ok{OK} \\
LRFD & La tabla de distribución contiene Nt torones & \ok{OK} \\
LRFD & Pérdida real $\approx$ asumida ($\pm$1 \%): no se requiere iterar & \ok{OK} \\
LRFD & fpe $\le$ 0.8 fpy & \ok{OK} \\
LRFD & Rectos + desviados = Nt & \ok{OK} \\
LRFD & Fibra superior (con torones desviados) & \ok{OK} \\
LRFD & Fibra inferior (con torones desviados) & \ok{OK} \\
LRFD & Fibra superior en punto harp & \ok{OK} \\
LRFD & Fibra inferior en punto harp & \ok{OK} \\
LRFD & Fibra superior en centro del tramo & \ok{OK} \\
LRFD & Fibra inferior en centro del tramo & \ok{OK} \\
LRFD & Fibra superior viga - permanentes & \ok{OK} \\
LRFD & Fibra superior viga - todas las cargas & \ok{OK} \\
LRFD & Losa - permanentes & \ok{OK} \\
LRFD & Losa - todas las cargas & \ok{OK} \\
LRFD & Fibra inferior viga - tensión Servicio III & \ok{OK} \\
LRFD & Fatiga: 2(0.75 Mf/Sbc) < f compresión & \ok{OK} \\
LRFD & Comportamiento de sección rectangular (c < ts) & \ok{OK} \\
LRFD & fpe $\ge$ 0.5 fpu & \ok{OK} \\
LRFD & Resistencia a flexión Mr > Mu & \ok{OK} \\
LRFD & Máximo refuerzo: c/de $\le$ 0.42 & \ok{OK} \\
LRFD & Refuerzo mínimo & \ok{OK} \\
LRFD & dv $\ge$ 0.9 de & \ok{OK} \\
LRFD & dv $\ge$ 0.72 h & \ok{OK} \\
LRFD & vu < 0.125 f'c & \ok{OK} \\
LRFD & Separación $\le$ máxima & \ok{OK} \\
LRFD & Estribo proporcionado $\ge$ requerido & \ok{OK} \\
LRFD & Refuerzo $\ge$ mínimo & \ok{OK} \\
LRFD & Vc + Vs $\le$ 0.25 f'c bv dv & \ok{OK} \\
LRFD & Vn/Acv < 0.1 ksi: se puede despreciar el mínimo & \ok{OK} \\
LRFD & Vni $\le$ K1 f'c Acv & \ok{OK} \\
LRFD & Vni $\le$ K2 Acv & \ok{OK} \\
LRFD & Fuerza disponible > V (criterio del PDF) & \ok{CUMPLE} \\
LRFD & Refuerzo en zona de anclaje & \ok{OK} \\
LRFD & Deflexión por carga viva $\le$ L/800 & \ok{OK} \\
PAC & Distribución transversal = ancho del tablero & \ok{OK} \\
PAC & Espesor asumido $\ge$ mínimo & \ok{OK} \\
PPR & Distribución transversal = ancho del tablero & \ok{OK} \\
PPR & Espesor asumido $\ge$ mínimo & \ok{OK} \\
VC & Ancho asumido $\ge$ ancho mínimo & \ok{OK} \\
VC & (v - vh) $\le$ límite 1 & \ok{CUMPLE} \\
VC & (v - vh) $\le$ límite 2 & \ok{CUMPLE} \\
VC & Separación escogida $\le$ requerida & \ok{CUMPLE} \\
VC & Acero proporcionado $\ge$ requerido & \ok{CUMPLE} \\
VC & Ancho mínimo < ancho de viga & \ok{CUMPLE} \\
VC & Refuerzo de temperatura & \ok{CUMPLE} \\
VC & Separación escogida $\le$ máxima & \ok{CUMPLE} \\
VC & 2 Ask < 0.5 AsT & \ok{CUMPLE} \\
VC & Refuerzo superficial & \ok{CUMPLE} \\
VC & Acero superficial proporcionado $\ge$ requerido & \ok{CUMPLE} \\
EST & $\phi$ $\ge$ i + $\theta$MO (Ec. 11.6.5.3-1) & \ok{CUMPLE} \\
EST & Volcamiento - Resistencia I - Máx. & \ok{OK} \\
EST & Volcamiento - Resistencia I - Mín. & \ok{OK} \\
EST & Volcamiento - Evento Extremo I - Máx. & \ok{OK} \\
EST & Volcamiento - Evento Extremo I - Mín. & \ok{OK} \\
EST & Esfuerzo sobre el suelo - Resistencia I - Máx. & \ok{OK} \\
EST & Esfuerzo sobre el suelo - Resistencia I - Mín. & \ok{OK} \\
EST & Esfuerzo sobre el suelo - Evento Extremo I - Máx. & \ok{OK} \\
EST & Esfuerzo sobre el suelo - Evento Extremo I - Mín. & \ok{OK} \\
EST & Esfuerzo sobre el suelo - Servicio I & \ok{OK} \\
EST & Deslizamiento - Resistencia I - Máx. & \ok{OK} \\
EST & Deslizamiento - Resistencia I - Mín. & \ok{OK} \\
EST & Deslizamiento - Evento Extremo I - Máx. & \ok{OK} \\
EST & Deslizamiento - Evento Extremo I - Mín. & \nook{NO OK} \\
EST & $\varepsilon$t $\ge$ 0.005 (tensión controlada) - pantalla & \ok{OK} \\
EST & Acero proporcionado $\ge$ requerido - pantalla & \ok{OK} \\
EST & Capacidad $\ge$ momento último - pantalla & \ok{CUMPLE} \\
EST & As mínimo - pantalla & \ok{CUMPLE} \\
EST & 0.233 $\le$ As $\le$ 1.27 mm$^{2}$/mm - pantalla & \ok{OK} \\
EST & $\varepsilon$t $\ge$ 0.005 - punta & \ok{OK} \\
EST & Acero proporcionado $\ge$ requerido - punta & \ok{CUMPLE} \\
EST & Capacidad de momento - punta & \ok{CUMPLE} \\
EST & 3 dv > l8 & \ok{OK} \\
EST & Cortante en la punta & \ok{OK} \\
EST & $\varepsilon$t $\ge$ 0.005 - talón & \ok{OK} \\
EST & Acero proporcionado $\ge$ requerido - talón & \ok{CUMPLE} \\
EST & Capacidad de momento - talón & \ok{CUMPLE} \\
EST & 0.233 $\le$ As $\le$ 1.27 mm$^{2}$/mm - zapata & \ok{OK} \\
\bottomrule
\end{longtable}
\clearpage

\section{Viga interior preforzada - AASHTO Estándar}

\nota{Puente tipo losa-vigas preforzadas. Viga AASHTO-PCI Bulb-Tee BT-72. Diseño según especificaciones AASHTO Estándar (carga HS20).}

\subsection{2. Materiales}

\begin{longtable}{@{}l p{0.52\linewidth} r l l@{}}
\toprule
\textbf{Dato} & \textbf{Descripción} & \textbf{Valor} & \textbf{Unidad} & \textbf{Ref.}\\
\midrule
$t_{\mathrm{s}}$ & Espesor estructural de la losa & \cellcolor{fondodato}19 & $\mathrm{cm}$ & {\scriptsize } \\
$t_{\mathrm{ws}}$ & Sobre-espesor de desgaste & \cellcolor{fondodato}1 & $\mathrm{cm}$ & {\scriptsize } \\
\bottomrule
\end{longtable}

\eqitem{Espesor total de la losa}{}{t = t_{\mathrm{s}} + t_{\mathrm{ws}} = \mathbf{20}\,\mathrm{cm}}

\begin{longtable}{@{}l p{0.52\linewidth} r l l@{}}
\toprule
\textbf{Dato} & \textbf{Descripción} & \textbf{Valor} & \textbf{Unidad} & \textbf{Ref.}\\
\midrule
$f'_{c,s}$ & f'c de diseño de la losa & \cellcolor{fondodato}280 & $\mathrm{kgf/cm^{2}}$ & {\scriptsize } \\
$f'_{c,g}$ & f'c de diseño de la viga (BT-72) & \cellcolor{fondodato}450 & $\mathrm{kgf/cm^{2}}$ & {\scriptsize } \\
$f'_{ci}$ & f'c de la viga en la transferencia & \cellcolor{fondodato}380 & $\mathrm{kgf/cm^{2}}$ & {\scriptsize } \\
$W_{\mathrm{c}}$ & Peso unitario del hormigón & \cellcolor{fondodato}2.400 & $\mathrm{tonf/m^{3}}$ & {\scriptsize STD A. 3.3.6} \\
$L_{\mathrm{2}}$ & Longitud total de la viga & \cellcolor{fondodato}36.900 & $\mathrm{m}$ & {\scriptsize } \\
$L_{\mathrm{1}}$ & Tramo de diseño (entre apoyos) & \cellcolor{fondodato}36.600 & $\mathrm{m}$ & {\scriptsize } \\
$d_t$ & Diámetro del torón (7 alambres, baja relajación) & \cellcolor{fondodato}12.700 & $\mathrm{mm}$ & {\scriptsize } \\
$A_{\mathrm{st}}$ & Área de un torón & \cellcolor{fondodato}98.700 & $\mathrm{mm^{2}}$ & {\scriptsize } \\
$f'_s$ & Esfuerzo último del torón f's & \cellcolor{fondodato}19 & $\mathrm{tonf/cm^{2}}$ & {\scriptsize } \\
\bottomrule
\end{longtable}

\eqitem{Resistencia a la fluencia}{STD A. 9.1.2}{f_{\mathrm{py}} = 0.9 \cdot f'_s = \mathbf{17.100}\,\mathrm{tonf/cm^{2}}}

\eqitem{Pretensado inicial}{STD A. 9.15.1}{f_{\mathrm{si}} = 0.75 \cdot f'_s = \mathbf{14.250}\,\mathrm{tonf/cm^{2}}}

\begin{longtable}{@{}l p{0.52\linewidth} r l l@{}}
\toprule
\textbf{Dato} & \textbf{Descripción} & \textbf{Valor} & \textbf{Unidad} & \textbf{Ref.}\\
\midrule
$E_{\mathrm{p}}$ & Módulo de elasticidad del torón & \cellcolor{fondodato}2\,003.75 & $\mathrm{tonf/cm^{2}}$ & {\scriptsize STD 9.16.2.1.2} \\
$f_{\mathrm{y}}$ & Fluencia de barras de refuerzo & \cellcolor{fondodato}4.200 & $\mathrm{tonf/cm^{2}}$ & {\scriptsize } \\
$f_{\mathrm{ws}}$ & Espesor de la capa de rodadura & \cellcolor{fondodato}5 & $\mathrm{cm}$ & {\scriptsize } \\
$W_{c,ha}$ & Peso unitario del hormigón asfáltico & \cellcolor{fondodato}2.400 & $\mathrm{tonf/m^{3}}$ & {\scriptsize } \\
$W_{\mathrm{b}}$ & Peso de la barrera New Jersey (por lado) & \cellcolor{fondodato}0.4465 & $\mathrm{tonf/m}$ & {\scriptsize } \\
\bottomrule
\end{longtable}

\subsection{3. Propiedades de la sección transversal de una viga interior}

\subsubsection{3.1 Sección no compuesta}

\begin{longtable}{@{}l p{0.52\linewidth} r l l@{}}
\toprule
\textbf{Dato} & \textbf{Descripción} & \textbf{Valor} & \textbf{Unidad} & \textbf{Ref.}\\
\midrule
$A_{\mathrm{g}}$ & Área de la viga & \cellcolor{fondodato}0.4948 & $\mathrm{m^{2}}$ & {\scriptsize } \\
$h_{\mathrm{g}}$ & Altura de la viga & \cellcolor{fondodato}182.88 & $\mathrm{cm}$ & {\scriptsize } \\
$I_{\mathrm{g}}$ & Inercia de la viga & \cellcolor{fondodato}0.2272 & $\mathrm{m^{4}}$ & {\scriptsize } \\
$y_{\mathrm{b}}$ & Distancia del CG a la fibra inferior & \cellcolor{fondodato}92.964 & $\mathrm{cm}$ & {\scriptsize } \\
\bottomrule
\end{longtable}

\eqitem{Distancia del CG a la fibra superior}{}{y_{\mathrm{t}} = h_{\mathrm{g}} - y_{\mathrm{b}} = \mathbf{89.916}\,\mathrm{cm}}

\eqitem{Módulo de sección, fibra inferior}{}{S_{\mathrm{b}} = \dfrac{I_{\mathrm{g}}}{y_{\mathrm{b}}} = \mathbf{0.2444}\,\mathrm{m^{3}}}

\eqitem{Módulo de sección, fibra superior}{}{S_{\mathrm{t}} = \dfrac{I_{\mathrm{g}}}{y_{\mathrm{t}}} = \mathbf{0.2527}\,\mathrm{m^{3}}}

\eqitem{Peso propio de la viga}{}{W_{\mathrm{g}} = A_{\mathrm{g}} \cdot W_{\mathrm{c}} = \mathbf{1.188}\,\mathrm{tonf/m}}

\nota{Módulo de elasticidad: Ec = 33000$\cdot$Wc\textasciicircum{}1.5$\cdot$$\surd$f'c  (Wc en kip/ft$^{3}$, f'c en ksi) [STD A. 8.7.1]}

\eqitem{M.E. de la losa}{STD A. 8.7.1}{E_{\mathrm{c}} = 33000 \cdot {\left(\dfrac{W_{\mathrm{c}}}{\dfrac{\mathrm{kip}}{{\mathrm{ft}}^{3}}}\right)}^{1.5} \cdot \sqrt{\dfrac{f'_{c,s}}{\mathrm{ksi}}} \cdot \mathrm{ksi} = \mathbf{268.52}\,\mathrm{tonf/cm^{2}}}

\eqitem{M.E. de la viga en la transferencia}{}{E_{\mathrm{ci}} = 33000 \cdot {\left(\dfrac{W_{\mathrm{c}}}{\dfrac{\mathrm{kip}}{{\mathrm{ft}}^{3}}}\right)}^{1.5} \cdot \sqrt{\dfrac{f'_{ci}}{\mathrm{ksi}}} \cdot \mathrm{ksi} = \mathbf{312.82}\,\mathrm{tonf/cm^{2}}}

\eqitem{M.E. de la viga (cargas de servicio)}{}{E_{\mathrm{cs}} = 33000 \cdot {\left(\dfrac{W_{\mathrm{c}}}{\dfrac{\mathrm{kip}}{{\mathrm{ft}}^{3}}}\right)}^{1.5} \cdot \sqrt{\dfrac{f'_{c,g}}{\mathrm{ksi}}} \cdot \mathrm{ksi} = \mathbf{340.41}\,\mathrm{tonf/cm^{2}}}

\subsubsection{3.2 Sección compuesta}

\nota{Ancho de alma efectivo de la viga: el menor de be1 (ancho del ala superior) y be2 = 6(2)(h1+h2)+tw+2(h3). [STD A. 9.8.3.1]}

\begin{longtable}{@{}l p{0.52\linewidth} r l l@{}}
\toprule
\textbf{Dato} & \textbf{Descripción} & \textbf{Valor} & \textbf{Unidad} & \textbf{Ref.}\\
\midrule
$h_{\mathrm{tf1}}$ & Dimensión del ala h1 & \cellcolor{fondodato}8.890 & $\mathrm{cm}$ & {\scriptsize } \\
$h_{\mathrm{tf2}}$ & Dimensión del ala h2 & \cellcolor{fondodato}5.080 & $\mathrm{cm}$ & {\scriptsize } \\
$h_{\mathrm{tf3}}$ & Dimensión del ala h3 & \cellcolor{fondodato}5.080 & $\mathrm{cm}$ & {\scriptsize } \\
$t_{\mathrm{w}}$ & Espesor del alma de la viga & \cellcolor{fondodato}15.240 & $\mathrm{cm}$ & {\scriptsize } \\
$b_{\mathrm{ew1}}$ & Ancho del ala superior de la viga & \cellcolor{fondodato}106.68 & $\mathrm{cm}$ & {\scriptsize } \\
\bottomrule
\end{longtable}

\eqitem{Ancho be2}{}{b_{\mathrm{ew2}} = 6 \cdot 2 \cdot \left(h_{\mathrm{tf1}} + h_{\mathrm{tf2}}\right) + t_{\mathrm{w}} + 2 \cdot h_{\mathrm{tf3}} = \mathbf{193.04}\,\mathrm{cm}}

\eqitem{Ancho de alma efectivo}{}{b_{\mathrm{ew}} = \min\left(b_{\mathrm{ew1}},\ b_{\mathrm{ew2}}\right) = \mathbf{106.68}\,\mathrm{cm}}

\nota{Ancho de ala efectivo de la sección compuesta: el menor de L/4, S y 6(2)ts+bew. [STD A. 9.8.3.2]}

\begin{longtable}{@{}l p{0.52\linewidth} r l l@{}}
\toprule
\textbf{Dato} & \textbf{Descripción} & \textbf{Valor} & \textbf{Unidad} & \textbf{Ref.}\\
\midrule
$S$ & Separación entre vigas & \cellcolor{fondodato}275 & $\mathrm{cm}$ & {\scriptsize } \\
\bottomrule
\end{longtable}

\eqitem{1/4 del tramo de diseño}{}{b_{\mathrm{ea1}} = 0.25 \cdot L_{\mathrm{1}} = \mathbf{915}\,\mathrm{cm}}

\eqitem{Separación entre vigas}{}{b_{\mathrm{ea2}} = S = \mathbf{275}\,\mathrm{cm}}

\eqitem{6(2)ts + bew}{}{b_{\mathrm{ea3}} = 6 \cdot 2 \cdot t_{\mathrm{s}} + b_{\mathrm{ew}} = \mathbf{334.68}\,\mathrm{cm}}

\eqitem{Ancho de ala efectivo}{}{w_{\mathrm{fe}} = \min\left(b_{\mathrm{ea1}},\ b_{\mathrm{ea2}},\ b_{\mathrm{ea3}}\right) = \mathbf{275}\,\mathrm{cm}}

\eqitem{Razón modular losa/viga}{}{n_{dg} = \dfrac{E_{\mathrm{c}}}{E_{\mathrm{cs}}} = \mathbf{0.7888}}

\subsubsection{3.2.3 Propiedades de la sección transformada}

\eqitem{Ancho de ala transformado}{}{w_{\mathrm{ft}} = n_{dg} \cdot w_{\mathrm{fe}} = \mathbf{2.169}\,\mathrm{m}}

\eqitem{Área del ala transformada}{}{a_{\mathrm{ft}} = w_{\mathrm{ft}} \cdot t_{\mathrm{s}} = \mathbf{0.4122}\,\mathrm{m^{2}}}

\eqitem{Inercia propia del ala}{}{I_{\mathrm{ft}} = \dfrac{w_{\mathrm{ft}} \cdot {t_{\mathrm{s}}}^{3}}{12} = \mathbf{1.240\times10^{-3}}\,\mathrm{m^{4}}}

\begin{longtable}{@{}l p{0.52\linewidth} r l l@{}}
\toprule
\textbf{Dato} & \textbf{Descripción} & \textbf{Valor} & \textbf{Unidad} & \textbf{Ref.}\\
\midrule
$t_{\mathrm{h}}$ & Espesor mínimo de cartela (centro del tramo) & \cellcolor{fondodato}1 & $\mathrm{cm}$ & {\scriptsize } \\
\bottomrule
\end{longtable}

\eqitem{Ancho de cartela transformado}{}{w_{\mathrm{th}} = n_{dg} \cdot b_{\mathrm{ew1}} = \mathbf{0.8415}\,\mathrm{m}}

\eqitem{Área de cartela transformada}{}{a_{\mathrm{th}} = w_{\mathrm{th}} \cdot t_{\mathrm{h}} = \mathbf{8.415\times10^{-3}}\,\mathrm{m^{2}}}

\eqitem{Inercia propia de la cartela}{}{I_{\mathrm{at}} = \dfrac{w_{\mathrm{th}} \cdot {t_{\mathrm{h}}}^{3}}{12} = \mathbf{7.013}\,\mathrm{cm^{4}}}

\begin{center}
{\small\textbf{Propiedades de la sección compuesta (parte 1)}}\\[3pt]
\adjustbox{max width=\linewidth}{\begin{tabular}{ccccc}
\toprule
\begin{tabular}[b]{@{}c@{}}\textbf{Elemento}\end{tabular} & \begin{tabular}[b]{@{}c@{}}\textbf{Área}\\ {\scriptsize $\mathrm{m^{2}}$}\end{tabular} & \begin{tabular}[b]{@{}c@{}}\textbf{yb}\\ {\scriptsize $\mathrm{m}$}\end{tabular} & \begin{tabular}[b]{@{}c@{}}\textbf{A$\cdot$yb}\\ {\scriptsize $\mathrm{m^{3}}$}\end{tabular} & \begin{tabular}[b]{@{}c@{}}\textbf{I propia}\\ {\scriptsize $\mathrm{m^{4}}$}\end{tabular} \\
\midrule
Viga & $0.4948$ & $0.9296$ & $0.4600$ & $0.2272$ \\
Cartela & $8.415\times10^{-3}$ & $1.834$ & $0.0154$ & $7.013\times10^{-8}$ \\
Losa & $0.4122$ & $1.934$ & $0.7970$ & $1.240\times10^{-3}$ \\
\bottomrule
\end{tabular}}
\end{center}

\eqitem{Área total de la sección compuesta}{}{A_{\mathrm{c}} = A_{\mathrm{v}} + A_{\mathrm{h}} + A_{\mathrm{l}} = \mathbf{0.9154}\,\mathrm{m^{2}}}

\eqitem{CG de la sección compuesta - fibra inferior de la viga}{}{y_{\mathrm{bc}} = \dfrac{A_{\mathrm{y,v}} + A_{\mathrm{y,h}} + A_{\mathrm{y,l}}}{A_{\mathrm{c}}} = \mathbf{1.390}\,\mathrm{m}}

\begin{center}
{\small\textbf{Propiedades de la sección compuesta (parte 2: Steiner)}}\\[3pt]
\adjustbox{max width=\linewidth}{\begin{tabular}{ccc}
\toprule
\begin{tabular}[b]{@{}c@{}}\textbf{Elemento}\end{tabular} & \begin{tabular}[b]{@{}c@{}}\textbf{A(ybc - yb)$^{2}$}\\ {\scriptsize $\mathrm{m^{4}}$}\end{tabular} & \begin{tabular}[b]{@{}c@{}}\textbf{I + A(ybc - yb)$^{2}$}\\ {\scriptsize $\mathrm{m^{4}}$}\end{tabular} \\
\midrule
Viga & $0.1049$ & $0.3321$ \\
Cartela & $1.657\times10^{-3}$ & $1.657\times10^{-3}$ \\
Losa & $0.1218$ & $0.1231$ \\
\bottomrule
\end{tabular}}
\end{center}

\eqitem{Inercia de la sección compuesta}{}{I_{\mathrm{c}} = I_{\mathrm{t,v}} + I_{\mathrm{t,h}} + I_{\mathrm{t,l}} = \mathbf{0.4568}\,\mathrm{m^{4}}}

\eqitem{Altura total de la sección compuesta}{}{h_{\mathrm{sc}} = h_{\mathrm{g}} + t_{\mathrm{h}} + t_{\mathrm{s}} = \mathbf{2.029}\,\mathrm{m}}

\eqitem{CG - fibra superior de la viga}{}{y_{\mathrm{tg}} = h_{\mathrm{g}} - y_{\mathrm{bc}} = \mathbf{43.872}\,\mathrm{cm}}

\eqitem{CG - fibra superior de la losa}{}{y_{\mathrm{tc}} = h_{\mathrm{sc}} - y_{\mathrm{bc}} = \mathbf{63.872}\,\mathrm{cm}}

\eqitem{Módulo de sección compuesta, fibra inferior viga}{}{S_{\mathrm{bc}} = \dfrac{I_{\mathrm{c}}}{y_{\mathrm{bc}}} = \mathbf{0.3286}\,\mathrm{m^{3}}}

\eqitem{Módulo de sección compuesta, fibra superior viga}{}{S_{\mathrm{tg}} = \dfrac{I_{\mathrm{c}}}{y_{\mathrm{tg}}} = \mathbf{1.041}\,\mathrm{m^{3}}}

\eqitem{Módulo de sección compuesta, fibra superior losa}{}{S_{\mathrm{tc}} = \dfrac{1}{n_{dg}} \cdot \dfrac{I_{\mathrm{c}}}{y_{\mathrm{tc}}} = \mathbf{0.9067}\,\mathrm{m^{3}}}

\subsection{4. Fuerzas cortantes y momentos flectores}

\subsubsection{4.1 Cargas muertas (CM) [STD A. 3.3]}

\nota{Cargas sobre la sección no compuesta: viga, losa y cartela.}

\eqitem{Peso de la losa}{}{W_{\mathrm{d}} = t \cdot S \cdot W_{\mathrm{c}} = \mathbf{1.320}\,\mathrm{tonf/m}}

\eqitem{Peso de la cartela}{}{W_{\mathrm{ha}} = t_{\mathrm{h}} \cdot b_{\mathrm{ew1}} \cdot W_{\mathrm{c}} = \mathbf{0.0256}\,\mathrm{tonf/m}}

\eqitem{Peso de losa + cartela}{}{W_{\mathrm{s}} = W_{\mathrm{d}} + W_{\mathrm{ha}} = \mathbf{1.346}\,\mathrm{tonf/m}}

\nota{Cargas sobre la sección compuesta (barreras y capa de rodadura), distribuidas equitativamente entre todas las vigas [STD A. 3.23.2.3.1.1].}

\begin{longtable}{@{}l p{0.52\linewidth} r l l@{}}
\toprule
\textbf{Dato} & \textbf{Descripción} & \textbf{Valor} & \textbf{Unidad} & \textbf{Ref.}\\
\midrule
$N_{\mathrm{b}}$ & Número de vigas & \cellcolor{fondodato}6 & $$ & {\scriptsize } \\
\bottomrule
\end{longtable}

\eqitem{Peso de las 2 barreras por viga}{}{\mathit{WDC}_{\mathrm{b}} = \dfrac{W_{\mathrm{b}} \cdot 2}{N_{\mathrm{b}}} = \mathbf{0.1488}\,\mathrm{tonf/m}}

\begin{longtable}{@{}l p{0.52\linewidth} r l l@{}}
\toprule
\textbf{Dato} & \textbf{Descripción} & \textbf{Valor} & \textbf{Unidad} & \textbf{Ref.}\\
\midrule
$w$ & Ancho de la capa de rodadura & \cellcolor{fondodato}14.634 & $\mathrm{m}$ & {\scriptsize } \\
\bottomrule
\end{longtable}

\eqitem{Peso de la capa de rodadura por viga}{}{\mathit{WDW} = \dfrac{W_{c,ha} \cdot f_{\mathrm{ws}} \cdot w}{N_{\mathrm{b}}} = \mathbf{0.2927}\,\mathrm{tonf/m}}

\nota{Viga simplemente apoyada con carga uniforme q: V(x) = q(L/2 - x); M(x) = q$\cdot$x(L - x)/2.}

\begin{center}
{\small\textbf{Cortantes (tonf) y momentos (tonf$\cdot$m) no mayorados por CM - viga interior}}\\[3pt]
\adjustbox{max width=\linewidth}{\begin{tabular}{cccccccccc}
\toprule
\begin{tabular}[b]{@{}c@{}}\textbf{x/L}\end{tabular} & \begin{tabular}[b]{@{}c@{}}\textbf{x}\\ {\scriptsize $\mathrm{m}$}\end{tabular} & \begin{tabular}[b]{@{}c@{}}\textbf{Vg}\\ {\scriptsize $\mathrm{tonf}$}\end{tabular} & \begin{tabular}[b]{@{}c@{}}\textbf{Mg}\\ {\scriptsize $\mathrm{tonf\cdot m}$}\end{tabular} & \begin{tabular}[b]{@{}c@{}}\textbf{Vs}\\ {\scriptsize $\mathrm{tonf}$}\end{tabular} & \begin{tabular}[b]{@{}c@{}}\textbf{Ms}\\ {\scriptsize $\mathrm{tonf\cdot m}$}\end{tabular} & \begin{tabular}[b]{@{}c@{}}\textbf{Vb}\\ {\scriptsize $\mathrm{tonf}$}\end{tabular} & \begin{tabular}[b]{@{}c@{}}\textbf{Mb}\\ {\scriptsize $\mathrm{tonf\cdot m}$}\end{tabular} & \begin{tabular}[b]{@{}c@{}}\textbf{Vws}\\ {\scriptsize $\mathrm{tonf}$}\end{tabular} & \begin{tabular}[b]{@{}c@{}}\textbf{Mws}\\ {\scriptsize $\mathrm{tonf\cdot m}$}\end{tabular} \\
\midrule
\cellcolor{fondodato}$0$ & $0$ & $21.732$ & $0$ & $24.625$ & $0$ & $2.724$ & $0$ & $5.356$ & $0$ \\
\cellcolor{fondodato}$0.0278$ & $1.017$ & $20.523$ & $21.497$ & $23.255$ & $24.358$ & $2.572$ & $2.694$ & $5.058$ & $5.298$ \\
\cellcolor{fondodato}$0.1000$ & $3.660$ & $17.385$ & $71.584$ & $19.700$ & $81.113$ & $2.179$ & $8.972$ & $4.285$ & $17.643$ \\
\cellcolor{fondodato}$0.2000$ & $7.320$ & $13.039$ & $127.26$ & $14.775$ & $144.20$ & $1.634$ & $15.950$ & $3.214$ & $31.365$ \\
\cellcolor{fondodato}$0.3000$ & $10.980$ & $8.693$ & $167.03$ & $9.850$ & $189.26$ & $1.089$ & $20.934$ & $2.142$ & $41.167$ \\
\cellcolor{fondodato}$0.4000$ & $14.640$ & $4.346$ & $190.89$ & $4.925$ & $216.30$ & $0.5447$ & $23.925$ & $1.071$ & $47.047$ \\
\cellcolor{fondodato}$0.5000$ & $18.300$ & $0$ & $198.84$ & $0$ & $225.31$ & $0$ & $24.921$ & $0$ & $49.008$ \\
\bottomrule
\end{tabular}}
\end{center}

\subsubsection{4.2 Carga viva (CV): camión HS20 [STD A. 3.7.1.1]}

\eqitem{Factor de distribución (ruedas/viga)}{STD T. 3.23.1}{\mathit{DF}_{\mathrm{1}} = \dfrac{S}{5.5\,\mathrm{ft}} = \mathbf{1.640}}

\eqitem{Factor de distribución (carriles/viga)}{}{\mathit{DF} = \dfrac{\mathit{DF}_{\mathrm{1}}}{2} = \mathbf{0.8202}}

\eqitem{Factor de distribución para cortante}{}{\mathit{FDV} = \mathit{DF} = \mathbf{0.8202}}

\eqitem{Factor de distribución para momento}{}{\mathit{FDM} = \mathit{DF} = \mathbf{0.8202}}

\eqitem{Impacto I = 50/(L+125) (L en ft)}{STD Eq. 3-1}{I_c = \dfrac{50\,\mathrm{ft}}{L_{\mathrm{1}} + 125\,\mathrm{ft}} = \mathbf{0.2040}}

\eqitem{Impacto no mayor a 0.3}{STD A. 3.8.2.1}{I_c \le 0.3\;\;\Rightarrow\;\;\text{\ok{OK}}}

\eqitem{Impacto adoptado}{}{I = \min\left(I_c,\ 0.3\right) = \mathbf{0.2040}}

\begin{longtable}{@{}l p{0.52\linewidth} r l l@{}}
\toprule
\textbf{Dato} & \textbf{Descripción} & \textbf{Valor} & \textbf{Unidad} & \textbf{Ref.}\\
\midrule
$P$ & Carga por eje trasero del camión HS20 & \cellcolor{fondodato}32 & $\mathrm{kip}$ & {\scriptsize } \\
$l_1$ & Distancia entre ejes (delantero-central) & \cellcolor{fondodato}14 & $\mathrm{ft}$ & {\scriptsize } \\
$l_2$ & Distancia entre ejes (central-trasero) & \cellcolor{fondodato}14 & $\mathrm{ft}$ & {\scriptsize } \\
$\mathit{WLL}$ & Carga de carril (referencia) & \cellcolor{fondodato}0.6400 & $\mathrm{kip/ft}$ & {\scriptsize } \\
\bottomrule
\end{longtable}

\nota{Envolvente del camión móvil (ejes P/4, P, P) calculada con líneas de influencia de la viga simple: cada eje se coloca sobre la sección y se prueban ambos sentidos de circulación.}

\begin{center}
{\small\textbf{Cortantes y momentos no mayorados por camión HS20 con impacto}}\\[3pt]
\adjustbox{max width=\linewidth}{\begin{tabular}{cccc}
\toprule
\begin{tabular}[b]{@{}c@{}}\textbf{x/L}\end{tabular} & \begin{tabular}[b]{@{}c@{}}\textbf{x}\\ {\scriptsize $\mathrm{m}$}\end{tabular} & \begin{tabular}[b]{@{}c@{}}\textbf{V L+I}\\ {\scriptsize $\mathrm{tonf}$}\end{tabular} & \begin{tabular}[b]{@{}c@{}}\textbf{M L+I}\\ {\scriptsize $\mathrm{tonf\cdot m}$}\end{tabular} \\
\midrule
$0$ & $0$ & $29.745$ & $0$ \\
$0.0278$ & $1.017$ & $28.848$ & $29.353$ \\
$0.1000$ & $3.660$ & $26.520$ & $97.063$ \\
$0.2000$ & $7.320$ & $23.295$ & $170.52$ \\
$0.3000$ & $10.980$ & $20.070$ & $220.36$ \\
$0.4000$ & $14.640$ & $16.844$ & $249.66$ \\
$0.5000$ & $18.300$ & $13.619$ & $256.88$ \\
\bottomrule
\end{tabular}}
\end{center}

\subsection{5. Estimación del pretensado}

\subsubsection{5.1 Esfuerzos por carga de servicio en el centro del tramo}

\eqitem{Momento por peso de la viga}{}{M_{\mathrm{g}} = M_{\mathrm{g,6}} = \mathbf{198.84}\,\mathrm{tonf\cdot m}}

\eqitem{Momento por losa y cartela}{}{M_{\mathrm{s}} = M_{\mathrm{s,6}} = \mathbf{225.31}\,\mathrm{tonf\cdot m}}

\eqitem{Momento por barreras}{}{M_{\mathrm{b}} = M_{\mathrm{b,6}} = \mathbf{24.921}\,\mathrm{tonf\cdot m}}

\eqitem{Momento por capa de rodadura}{}{M_{\mathrm{ws}} = M_{\mathrm{ws,6}} = \mathbf{49.008}\,\mathrm{tonf\cdot m}}

\eqitem{Momento por camión + impacto}{}{\mathit{MLT} = \mathit{MLT}_{\mathrm{6}} = \mathbf{256.88}\,\mathrm{tonf\cdot m}}

\eqitem{Esfuerzo de tensión en la fibra inferior}{}{f_{\mathrm{tb}} = \dfrac{M_{\mathrm{g}} + M_{\mathrm{s}}}{S_{\mathrm{b}}} + \dfrac{M_{\mathrm{b}} + M_{\mathrm{ws}} + \mathit{MLT}}{S_{\mathrm{bc}}} = \mathbf{0.2742}\,\mathrm{tonf/cm^{2}}}

\subsubsection{5.2 Esfuerzos límites permisibles}

\eqitem{Tensión permisible en zona precomprimida}{STD A. 9.15.2.2}{F_{\mathrm{b}} = -6 \cdot \sqrt{\dfrac{f'_{c,g}}{\mathrm{psi}}} \cdot \mathrm{psi} = \mathbf{-0.0337}\,\mathrm{tonf/cm^{2}}}

\subsubsection{5.3 Número de torones}

\eqitem{Precompresión requerida en la fibra inferior}{}{f_b = f_{\mathrm{tb}} + F_{\mathrm{b}} = \mathbf{0.2405}\,\mathrm{tonf/cm^{2}}}

\eqitem{Distancia estimada CG torones - fibra inferior (0.05h)}{}{y_{\mathrm{bs,e}} = 0.05 \cdot h_{\mathrm{g}} = \mathbf{9.144}\,\mathrm{cm}}

\begin{longtable}{@{}l p{0.52\linewidth} r l l@{}}
\toprule
\textbf{Dato} & \textbf{Descripción} & \textbf{Valor} & \textbf{Unidad} & \textbf{Ref.}\\
\midrule
$y_{\mathrm{bs,a}}$ & Distancia asumida CG torones - fibra inferior & \cellcolor{fondodato}10 & $\mathrm{cm}$ & {\scriptsize } \\
\bottomrule
\end{longtable}

\eqitem{Excentricidad del torón en el centro del tramo}{}{e_{\mathrm{c1}} = y_{\mathrm{b}} - y_{\mathrm{bs,a}} = \mathbf{82.964}\,\mathrm{cm}}

\eqitem{Fuerza de pretensado total después de pérdidas}{}{P_{\mathrm{se}} = \dfrac{f_b \cdot A_{\mathrm{g}} \cdot S_{\mathrm{b}}}{S_{\mathrm{b}} + A_{\mathrm{g}} \cdot e_{\mathrm{c1}}} = \mathbf{444.01}\,\mathrm{tonf}}

\begin{longtable}{@{}l p{0.52\linewidth} r l l@{}}
\toprule
\textbf{Dato} & \textbf{Descripción} & \textbf{Valor} & \textbf{Unidad} & \textbf{Ref.}\\
\midrule
$p$ & Pérdida final asumida (\% de fsi) & \cellcolor{fondodato}25 & $\%$ & {\scriptsize } \\
\bottomrule
\end{longtable}

\eqitem{Pérdida final}{}{p_{\mathrm{f}} = p \cdot f_{\mathrm{si}} = \mathbf{3.562}\,\mathrm{tonf/cm^{2}}}

\eqitem{Fuerza por torón después de pérdidas}{}{F_{se} = A_{\mathrm{st}} \cdot \left(f_{\mathrm{si}} - p_{\mathrm{f}}\right) = \mathbf{10.549}\,\mathrm{tonf}}

\eqitem{Número de torones requerido}{}{N_{\mathrm{ti}} = \dfrac{P_{\mathrm{se}}}{F_{se}} = \mathbf{42.092}}

\eqitem{Número mínimo entero de torones}{}{N_{\mathrm{t,min}} = \left\lceil N_{\mathrm{ti}}\right\rceil = \mathbf{43}}

\begin{longtable}{@{}l p{0.52\linewidth} r l l@{}}
\toprule
\textbf{Dato} & \textbf{Descripción} & \textbf{Valor} & \textbf{Unidad} & \textbf{Ref.}\\
\midrule
$N_{\mathrm{t}}$ & Número de torones seleccionado & \cellcolor{fondodato}44 & $$ & {\scriptsize } \\
\bottomrule
\end{longtable}

\eqitem{Torones seleccionados $\ge$ requeridos}{}{N_{\mathrm{t}} \ge N_{\mathrm{ti}}\;\;\Rightarrow\;\;\text{\ok{OK}}}

\subsubsection{5.4 Distribución de torones (centro del tramo)}

\begin{center}
{\small\textbf{Distribución de torones en el centro del tramo}}\\[3pt]
\adjustbox{max width=\linewidth}{\begin{tabular}{cccc}
\toprule
\begin{tabular}[b]{@{}c@{}}\textbf{Fila}\end{tabular} & \begin{tabular}[b]{@{}c@{}}\textbf{\# Torones (A)}\end{tabular} & \begin{tabular}[b]{@{}c@{}}\textbf{Dist. al fondo (B)}\\ {\scriptsize $\mathrm{cm}$}\end{tabular} & \begin{tabular}[b]{@{}c@{}}\textbf{(A)$\cdot$(B)}\end{tabular} \\
\midrule
y1 & \cellcolor{fondodato}$12$ & \cellcolor{fondodato}$5$ & $60$ \\
y2 & \cellcolor{fondodato}$12$ & \cellcolor{fondodato}$10$ & $120$ \\
y3 & \cellcolor{fondodato}$8$ & \cellcolor{fondodato}$15$ & $120$ \\
y4 & \cellcolor{fondodato}$4$ & \cellcolor{fondodato}$20$ & $80$ \\
y5 & \cellcolor{fondodato}$2$ & \cellcolor{fondodato}$25$ & $50$ \\
y6 & \cellcolor{fondodato}$2$ & \cellcolor{fondodato}$30$ & $60$ \\
y7 & \cellcolor{fondodato}$2$ & \cellcolor{fondodato}$35$ & $70$ \\
y8 & \cellcolor{fondodato}$2$ & \cellcolor{fondodato}$40$ & $80$ \\
\bottomrule
\end{tabular}}
\end{center}

\eqitem{Total de torones en la tabla}{}{N_{\mathrm{t,tab}} = \sum A = \mathbf{44}}

\eqitem{La tabla de distribución contiene Nt torones}{}{N_{\mathrm{t,tab}} \ge N_{\mathrm{t}} \;\wedge\; N_{\mathrm{t,tab}} \le N_{\mathrm{t}}\;\;\Rightarrow\;\;\text{\ok{OK}}}

\eqitem{Distancia CG torones - fibra inferior}{}{y_{\mathrm{bs}} = \dfrac{\sum A\cdot B}{\sum A} = \mathbf{14.545}\,\mathrm{cm}}

\eqitem{Excentricidad de los torones}{}{e_{\mathrm{c}} = y_{\mathrm{b}} - y_{\mathrm{bs}} = \mathbf{78.419}\,\mathrm{cm}}

\subsection{6. Pérdidas del pretensado [STD A. 9.16.2]}

\nota{fpT = SH + ES + CRc + CRs  [STD Eq. 9.3]}

\subsubsection{6.1 Contracción del hormigón [STD A. 9.16.2.1.1]}

\begin{longtable}{@{}l p{0.52\linewidth} r l l@{}}
\toprule
\textbf{Dato} & \textbf{Descripción} & \textbf{Valor} & \textbf{Unidad} & \textbf{Ref.}\\
\midrule
$\mathit{RH}$ & Humedad relativa (\%) & \cellcolor{fondodato}70 & $$ & {\scriptsize } \\
\bottomrule
\end{longtable}

\eqitem{Pérdida por contracción}{STD Eq. 9.4}{\mathit{SH} = 17000\,\mathrm{psi} - 150\,\mathrm{psi} \cdot \mathit{RH} = \mathbf{4\,569.95}\,\mathrm{tonf/m^{2}}}

\subsubsection{6.2 Acortamiento elástico [STD A. 9.16.2.1.2]}

\eqitem{Fuerza de pretensado después de pérdidas iniciales}{}{P_{si} = N_{\mathrm{t}} \cdot A_{\mathrm{st}} \cdot 0.69 \cdot f'_s = \mathbf{569.34}\,\mathrm{tonf}}

\eqitem{Esfuerzo en el CG de los torones}{}{f_{\mathrm{cir}} = \dfrac{P_{si}}{A_{\mathrm{g}}} + \dfrac{P_{si} \cdot {e_{\mathrm{c}}}^{2}}{I_{\mathrm{g}}} - \dfrac{M_{\mathrm{g}} \cdot e_{\mathrm{c}}}{I_{\mathrm{g}}} = \mathbf{2\,005.33}\,\mathrm{tonf/m^{2}}}

\eqitem{Pérdida por acortamiento elástico}{STD Eq. 9.6}{\mathit{ES} = \dfrac{E_{\mathrm{p}}}{E_{\mathrm{ci}}} \cdot f_{\mathrm{cir}} = \mathbf{12\,845.13}\,\mathrm{tonf/m^{2}}}

\subsubsection{6.3 Flujo plástico del hormigón [STD A. 9.16.2.1.3]}

\eqitem{Esfuerzo en el CG de torones por CM}{}{f_{\mathrm{cds}} = \dfrac{M_{\mathrm{s}} \cdot e_{\mathrm{c}}}{I_{\mathrm{g}}} + \dfrac{\left(M_{\mathrm{b}} + M_{\mathrm{ws}}\right) \cdot \left(y_{\mathrm{bc}} - y_{\mathrm{bs}}\right)}{I_{\mathrm{c}}} = \mathbf{979.09}\,\mathrm{tonf/m^{2}}}

\eqitem{Pérdida por flujo plástico}{STD Eq. 9.9}{\mathit{CR}_{\mathrm{c}} = 12 \cdot f_{\mathrm{cir}} - 7 \cdot f_{\mathrm{cds}} = \mathbf{17\,210.33}\,\mathrm{tonf/m^{2}}}

\subsubsection{6.4 Relajación de los torones [STD A. 9.16.2.1.4]}

\eqitem{Pérdida por relajación (baja relajación)}{STD Eq. 9.10A}{\mathit{CR}_{\mathrm{s}} = 5\,\mathrm{ksi} - 0.1 \cdot \mathit{ES} - 0.05 \cdot \left(\mathit{SH} + \mathit{CR}_{\mathrm{c}}\right) = \mathbf{1\,141.82}\,\mathrm{tonf/m^{2}}}

\subsubsection{6.5 Pérdidas totales en la transferencia}

\eqitem{Esfuerzo efectivo inicial de pretensado}{}{f_{\mathrm{si,i}} = f_{\mathrm{si}} - \mathit{ES} = \mathbf{129\,654.87}\,\mathrm{tonf/m^{2}}}

\eqitem{Fuerza efectiva después de pérdidas iniciales}{}{P_{si,i} = N_{\mathrm{t}} \cdot A_{\mathrm{st}} \cdot f_{\mathrm{si,i}} = \mathbf{563.07}\,\mathrm{tonf}}

\subsubsection{6.6 Pérdidas totales en cargas de servicio}

\eqitem{Pérdidas totales}{STD Eq. 9.3}{f_{\mathrm{pT}} = \mathit{SH} + \mathit{ES} + \mathit{CR}_{\mathrm{c}} + \mathit{CR}_{\mathrm{s}} = \mathbf{35\,767.23}\,\mathrm{tonf/m^{2}}}

\eqitem{Porcentaje de pérdidas totales}{}{p_{\mathrm{T}} = \dfrac{f_{\mathrm{pT}}}{0.75 \cdot f'_s} = \mathbf{0.2510}}

\eqitem{Esfuerzo de pretensado final efectivo}{}{f_{\mathrm{se}} = 0.75 \cdot f'_s - f_{\mathrm{pT}} = \mathbf{106\,732.77}\,\mathrm{tonf/m^{2}}}

\eqitem{Fuerza de pretensado efectiva}{}{P_{se,c} = N_{\mathrm{t}} \cdot A_{\mathrm{st}} \cdot f_{\mathrm{se}} = \mathbf{463.52}\,\mathrm{tonf}}

\subsection{7. Esfuerzos en el hormigón durante la transferencia}

\subsubsection{7.1 Esfuerzos límites [STD A. 9.15.2.1]}

\eqitem{Compresión permisible}{}{f_{\mathrm{lc}} = 0.6 \cdot f'_{ci} = \mathbf{2\,280}\,\mathrm{tonf/m^{2}}}

\eqitem{Tensión máxima}{}{f_{\mathrm{lt1}} = -7.5 \cdot \sqrt{\dfrac{f'_{ci}}{\mathrm{psi}}} \cdot \mathrm{psi} = \mathbf{-387.66}\,\mathrm{tonf/m^{2}}}

\eqitem{Tensión sin refuerzo adherido (0.2 ksi)}{}{f_{\mathrm{lt2}} = -0.2 \cdot \mathrm{ksi} = \mathbf{-140.61}\,\mathrm{tonf/m^{2}}}

\eqitem{Tensión sin refuerzo adherido (3$\surd$f'ci)}{}{f_{\mathrm{lt3}} = -3 \cdot \sqrt{\dfrac{f'_{ci}}{\mathrm{psi}}} \cdot \mathrm{psi} = \mathbf{-155.06}\,\mathrm{tonf/m^{2}}}

\subsubsection{7.2 Esfuerzos en la longitud de transferencia}

\eqitem{Longitud de transferencia (50 d)}{STD A. 9.20.2.4}{L_{\mathrm{transf}} = 50 \cdot d_t = \mathbf{0.6350}\,\mathrm{m}}

\eqitem{Sección de análisis}{}{x_{\mathrm{t}} = L_{\mathrm{transf}} = \mathbf{0.6350}\,\mathrm{m}}

\eqitem{Momento por peso propio (viga de longitud L2)}{}{M_{\mathrm{glt}} = 0.5 \cdot W_{\mathrm{g}} \cdot x_{\mathrm{t}} \cdot \left(L_{\mathrm{2}} - x_{\mathrm{t}}\right) = \mathbf{13.673}\,\mathrm{tonf\cdot m}}

\eqitem{Esfuerzo en fibra superior}{}{f_{\mathrm{t,lt}} = \dfrac{P_{si,i}}{A_{\mathrm{g}}} - \dfrac{P_{si,i} \cdot e_{\mathrm{c}}}{S_{\mathrm{t}}} + \dfrac{M_{\mathrm{glt}}}{S_{\mathrm{t}}} = \mathbf{-555.38}\,\mathrm{tonf/m^{2}}}

\eqitem{Fibra superior (torones rectos)}{}{f_{\mathrm{t,lt}} \ge f_{\mathrm{lt1}}\;\;\Rightarrow\;\;\text{\nook{NO OK}}}

\eqitem{Esfuerzo en fibra inferior}{}{f_{\mathrm{b,lt}} = \dfrac{P_{si,i}}{A_{\mathrm{g}}} + \dfrac{P_{si,i} \cdot e_{\mathrm{c}}}{S_{\mathrm{b}}} - \dfrac{M_{\mathrm{glt}}}{S_{\mathrm{b}}} = \mathbf{2\,888.71}\,\mathrm{tonf/m^{2}}}

\eqitem{Fibra inferior (torones rectos)}{}{f_{\mathrm{b,lt}} \le f_{\mathrm{lc}}\;\;\Rightarrow\;\;\text{\nook{NO OK}}}

\nota{Los esfuerzos superan los límites: se desvían torones en los puntos 0.4L (patrón 'harped').}

\begin{longtable}{@{}l p{0.52\linewidth} r l l@{}}
\toprule
\textbf{Dato} & \textbf{Descripción} & \textbf{Valor} & \textbf{Unidad} & \textbf{Ref.}\\
\midrule
$N_{c,t}$ & Torones desviados & \cellcolor{fondodato}10 & $$ & {\scriptsize } \\
\bottomrule
\end{longtable}

\subsubsection{Patrón de torones desviados}

\begin{center}
{\small\textbf{Torones rectos en el extremo (medidos desde el fondo)}}\\[3pt]
\adjustbox{max width=\linewidth}{\begin{tabular}{cccc}
\toprule
\begin{tabular}[b]{@{}c@{}}\textbf{Fila}\end{tabular} & \begin{tabular}[b]{@{}c@{}}\textbf{\# Torones (A)}\end{tabular} & \begin{tabular}[b]{@{}c@{}}\textbf{Dist. al fondo (B)}\\ {\scriptsize $\mathrm{cm}$}\end{tabular} & \begin{tabular}[b]{@{}c@{}}\textbf{(A)$\cdot$(B)}\end{tabular} \\
\midrule
y1 & \cellcolor{fondodato}$12$ & \cellcolor{fondodato}$5$ & $60$ \\
y2 & \cellcolor{fondodato}$12$ & \cellcolor{fondodato}$10$ & $120$ \\
y3 & \cellcolor{fondodato}$8$ & \cellcolor{fondodato}$15$ & $120$ \\
y4 & \cellcolor{fondodato}$2$ & \cellcolor{fondodato}$20$ & $40$ \\
\bottomrule
\end{tabular}}
\end{center}

\eqitem{Torones inferiores rectos}{}{N_{e,b} = \sum A = \mathbf{34}}

\eqitem{CG de torones rectos - fibra inferior}{}{y_{\mathrm{be}} = \dfrac{\left(\mathit{AD}_{\mathrm{e,1}} + \mathit{AD}_{\mathrm{e,2}} + \mathit{AD}_{\mathrm{e,3}} + \mathit{AD}_{\mathrm{e,4}}\right) \cdot \mathrm{cm}}{N_{e,b}} = \mathbf{10}\,\mathrm{cm}}

\begin{center}
{\small\textbf{Torones desviados en el extremo (medidos desde arriba)}}\\[3pt]
\adjustbox{max width=\linewidth}{\begin{tabular}{cccc}
\toprule
\begin{tabular}[b]{@{}c@{}}\textbf{Fila}\end{tabular} & \begin{tabular}[b]{@{}c@{}}\textbf{\# Torones (A)}\end{tabular} & \begin{tabular}[b]{@{}c@{}}\textbf{Dist. desde arriba (B)}\\ {\scriptsize $\mathrm{cm}$}\end{tabular} & \begin{tabular}[b]{@{}c@{}}\textbf{(A)$\cdot$(B)}\end{tabular} \\
\midrule
y5 & \cellcolor{fondodato}$2$ & \cellcolor{fondodato}$25$ & $50$ \\
y6 & \cellcolor{fondodato}$2$ & \cellcolor{fondodato}$20$ & $40$ \\
y7 & \cellcolor{fondodato}$2$ & \cellcolor{fondodato}$15$ & $30$ \\
y8 & \cellcolor{fondodato}$2$ & \cellcolor{fondodato}$10$ & $20$ \\
y9 & \cellcolor{fondodato}$2$ & \cellcolor{fondodato}$5$ & $10$ \\
\bottomrule
\end{tabular}}
\end{center}

\eqitem{Torones desviados}{}{N_{e,t} = \sum A = \mathbf{10}}

\eqitem{CG de torones desviados - fibra superior (extremo)}{}{y_{\mathrm{te}} = \dfrac{\left(\mathit{AD}_{\mathrm{t,5}} + \mathit{AD}_{\mathrm{t,6}} + \mathit{AD}_{\mathrm{t,7}} + \mathit{AD}_{\mathrm{t,8}} + \mathit{AD}_{\mathrm{t,9}}\right) \cdot \mathrm{cm}}{N_{e,t}} = \mathbf{15}\,\mathrm{cm}}

\begin{center}
{\small\textbf{Torones desviados en el punto 'harp' (medidos desde el fondo)}}\\[3pt]
\adjustbox{max width=\linewidth}{\begin{tabular}{cccc}
\toprule
\begin{tabular}[b]{@{}c@{}}\textbf{Fila}\end{tabular} & \begin{tabular}[b]{@{}c@{}}\textbf{\# Torones (A)}\end{tabular} & \begin{tabular}[b]{@{}c@{}}\textbf{Dist. al fondo (B)}\\ {\scriptsize $\mathrm{cm}$}\end{tabular} & \begin{tabular}[b]{@{}c@{}}\textbf{(A)$\cdot$(B)}\end{tabular} \\
\midrule
y4 & \cellcolor{fondodato}$2$ & \cellcolor{fondodato}$20$ & $40$ \\
y5 & \cellcolor{fondodato}$2$ & \cellcolor{fondodato}$25$ & $50$ \\
y6 & \cellcolor{fondodato}$2$ & \cellcolor{fondodato}$30$ & $60$ \\
y7 & \cellcolor{fondodato}$2$ & \cellcolor{fondodato}$35$ & $70$ \\
y8 & \cellcolor{fondodato}$2$ & \cellcolor{fondodato}$40$ & $80$ \\
\bottomrule
\end{tabular}}
\end{center}

\eqitem{Torones desviados en el punto harp}{}{N_{m,b} = \sum A = \mathbf{10}}

\eqitem{CG de torones desviados - fibra inferior (punto harp)}{}{y_{\mathrm{ph}} = \dfrac{\left(\mathit{AD}_{\mathrm{h,4}} + \mathit{AD}_{\mathrm{h,5}} + \mathit{AD}_{\mathrm{h,6}} + \mathit{AD}_{\mathrm{h,7}} + \mathit{AD}_{\mathrm{h,8}}\right) \cdot \mathrm{cm}}{N_{m,b}} = \mathbf{30}\,\mathrm{cm}}

\eqitem{Ubicación del punto harp (0.4 L2)}{}{l_{\mathrm{lt}} = 0.4 \cdot L_{\mathrm{2}} = \mathbf{14.760}\,\mathrm{m}}

\eqitem{CG de torones desviados - fibra superior (long. transf.)}{}{y_{\mathrm{lt}} = y_{\mathrm{te}} + \dfrac{h_{\mathrm{g}} - y_{\mathrm{ph}} - y_{\mathrm{te}}}{l_{\mathrm{lt}}} \cdot L_{\mathrm{transf}} = \mathbf{20.932}\,\mathrm{cm}}

\begin{longtable}{@{}l p{0.52\linewidth} r l l@{}}
\toprule
\textbf{Dato} & \textbf{Descripción} & \textbf{Valor} & \textbf{Unidad} & \textbf{Ref.}\\
\midrule
$d_{\mathrm{str}}$ & Distancia entre torones & \cellcolor{fondodato}5 & $\mathrm{cm}$ & {\scriptsize } \\
\bottomrule
\end{longtable}

\eqitem{Torones rectos}{}{N_{blt} = N_{e,b} = \mathbf{34}}

\eqitem{Torones desviados}{}{N_{tlt} = N_{\mathrm{t}} - N_{blt} = \mathbf{10}}

\eqitem{Torones desviados coinciden con la tabla del extremo}{}{N_{tlt} \ge N_{e,t} \;\wedge\; N_{tlt} \le N_{e,t}\;\;\Rightarrow\;\;\text{\ok{OK}}}

\eqitem{CG de todos los torones (long. transf.)}{}{y_{\mathrm{lt}} = \dfrac{N_{blt} \cdot y_{\mathrm{be}} + N_{tlt} \cdot \left(h_{\mathrm{g}} - y_{\mathrm{lt}}\right)}{N_{\mathrm{t}}} = \mathbf{44.534}\,\mathrm{cm}}

\begin{center}
{\small\textbf{Patrón de torones en el extremo (desde el fondo)}}\\[3pt]
\adjustbox{max width=\linewidth}{\begin{tabular}{cccc}
\toprule
\begin{tabular}[b]{@{}c@{}}\textbf{Fila}\end{tabular} & \begin{tabular}[b]{@{}c@{}}\textbf{\# Torones (A)}\end{tabular} & \begin{tabular}[b]{@{}c@{}}\textbf{Dist. al fondo (cm) (B)}\end{tabular} & \begin{tabular}[b]{@{}c@{}}\textbf{(A)$\cdot$(B)}\end{tabular} \\
\midrule
y1 & $12$ & $5$ & $60$ \\
y2 & $12$ & $10$ & $120$ \\
y3 & $8$ & $15$ & $120$ \\
y4 & $2$ & $20$ & $40$ \\
y6 & $2$ & $157.88$ & $315.76$ \\
y7 & $2$ & $162.88$ & $325.76$ \\
y8 & $2$ & $167.88$ & $335.76$ \\
y9 & $2$ & $172.88$ & $345.76$ \\
y10 & $2$ & $177.88$ & $355.76$ \\
\bottomrule
\end{tabular}}
\end{center}

\eqitem{CG de todos los torones - fondo (extremo)}{}{y_{\mathrm{bs,ext}} = \dfrac{\left(\mathit{AE}_{\mathrm{x,1}} + \mathit{AE}_{\mathrm{x,2}} + \mathit{AE}_{\mathrm{x,3}} + \mathit{AE}_{\mathrm{x,4}} + \mathit{AE}_{\mathrm{x,5}} + \mathit{AE}_{\mathrm{x,6}} + \mathit{AE}_{\mathrm{x,7}} + \mathit{AE}_{\mathrm{x,8}} + \mathit{AE}_{\mathrm{x,9}}\right) \cdot \mathrm{cm}}{N_{\mathrm{t}}} = \mathbf{45.882}\,\mathrm{cm}}

\eqitem{Excentricidad en la longitud de transferencia}{}{e_{lt} = y_{\mathrm{b}} - y_{\mathrm{lt}} = \mathbf{48.430}\,\mathrm{cm}}

\nota{Recálculo de esfuerzos en la sección de longitud de transferencia con torones desviados.}

\eqitem{Esfuerzo en fibra superior}{}{f_{\mathrm{t,ext}} = \dfrac{P_{si,i}}{A_{\mathrm{g}}} - \dfrac{P_{si,i} \cdot e_{lt}}{S_{\mathrm{t}}} + \dfrac{M_{\mathrm{glt}}}{S_{\mathrm{t}}} = \mathbf{112.87}\,\mathrm{tonf/m^{2}}}

\eqitem{Fibra superior (con torones desviados)}{}{f_{\mathrm{t,ext}} \le f_{\mathrm{lc}}\;\;\Rightarrow\;\;\text{\ok{OK}}}

\eqitem{Esfuerzo en fibra inferior}{}{f_{\mathrm{b,ext}} = \dfrac{P_{si,i}}{A_{\mathrm{g}}} + \dfrac{P_{si,i} \cdot e_{lt}}{S_{\mathrm{b}}} - \dfrac{M_{\mathrm{glt}}}{S_{\mathrm{b}}} = \mathbf{2\,197.81}\,\mathrm{tonf/m^{2}}}

\eqitem{Fibra inferior (con torones desviados)}{}{f_{\mathrm{b,ext}} \le f_{\mathrm{lc}}\;\;\Rightarrow\;\;\text{\ok{OK}}}

\subsubsection{7.3 Esfuerzos en los puntos 'harp'}

\eqitem{Ubicación del punto harp}{}{x_{\mathrm{ph}} = l_{\mathrm{lt}} = \mathbf{14.760}\,\mathrm{m}}

\eqitem{Momento por peso de la viga}{}{M_{\mathrm{gph}} = 0.5 \cdot W_{\mathrm{g}} \cdot x_{\mathrm{ph}} \cdot \left(L_{\mathrm{2}} - x_{\mathrm{ph}}\right) = \mathbf{194.03}\,\mathrm{tonf\cdot m}}

\eqitem{Esfuerzo en fibra superior}{}{f_{\mathrm{t,ph}} = \dfrac{P_{si,i}}{A_{\mathrm{g}}} - \dfrac{P_{si,i} \cdot e_{\mathrm{c}}}{S_{\mathrm{t}}} + \dfrac{M_{\mathrm{gph}}}{S_{\mathrm{t}}} = \mathbf{158.41}\,\mathrm{tonf/m^{2}}}

\eqitem{Fibra superior en punto harp}{}{f_{\mathrm{t,ph}} \le f_{\mathrm{lc}}\;\;\Rightarrow\;\;\text{\ok{OK}}}

\eqitem{Esfuerzo en fibra inferior}{}{f_{\mathrm{b,ph}} = \dfrac{P_{si,i}}{A_{\mathrm{g}}} + \dfrac{P_{si,i} \cdot e_{\mathrm{c}}}{S_{\mathrm{b}}} - \dfrac{M_{\mathrm{gph}}}{S_{\mathrm{b}}} = \mathbf{2\,150.73}\,\mathrm{tonf/m^{2}}}

\eqitem{Fibra inferior en punto harp}{}{f_{\mathrm{b,ph}} \le f_{\mathrm{lc}}\;\;\Rightarrow\;\;\text{\ok{OK}}}

\subsubsection{7.4 Esfuerzos en el centro del tramo}

\eqitem{Centro del tramo}{}{x_{\mathrm{cl}} = 0.5 \cdot L_{\mathrm{2}} = \mathbf{18.450}\,\mathrm{m}}

\eqitem{Momento por peso de la viga}{}{M_{\mathrm{gcl}} = 0.5 \cdot W_{\mathrm{g}} \cdot x_{\mathrm{cl}} \cdot \left(L_{\mathrm{2}} - x_{\mathrm{cl}}\right) = \mathbf{202.12}\,\mathrm{tonf\cdot m}}

\eqitem{Esfuerzo en fibra superior}{}{f_{\mathrm{t,cl}} = \dfrac{P_{si,i}}{A_{\mathrm{g}}} - \dfrac{P_{si,i} \cdot e_{\mathrm{c}}}{S_{\mathrm{t}}} + \dfrac{M_{\mathrm{gcl}}}{S_{\mathrm{t}}} = \mathbf{190.40}\,\mathrm{tonf/m^{2}}}

\eqitem{Fibra superior en centro del tramo}{}{f_{\mathrm{t,cl}} \le f_{\mathrm{lc}}\;\;\Rightarrow\;\;\text{\ok{OK}}}

\eqitem{Esfuerzo en fibra inferior}{}{f_{\mathrm{b,cl}} = \dfrac{P_{si,i}}{A_{\mathrm{g}}} + \dfrac{P_{si,i} \cdot e_{\mathrm{c}}}{S_{\mathrm{b}}} - \dfrac{M_{\mathrm{gcl}}}{S_{\mathrm{b}}} = \mathbf{2\,117.65}\,\mathrm{tonf/m^{2}}}

\eqitem{Fibra inferior en centro del tramo}{}{f_{\mathrm{b,cl}} \le f_{\mathrm{lc}}\;\;\Rightarrow\;\;\text{\ok{OK}}}

\subsubsection{7.5 Fuerza de sujeción}

\eqitem{Esfuerzo inicial máximo antes de pérdidas}{}{f_s = 0.8 \cdot f'_s = \mathbf{152\,000}\,\mathrm{tonf/m^{2}}}

\eqitem{Fuerza inicial por torón}{}{F_s = A_{\mathrm{st}} \cdot f_s = \mathbf{15.002}\,\mathrm{tonf}}

\eqitem{Ángulo 'harp'}{}{\psi = \arctan\left(\dfrac{h_{\mathrm{g}} - y_{\mathrm{ph}} - y_{\mathrm{te}}}{l_{\mathrm{lt}}}\right) = \mathbf{5.337}^{\circ}}

\begin{longtable}{@{}l p{0.52\linewidth} r l l@{}}
\toprule
\textbf{Dato} & \textbf{Descripción} & \textbf{Valor} & \textbf{Unidad} & \textbf{Ref.}\\
\midrule
$f_f$ & Factor de fricción & \cellcolor{fondodato}1.050 & $$ & {\scriptsize } \\
\bottomrule
\end{longtable}

\eqitem{Fuerza de sujeción por torón}{}{F_{\mathrm{suj}} = f_f \cdot F_s \cdot \sin\left(\psi\right) = \mathbf{1.465}\,\mathrm{tonf}}

\eqitem{Fuerza de sujeción total}{}{F_{\mathrm{sujtot}} = N_{tlt} \cdot F_{\mathrm{suj}} = \mathbf{14.651}\,\mathrm{tonf}}

\subsection{8. Esfuerzos en el hormigón debido a las cargas de servicio}

\subsubsection{8.1 Esfuerzos límites [STD A. 9.15.2.2]}

\eqitem{Caso I - compresión en viga}{}{f_{\mathrm{pcsl2}} = 0.6 \cdot f'_{c,g} = \mathbf{2\,700}\,\mathrm{tonf/m^{2}}}

\eqitem{Caso I - compresión en losa}{}{f_{\mathrm{pcsl1}} = 0.6 \cdot f'_{c,s} = \mathbf{1\,680}\,\mathrm{tonf/m^{2}}}

\eqitem{Tensión en viga}{}{f_{\mathrm{pttsl2}} = -6 \cdot \sqrt{\dfrac{f'_{c,g}}{\mathrm{psi}}} \cdot \mathrm{psi} = \mathbf{-337.49}\,\mathrm{tonf/m^{2}}}

\eqitem{Caso II - compresión en viga}{}{f_{\mathrm{ptcsl2}} = 0.4 \cdot f'_{c,g} = \mathbf{1\,800}\,\mathrm{tonf/m^{2}}}

\eqitem{Caso II - compresión en losa}{}{f_{\mathrm{ptcsl1}} = 0.4 \cdot f'_{c,s} = \mathbf{1\,120}\,\mathrm{tonf/m^{2}}}

\eqitem{Caso III - compresión en viga}{}{f_{\mathrm{ciii2}} = 0.4 \cdot f'_{c,g} = \mathbf{1\,800}\,\mathrm{tonf/m^{2}}}

\eqitem{Caso III - compresión en losa}{}{f_{\mathrm{ciii1}} = 0.4 \cdot f'_{c,s} = \mathbf{1\,120}\,\mathrm{tonf/m^{2}}}

\subsubsection{8.2 Esfuerzos en el centro del tramo}

\eqitem{Fibra superior de la viga - Caso I}{}{f_{\mathrm{tg1}} = \dfrac{P_{se,c}}{A_{\mathrm{g}}} - \dfrac{P_{se,c} \cdot e_{\mathrm{c}}}{S_{\mathrm{t}}} + \dfrac{M_{\mathrm{g}} + M_{\mathrm{s}}}{S_{\mathrm{t}}} + \dfrac{M_{\mathrm{ws}} + M_{\mathrm{b}}}{S_{\mathrm{tg}}} + \dfrac{\mathit{MLT}}{S_{\mathrm{tg}}} = \mathbf{1\,494.58}\,\mathrm{tonf/m^{2}}}

\eqitem{Fibra superior viga - Caso I}{}{f_{\mathrm{tg1}} \le f_{\mathrm{pcsl2}}\;\;\Rightarrow\;\;\text{\ok{OK}}}

\eqitem{Fibra superior de la viga - Caso II}{}{f_{\mathrm{tg2}} = \dfrac{P_{se,c}}{A_{\mathrm{g}}} - \dfrac{P_{se,c} \cdot e_{\mathrm{c}}}{S_{\mathrm{t}}} + \dfrac{M_{\mathrm{g}} + M_{\mathrm{s}}}{S_{\mathrm{t}}} + \dfrac{M_{\mathrm{ws}} + M_{\mathrm{b}}}{S_{\mathrm{tg}}} = \mathbf{1\,247.90}\,\mathrm{tonf/m^{2}}}

\eqitem{Fibra superior viga - Caso II}{}{f_{\mathrm{tg2}} \le f_{\mathrm{ptcsl2}}\;\;\Rightarrow\;\;\text{\ok{OK}}}

\eqitem{Fibra superior de la viga - Caso III}{}{f_{\mathrm{tg3}} = \dfrac{f_{\mathrm{tg2}}}{2} + \dfrac{\mathit{MLT}}{S_{\mathrm{tg}}} = \mathbf{870.63}\,\mathrm{tonf/m^{2}}}

\eqitem{Fibra superior viga - Caso III}{}{f_{\mathrm{tg3}} \le f_{\mathrm{ciii2}}\;\;\Rightarrow\;\;\text{\ok{OK}}}

\eqitem{Fibra inferior de la viga}{}{f_b = \dfrac{P_{se,c}}{A_{\mathrm{g}}} + \dfrac{P_{se,c} \cdot e_{\mathrm{c}}}{S_{\mathrm{b}}} - \dfrac{M_{\mathrm{g}} + M_{\mathrm{s}}}{S_{\mathrm{b}}} - \dfrac{M_{\mathrm{ws}} + M_{\mathrm{b}}}{S_{\mathrm{bc}}} - \dfrac{\mathit{MLT}}{S_{\mathrm{bc}}} = \mathbf{-318.06}\,\mathrm{tonf/m^{2}}}

\eqitem{Fibra inferior viga - compresión}{}{f_b \le f_{\mathrm{pcsl2}}\;\;\Rightarrow\;\;\text{\ok{OK}}}

\eqitem{Fibra inferior viga - tensión ($\ge$ -6$\surd$f'c)}{}{f_b \ge f_{\mathrm{pttsl2}}\;\;\Rightarrow\;\;\text{\ok{OK}}}

\eqitem{Fibra superior de la losa - Caso I}{}{f_{\mathrm{tg1t}} = \dfrac{M_{\mathrm{ws}} + M_{\mathrm{b}}}{S_{\mathrm{tc}}} + \dfrac{\mathit{MLT}}{S_{\mathrm{tc}}} = \mathbf{364.83}\,\mathrm{tonf/m^{2}}}

\eqitem{Losa - Caso I}{}{f_{\mathrm{tg1t}} \le f_{\mathrm{pcsl1}}\;\;\Rightarrow\;\;\text{\ok{OK}}}

\eqitem{Fibra superior de la losa - Caso II}{}{f_{\mathrm{tg2t}} = \dfrac{M_{\mathrm{ws}} + M_{\mathrm{b}}}{S_{\mathrm{tc}}} = \mathbf{81.532}\,\mathrm{tonf/m^{2}}}

\eqitem{Losa - Caso II}{}{f_{\mathrm{tg2t}} \le f_{\mathrm{ptcsl1}}\;\;\Rightarrow\;\;\text{\ok{OK}}}

\eqitem{Fibra superior de la losa - Caso III}{}{f_{\mathrm{tg3t}} = \dfrac{0.5 \cdot \left(M_{\mathrm{ws}} + M_{\mathrm{b}}\right)}{S_{\mathrm{tc}}} + \dfrac{\mathit{MLT}}{S_{\mathrm{tc}}} = \mathbf{324.06}\,\mathrm{tonf/m^{2}}}

\eqitem{Losa - Caso III}{}{f_{\mathrm{tg3t}} \le f_{\mathrm{ciii1}}\;\;\Rightarrow\;\;\text{\ok{OK}}}

\subsection{9. Resistencia a la flexión [STD A. 9.17]}

\nota{Momento de diseño: Grupo I - combinación LFD [STD T. 3.22.1A].}

\begin{longtable}{@{}l p{0.52\linewidth} r l l@{}}
\toprule
\textbf{Dato} & \textbf{Descripción} & \textbf{Valor} & \textbf{Unidad} & \textbf{Ref.}\\
\midrule
$\gamma$ & Factor $\gamma$ & \cellcolor{fondodato}1.300 & $$ & {\scriptsize } \\
$\beta_D$ & Factor $\beta$D & \cellcolor{fondodato}1 & $$ & {\scriptsize } \\
$\beta_L$ & Factor $\beta$L & \cellcolor{fondodato}1.670 & $$ & {\scriptsize } \\
\bottomrule
\end{longtable}

\eqitem{Momento último}{}{M_{\mathrm{u}} = \gamma \cdot \left(\beta_D \cdot \left(M_{\mathrm{g}} + M_{\mathrm{b}} + M_{\mathrm{s}} + M_{\mathrm{ws}}\right) + \beta_L \cdot \mathit{MLT}\right) = \mathbf{1\,205.19}\,\mathrm{tonf\cdot m}}

\begin{longtable}{@{}l p{0.52\linewidth} r l l@{}}
\toprule
\textbf{Dato} & \textbf{Descripción} & \textbf{Valor} & \textbf{Unidad} & \textbf{Ref.}\\
\midrule
$\gamma'$ & $\gamma$' para torón de baja relajación & \cellcolor{fondodato}0.2800 & $$ & {\scriptsize STD A. 9.1.2} \\
\bottomrule
\end{longtable}

\eqitem{Factor del bloque de compresión $\beta$1}{STD A. 8.16.2.7}{\beta_{\mathrm{1}} = \begin{cases}0.85 & \text{si } \operatorname{red}\left(\dfrac{f'_{c,s}}{\dfrac{\mathrm{kgf}}{{\mathrm{cm}}^{2}}}\right) \le 280\\ \begin{cases}0.65 & \text{si } \operatorname{red}\left(\dfrac{f'_{c,s}}{\dfrac{\mathrm{kgf}}{{\mathrm{cm}}^{2}}}\right) > 550\\ 0.85 - \dfrac{0.05 \cdot \left(f'_{c,s} - \dfrac{280\,\mathrm{kgf}}{{\mathrm{cm}}^{2}}\right)}{\dfrac{70\,\mathrm{kgf}}{{\mathrm{cm}}^{2}}} & \text{en otro caso}\end{cases} & \text{en otro caso}\end{cases} = \mathbf{0.8500}}

\eqitem{Área de los torones A's}{}{A'_s = N_{\mathrm{t}} \cdot A_{\mathrm{st}} = \mathbf{4\,342.80}\,\mathrm{mm^{2}}}

\eqitem{Ancho efectivo del ala}{}{b = w_{\mathrm{fe}} = \mathbf{275}\,\mathrm{cm}}

\eqitem{Peralte al centroide de los torones}{}{d = h_{\mathrm{g}} + t_{\mathrm{h}} + t_{\mathrm{s}} - y_{\mathrm{bs}} = \mathbf{188.33}\,\mathrm{cm}}

\eqitem{Cuantía $\rho$'}{}{\rho' = \dfrac{A'_s}{b \cdot d} = \mathbf{8.385\times10^{-4}}}

\eqitem{Esfuerzo promedio en el torón (carga última)}{STD Eq. 9-17}{f^*_{su} = f'_s \cdot \left(1 - \dfrac{\dfrac{\gamma'}{\beta_{\mathrm{1}}} \cdot \rho' \cdot f'_s}{f'_{c,s}}\right) = \mathbf{186\,438.81}\,\mathrm{tonf/m^{2}}}

\eqitem{Profundidad del bloque de compresión}{STD Eq. 9-17.2}{a = \dfrac{A'_s \cdot f^*_{su}}{0.85 \cdot f'_{c,s} \cdot b} = \mathbf{12.371}\,\mathrm{cm}}

\eqitem{Bloque dentro del ala (sección rectangular)}{}{a < t_{\mathrm{s}}\;\;\Rightarrow\;\;\text{\ok{OK}}}

\begin{longtable}{@{}l p{0.52\linewidth} r l l@{}}
\toprule
\textbf{Dato} & \textbf{Descripción} & \textbf{Valor} & \textbf{Unidad} & \textbf{Ref.}\\
\midrule
$\phi$ & Factor de reducción de resistencia (flexión) & \cellcolor{fondodato}1 & $$ & {\scriptsize STD A. 9.14} \\
\bottomrule
\end{longtable}

\eqitem{Resistencia nominal a flexión}{STD Eq. 9-13}{M_{\mathrm{n}} = A'_s \cdot f^*_{su} \cdot d \cdot \left(1 - \dfrac{0.6 \cdot \rho' \cdot f^*_{su}}{f'_{c,s}}\right) = \mathbf{1\,473.80}\,\mathrm{tonf\cdot m}}

\eqitem{Resistencia a flexión mayorada}{}{M_{\mathrm{r}} = \phi \cdot M_{\mathrm{n}} = \mathbf{1\,473.80}\,\mathrm{tonf\cdot m}}

\eqitem{Resistencia a flexión Mr > Mu}{}{M_{\mathrm{r}} > M_{\mathrm{u}}\;\;\Rightarrow\;\;\text{\ok{OK}}}

\subsection{10. Límites del refuerzo}

\subsubsection{10.1 Máximo refuerzo [STD A. 9.18.1]}

\eqitem{Índice de refuerzo $\rho$'$\cdot$f*su/f'c}{STD Eq. 9-20}{I_{\mathrm{r}} = \dfrac{\rho' \cdot f^*_{su}}{f'_{c,s}} = \mathbf{0.0558}}

\eqitem{Refuerzo máximo}{}{I_{\mathrm{r}} < 0.36 \cdot \beta_{\mathrm{1}}\;\;\Rightarrow\;\;\text{\ok{OK}}}

\subsubsection{10.2 Mínimo refuerzo [STD A. 9.18.2]}

\eqitem{Módulo de ruptura}{STD A. 9.15.2.3}{f_{\mathrm{r}} = 7.5 \cdot \sqrt{\dfrac{f'_{c,g}}{\mathrm{psi}}} \cdot \mathrm{psi} = \mathbf{421.86}\,\mathrm{tonf/m^{2}}}

\eqitem{Compresión por pretensado efectivo en fibra inferior}{}{f_{\mathrm{pe}} = \dfrac{P_{se,c}}{A_{\mathrm{g}}} + \dfrac{P_{se,c} \cdot e_{\mathrm{c}}}{S_{\mathrm{b}}} = \mathbf{2\,424.06}\,\mathrm{tonf/m^{2}}}

\eqitem{Momento por CM en la sección no compuesta}{}{M_{\mathrm{dnc}} = M_{\mathrm{g}} + M_{\mathrm{s}} = \mathbf{424.16}\,\mathrm{tonf\cdot m}}

\eqitem{Momento de fisuración}{}{M_{\mathrm{cr}} = S_{\mathrm{bc}} \cdot \left(f_{\mathrm{r}} + f_{\mathrm{pe}}\right) - M_{\mathrm{dnc}} \cdot \left(\dfrac{S_{\mathrm{bc}}}{S_{\mathrm{b}}} - 1\right) = \mathbf{789.08}\,\mathrm{tonf\cdot m}}

\eqitem{1.2 Mcr}{}{1.2M_{cr} = 1.2 \cdot M_{\mathrm{cr}} = \mathbf{946.89}\,\mathrm{tonf\cdot m}}

\eqitem{Refuerzo mínimo: $\phi$Mn $\ge$ 1.2 Mcr}{}{1.2M_{cr} < M_{\mathrm{r}}\;\;\Rightarrow\;\;\text{\ok{OK}}}

\subsection{11. Diseño por corte [STD A. 9.20]}

\nota{Vu $\le$ $\phi$(Vc + Vs)  [STD Eq. 9-26]. La sección crítica se considera a hc/2 desde el eje del apoyo.}

\begin{longtable}{@{}l p{0.52\linewidth} r l l@{}}
\toprule
\textbf{Dato} & \textbf{Descripción} & \textbf{Valor} & \textbf{Unidad} & \textbf{Ref.}\\
\midrule
$\phi_c$ & Factor de reducción para cortante & \cellcolor{fondodato}0.9000 & $$ & {\scriptsize STD A. 9.14} \\
\bottomrule
\end{longtable}

\eqitem{Ubicación de la sección crítica (h/2)}{}{h_{\mathrm{c}} = 0.5 \cdot \left(h_{\mathrm{g}} + t\right) = \mathbf{101.44}\,\mathrm{cm}}

\nota{Las fuerzas en la sección crítica se toman de la fila x/L = 0.0278 de las tablas.}

\eqitem{Ancho del alma b'}{}{b' = t_{\mathrm{w}} = \mathbf{15.240}\,\mathrm{cm}}

\eqitem{Cortante por CM}{}{V_{\mathrm{d}} = V_{\mathrm{g,1}} + V_{\mathrm{s,1}} + V_{\mathrm{b,1}} + V_{\mathrm{ws,1}} = \mathbf{51.409}\,\mathrm{tonf}}

\eqitem{Cortante por CV + impacto}{}{\mathit{VLT}_{\mathrm{1}} = \mathit{VLT}_{\mathrm{1}} = \mathbf{28.848}\,\mathrm{tonf}}

\eqitem{Momento por CM}{}{M_{\mathrm{d}} = M_{\mathrm{g,1}} + M_{\mathrm{s,1}} + M_{\mathrm{b,1}} + M_{\mathrm{ws,1}} = \mathbf{53.848}\,\mathrm{tonf\cdot m}}

\eqitem{Momento por CV + impacto}{}{\mathit{MLT}_{\mathrm{1}} = \mathit{MLT}_{\mathrm{1}} = \mathbf{29.353}\,\mathrm{tonf\cdot m}}

\eqitem{Cortante mayorado}{}{V_{\mathrm{u}} = \gamma \cdot \left(\beta_D \cdot V_{\mathrm{d}} + \beta_L \cdot \mathit{VLT}_{\mathrm{1}}\right) = \mathbf{129.46}\,\mathrm{tonf}}

\eqitem{Momento mayorado}{}{M_{\mathrm{u1}} = \gamma \cdot \left(\beta_D \cdot M_{\mathrm{d}} + \beta_L \cdot \mathit{MLT}_{\mathrm{1}}\right) = \mathbf{133.73}\,\mathrm{tonf\cdot m}}

\eqitem{Cortante mayorado simultáneo con Mu}{}{V_{\mathrm{mu}} = V_{\mathrm{u}} = \mathbf{129.46}\,\mathrm{tonf}}

\eqitem{Momento mayorado por cargas externas}{}{M_{\mathrm{max}} = M_{\mathrm{u1}} - M_{\mathrm{d}} = \mathbf{79.879}\,\mathrm{tonf\cdot m}}

\eqitem{Cortante mayorado por cargas externas}{}{V_{\mathrm{i}} = V_{\mathrm{u}} - V_{\mathrm{d}} = \mathbf{78.053}\,\mathrm{tonf}}

\eqitem{Distancia eje del apoyo - extremo de la viga}{}{d_{\mathrm{ext}} = 0.5 \cdot \left(L_{\mathrm{2}} - L_{\mathrm{1}}\right) = \mathbf{15}\,\mathrm{cm}}

\eqitem{CG de torones desviados desde arriba (h/2)}{}{c_{\mathrm{g10t}} = y_{\mathrm{te}} + \dfrac{h_{\mathrm{g}} - y_{\mathrm{ph}} - y_{\mathrm{te}}}{0.4 \cdot L_{\mathrm{2}}} \cdot \left(h_{\mathrm{c}} + d_{\mathrm{ext}}\right) = \mathbf{25.877}\,\mathrm{cm}}

\eqitem{CG de todos los torones desde el fondo (h/2)}{}{c_{\mathrm{gtt}} = \dfrac{N_{m,b} \cdot \left(h_{\mathrm{g}} - c_{\mathrm{g10t}}\right) + N_{e,b} \cdot y_{\mathrm{be}}}{N_{\mathrm{t}}} = \mathbf{43.410}\,\mathrm{cm}}

\eqitem{Excentricidad de los torones en h/2}{}{e_t = y_{\mathrm{b}} - c_{\mathrm{gtt}} = \mathbf{49.554}\,\mathrm{cm}}

\eqitem{Compresión por pretensado (fibra inferior)}{}{f_{\mathrm{peb}} = \dfrac{P_{se,c}}{A_{\mathrm{g}}} + \dfrac{P_{se,c} \cdot e_t}{S_{\mathrm{b}}} = \mathbf{1\,876.62}\,\mathrm{tonf/m^{2}}}

\eqitem{Esfuerzo por CM no mayorada}{}{f_{\mathrm{d}} = \dfrac{M_{\mathrm{g,1}} + M_{\mathrm{s,1}}}{S_{\mathrm{b}}} + \dfrac{M_{\mathrm{b,1}} + M_{\mathrm{ws,1}}}{S_{\mathrm{bc}}} = \mathbf{211.95}\,\mathrm{tonf/m^{2}}}

\eqitem{Momento de fisuración por cargas externas}{STD Eq. 9.28}{M_{cr,e} = \left(6 \cdot \sqrt{\dfrac{f'_{c,g}}{\mathrm{psi}}} \cdot \mathrm{psi} + f_{\mathrm{peb}} - f_{\mathrm{d}}\right) \cdot S_{\mathrm{bc}} = \mathbf{658.00}\,\mathrm{tonf\cdot m}}

\eqitem{0.8 hc}{}{d_{\mathrm{1}} = 0.8 \cdot h_{\mathrm{sc}} = \mathbf{162.30}\,\mathrm{cm}}

\eqitem{Distancia fibra comprimida - CG del refuerzo}{}{d_{\mathrm{2}} = h_{\mathrm{sc}} - c_{\mathrm{gtt}} = \mathbf{159.47}\,\mathrm{cm}}

\eqitem{d $\ge$ 0.8 hc}{STD A. 9.20.2.2}{d_{\mathrm{2}} > 0.8 \cdot h_{\mathrm{sc}}\;\;\Rightarrow\;\;\text{\info{NO OK}}}

\eqitem{Peralte d adoptado}{}{d_{\mathrm{v1}} = \max\left(d_{\mathrm{2}},\ d_{\mathrm{1}}\right) = \mathbf{162.30}\,\mathrm{cm}}

\eqitem{Resistencia Vci}{STD Eq. 9-27}{V_{\mathrm{ci}} = 0.6 \cdot \sqrt{\dfrac{f'_{c,g}}{\mathrm{psi}}} \cdot \mathrm{psi} \cdot b' \cdot d_{\mathrm{v1}} + V_{\mathrm{d}} + \dfrac{V_{\mathrm{i}} \cdot M_{cr,e}}{M_{\mathrm{max}}} = \mathbf{702.71}\,\mathrm{tonf}}

\eqitem{Vci mínimo}{STD A. 9.20.2.2}{V_{\mathrm{cimin}} = 1.7 \cdot \sqrt{\dfrac{f'_{c,g}}{\mathrm{psi}}} \cdot \mathrm{psi} \cdot b' \cdot d_{\mathrm{v1}} = \mathbf{23.652}\,\mathrm{tonf}}

\eqitem{Vci $\ge$ Vci mín}{}{V_{\mathrm{cimin}} < V_{\mathrm{ci}}\;\;\Rightarrow\;\;\text{\ok{OK}}}

\eqitem{Compresión en el centroide de la sección compuesta}{}{f_{\mathrm{pc}} = \dfrac{P_{se,c}}{A_{\mathrm{g}}} - \dfrac{P_{se,c} \cdot e_t \cdot \left(y_{\mathrm{bc}} - y_{\mathrm{b}}\right)}{I_{\mathrm{g}}} + \dfrac{\left(M_{\mathrm{g,1}} + M_{\mathrm{s,1}}\right) \cdot \left(y_{\mathrm{bc}} - y_{\mathrm{b}}\right)}{I_{\mathrm{g}}} = \mathbf{564.21}\,\mathrm{tonf/m^{2}}}

\eqitem{Pretensado efectivo de torones desviados}{}{P_{\mathrm{seph}} = N_{e,t} \cdot A_{\mathrm{st}} \cdot f_{\mathrm{se}} = \mathbf{105.35}\,\mathrm{tonf}}

\eqitem{Componente vertical del pretensado}{}{V_{\mathrm{p}} = P_{\mathrm{seph}} \cdot \sin\left(\psi\right) = \mathbf{9.798}\,\mathrm{tonf}}

\eqitem{Resistencia Vcw}{STD Eq. 9-29}{V_{\mathrm{cw}} = \left(3.5 \cdot \sqrt{\dfrac{f'_{c,g}}{\mathrm{psi}}} \cdot \mathrm{psi} + 0.3 \cdot f_{\mathrm{pc}}\right) \cdot b' \cdot d_{\mathrm{v1}} + V_{\mathrm{p}} = \mathbf{100.36}\,\mathrm{tonf}}

\eqitem{Resistencia nominal del hormigón}{}{V_{\mathrm{c}} = \min\left(V_{\mathrm{ci}},\ V_{\mathrm{cw}}\right) = \mathbf{100.36}\,\mathrm{tonf}}

\eqitem{Resistencia requerida del refuerzo}{}{V_{\mathrm{s}} = \dfrac{V_{\mathrm{u}}}{\phi_c} - V_{\mathrm{c}} = \mathbf{43.486}\,\mathrm{tonf}}

\eqitem{Vs máximo}{STD A. 9.20.3.1}{V_{\mathrm{smax}} = 8 \cdot \sqrt{\dfrac{f'_{c,g}}{\mathrm{psi}}} \cdot \mathrm{psi} \cdot b' \cdot d_{\mathrm{v1}} = \mathbf{111.30}\,\mathrm{tonf}}

\eqitem{Vs $\le$ Vs máx}{}{V_{\mathrm{smax}} > V_{\mathrm{s}}\;\;\Rightarrow\;\;\text{\ok{OK}}}

\begin{longtable}{@{}l p{0.52\linewidth} r l l@{}}
\toprule
\textbf{Dato} & \textbf{Descripción} & \textbf{Valor} & \textbf{Unidad} & \textbf{Ref.}\\
\midrule
$s_{\mathrm{p}}$ & Separación de estribos escogida & \cellcolor{fondodato}30 & $\mathrm{cm}$ & {\scriptsize } \\
\bottomrule
\end{longtable}

\eqitem{Área de refuerzo requerida}{STD Eq. 9.30}{A_{\mathrm{vp}} = \dfrac{s_{\mathrm{p}} \cdot V_{\mathrm{s}}}{f_{\mathrm{y}} \cdot d_{\mathrm{v1}}} = \mathbf{1.914}\,\mathrm{cm^{2}}}

\eqitem{Área mínima de refuerzo transversal}{STD Eq. 9.31}{A_{\mathrm{vmin}} = \dfrac{50 \cdot b' \cdot s_{\mathrm{p}}}{\dfrac{f_{\mathrm{y}}}{\mathrm{psi}}} = \mathbf{0.3827}\,\mathrm{cm^{2}}}

\eqitem{Área requerida}{}{A_{\mathrm{vmax}} = \max\left(A_{\mathrm{vmin}},\ A_{\mathrm{vp}}\right) = \mathbf{1.914}\,\mathrm{cm^{2}}}

\eqitem{Límite 4$\surd$f'c b'd}{STD A. 9.20.3.2}{V_{\mathrm{sc}} = 4 \cdot \sqrt{\dfrac{f'_{c,g}}{\mathrm{psi}}} \cdot \mathrm{psi} \cdot b' \cdot d_{\mathrm{v1}} = \mathbf{55.652}\,\mathrm{tonf}}

\eqitem{Vs > 4$\surd$f'c b'd (espaciamiento reducido)}{}{V_{\mathrm{s}} > V_{\mathrm{sc}}\;\;\Rightarrow\;\;\text{\info{NO OK}}}

\eqitem{Espaciamiento máx. (I) 0.75 hc}{}{s_{\mathrm{c}} = 0.75 \cdot h_{\mathrm{sc}} = \mathbf{152.16}\,\mathrm{cm}}

\begin{longtable}{@{}l p{0.52\linewidth} r l l@{}}
\toprule
\textbf{Dato} & \textbf{Descripción} & \textbf{Valor} & \textbf{Unidad} & \textbf{Ref.}\\
\midrule
$s_{\mathrm{maxc}}$ & Espaciamiento máx. (II) & \cellcolor{fondodato}24 & $\mathrm{in}$ & {\scriptsize } \\
\bottomrule
\end{longtable}

\eqitem{Espaciamiento máximo}{}{S_{\mathrm{max}} = \min\left(s_{\mathrm{maxc}},\ s_{\mathrm{c}}\right) = \mathbf{60.960}\,\mathrm{cm}}

\eqitem{Separación $\le$ máxima}{}{S_{\mathrm{max}} > s_{\mathrm{p}}\;\;\Rightarrow\;\;\text{\ok{OK}}}

\begin{longtable}{@{}l p{0.52\linewidth} r l l@{}}
\toprule
\textbf{Dato} & \textbf{Descripción} & \textbf{Valor} & \textbf{Unidad} & \textbf{Ref.}\\
\midrule
$d_{\mathrm{b}}$ & Diámetro del estribo & \cellcolor{fondodato}1.200 & $\mathrm{cm}$ & {\scriptsize } \\
\bottomrule
\end{longtable}

\eqitem{Área de una barra}{}{A_{\mathrm{b}} = 0.25 \cdot \pi \cdot {d_{\mathrm{b}}}^{2} = \mathbf{1.131}\,\mathrm{cm^{2}}}

\eqitem{Área de estribo (2 ramas)}{}{A_{\mathrm{bt}} = 2 \cdot A_{\mathrm{b}} = \mathbf{2.262}\,\mathrm{cm^{2}}}

\eqitem{Estribo proporcionado $\ge$ requerido}{}{A_{\mathrm{bt}} > A_{\mathrm{vmax}}\;\;\Rightarrow\;\;\text{\ok{OK}}}

\subsection{12. Diseño del cortante horizontal [STD A. 9.20.4]}

\eqitem{Resistencia requerida al corte horizontal}{STD Eq. 9.31a}{V_{\mathrm{nh}} = \dfrac{V_{\mathrm{u}}}{\phi_c} = \mathbf{143.85}\,\mathrm{tonf}}

\eqitem{Caso (a\&b): 80 bv d}{STD Art. 9.20.4.3}{V_{\mathrm{nhab}} = 80\,\mathrm{psi} \cdot b_{\mathrm{ew1}} \cdot d_{\mathrm{2}} = \mathbf{95.687}\,\mathrm{tonf}}

\eqitem{Caso (a\&b): superficie rugosa}{}{V_{\mathrm{nhab}} > V_{\mathrm{nh}}\;\;\Rightarrow\;\;\text{\info{NO OK}}}

\eqitem{Caso (c): 350 bv d}{STD Art. 9.20.4.3}{V_{\mathrm{nhc}} = 350\,\mathrm{psi} \cdot b_{\mathrm{ew1}} \cdot d_{\mathrm{2}} = \mathbf{418.63}\,\mathrm{tonf}}

\eqitem{Cortante horizontal - Caso (c)}{}{V_{\mathrm{nhc}} > V_{\mathrm{nh}}\;\;\Rightarrow\;\;\text{\ok{OK}}}

\eqitem{Mínimo Avh}{STD A. 9.20.4.5}{A_{\mathrm{vh}} = \dfrac{50 \cdot b_{\mathrm{ew1}} \cdot s_{\mathrm{p}}}{\dfrac{f_{\mathrm{y}}}{\mathrm{psi}}} = \mathbf{2.679}\,\mathrm{cm^{2}}}

\eqitem{Refuerzo proporcionado $\ge$ Avh mínimo}{}{A_{\mathrm{bt}} \ge A_{\mathrm{vh}}\;\;\Rightarrow\;\;\text{\nook{NO OK}}}

\eqitem{Espaciamiento máx. 4b}{}{s_{\mathrm{maxvh1}} = 4 \cdot t_{\mathrm{w}} = \mathbf{60.960}\,\mathrm{cm}}

\begin{longtable}{@{}l p{0.52\linewidth} r l l@{}}
\toprule
\textbf{Dato} & \textbf{Descripción} & \textbf{Valor} & \textbf{Unidad} & \textbf{Ref.}\\
\midrule
$s_{\mathrm{maxc2}}$ & Espaciamiento máx. 24 in & \cellcolor{fondodato}24 & $\mathrm{in}$ & {\scriptsize } \\
\bottomrule
\end{longtable}

\eqitem{Espaciamiento máximo}{}{S_{\mathrm{maxvh}} = \min\left(s_{\mathrm{maxvh1}},\ s_{\mathrm{maxc2}}\right) = \mathbf{60.960}\,\mathrm{cm}}

\eqitem{Separación $\le$ máxima (corte horizontal)}{}{S_{\mathrm{maxvh}} > s_{\mathrm{p}}\;\;\Rightarrow\;\;\text{\ok{OK}}}

\subsection{13. Zona de anclaje pretensada [STD A. 9.22]}

\subsubsection{13.1 Mínimo refuerzo vertical [STD A. 9.22.1]}

\begin{longtable}{@{}l p{0.52\linewidth} r l l@{}}
\toprule
\textbf{Dato} & \textbf{Descripción} & \textbf{Valor} & \textbf{Unidad} & \textbf{Ref.}\\
\midrule
$f_s$ & Esfuerzo unitario de los estribos verticales & \cellcolor{fondodato}20 & $\mathrm{ksi}$ & {\scriptsize } \\
\bottomrule
\end{longtable}

\eqitem{Fuerza en torones antes de pérdidas}{}{P_{si} = N_{\mathrm{t}} \cdot A_{\mathrm{st}} \cdot f_{\mathrm{si}} = \mathbf{618.85}\,\mathrm{tonf}}

\eqitem{4\% de la fuerza de pretensado}{}{P_{\mathrm{r}} = 0.04 \cdot P_{si} = \mathbf{24.754}\,\mathrm{tonf}}

\eqitem{Distancia d/4 desde el extremo}{}{d_{\mathrm{025}} = 0.25 \cdot d_{\mathrm{2}} = \mathbf{39.868}\,\mathrm{cm}}

\eqitem{Área requerida de acero}{}{A_{\mathrm{sza}} = \dfrac{P_{\mathrm{r}}}{f_s} = \mathbf{17.604}\,\mathrm{cm^{2}}}

\begin{longtable}{@{}l p{0.52\linewidth} r l l@{}}
\toprule
\textbf{Dato} & \textbf{Descripción} & \textbf{Valor} & \textbf{Unidad} & \textbf{Ref.}\\
\midrule
$N_{\mathrm{bza}}$ & Pares de barras & \cellcolor{fondodato}5 & $$ & {\scriptsize } \\
$d_{\mathrm{bza}}$ & Diámetro de barra & \cellcolor{fondodato}1.600 & $\mathrm{cm}$ & {\scriptsize } \\
\bottomrule
\end{longtable}

\eqitem{Área de una barra}{}{A_{\mathrm{bza}} = \dfrac{\pi}{4} \cdot {d_{\mathrm{bza}}}^{2} = \mathbf{2.011}\,\mathrm{cm^{2}}}

\eqitem{Área total proporcionada}{}{A_{\mathrm{btza}} = N_{\mathrm{bza}} \cdot 2 \cdot A_{\mathrm{bza}} = \mathbf{20.106}\,\mathrm{cm^{2}}}

\eqitem{Refuerzo en zona de anclaje}{}{A_{\mathrm{btza}} > A_{\mathrm{sza}}\;\;\Rightarrow\;\;\text{\ok{OK}}}

\begin{longtable}{@{}l p{0.52\linewidth} r l l@{}}
\toprule
\textbf{Dato} & \textbf{Descripción} & \textbf{Valor} & \textbf{Unidad} & \textbf{Ref.}\\
\midrule
$s_{\mathrm{ep,za}}$ & Espaciamiento & \cellcolor{fondodato}3 & $\mathrm{in}$ & {\scriptsize } \\
$d_{\mathrm{za}}$ & Primera barra desde el extremo & \cellcolor{fondodato}2 & $\mathrm{in}$ & {\scriptsize } \\
\bottomrule
\end{longtable}

\subsubsection{13.2 Refuerzo de confinamiento [STD A. 9.22.2]}

\nota{Para una distancia d desde el extremo de la viga se coloca refuerzo para confinar los torones en el ala inferior.}

\subsection{14. Deflexión y contraflecha}

\subsubsection{14.1 Contraflecha por pretensado en la transferencia}

\eqitem{Diferencia de excentricidad centro - extremo}{}{e' = e_{\mathrm{c}} - \left(y_{\mathrm{b}} - y_{\mathrm{bs,ext}}\right) = \mathbf{31.336}\,\mathrm{cm}}

\eqitem{Distancia del extremo al punto harp}{}{a = l_{\mathrm{lt}} = \mathbf{14.760}\,\mathrm{m}}

\eqitem{Contraflecha por pretensado ($\uparrow$)}{}{\Delta_p = \dfrac{P_{si,i}}{E_{\mathrm{ci}} \cdot I_{\mathrm{g}}} \cdot \left(\dfrac{e_{\mathrm{c}} \cdot {L_{\mathrm{2}}}^{2}}{8} - \dfrac{e' \cdot {a}^{2}}{6}\right) = \mathbf{9.673}\,\mathrm{cm}}

\subsubsection{14.2 Deflexión por peso de la viga}

\eqitem{En la transferencia ($\downarrow$)}{}{\Delta_g = \dfrac{5 \cdot W_{\mathrm{g}} \cdot {L_{\mathrm{2}}}^{4}}{384 \cdot E_{\mathrm{ci}} \cdot I_{\mathrm{g}}} = \mathbf{4.034}\,\mathrm{cm}}

\eqitem{En el izaje ($\downarrow$)}{}{\Delta_{gi} = \dfrac{5 \cdot W_{\mathrm{g}} \cdot {L_{\mathrm{1}}}^{4}}{384 \cdot E_{\mathrm{ci}} \cdot I_{\mathrm{g}}} = \mathbf{3.904}\,\mathrm{cm}}

\subsubsection{14.3 Resumen}

\eqitem{En la transferencia ($\uparrow$)}{}{\Delta_{pg} = \Delta_p - \Delta_g = \mathbf{5.639}\,\mathrm{cm}}

\eqitem{Total en el izaje, factores PCI ($\uparrow$)}{}{\Delta_{pg,izaje} = 1.8 \cdot \Delta_p - 1.85 \cdot \Delta_{gi} = \mathbf{10.188}\,\mathrm{cm}}

\subsubsection{14.4 Deflexión por peso de losa y cartela}

\eqitem{Deflexión ($\downarrow$)}{}{\Delta_s = \dfrac{5 \cdot W_{\mathrm{s}} \cdot {L_{\mathrm{1}}}^{4}}{384 \cdot E_{\mathrm{cs}} \cdot I_{\mathrm{g}}} = \mathbf{4.065}\,\mathrm{cm}}

\subsubsection{14.5 Deflexión por barreras y capa de rodadura}

\eqitem{Deflexión ($\downarrow$)}{}{\Delta_{b+ws} = \dfrac{5 \cdot \left(\mathit{WDC}_{\mathrm{b}} + \mathit{WDW}\right) \cdot {L_{\mathrm{1}}}^{4}}{384 \cdot E_{\mathrm{cs}} \cdot I_{\mathrm{c}}} = \mathbf{0.6633}\,\mathrm{cm}}

\subsubsection{14.6 Deflexión por carga viva e impacto}

\eqitem{Deflexión permitida L/800}{}{\Delta_{LL,adm} = \dfrac{L_{\mathrm{1}}}{800} = \mathbf{4.575}\,\mathrm{cm}}

\eqitem{Número de carriles}{}{\mathit{NL} = \left\lfloor \dfrac{w}{3.65\,\mathrm{m}}\right\rfloor = \mathbf{4}}

\eqitem{Factor de distribución para deflexión}{}{\mathit{FDF} = \dfrac{\mathit{NL}}{N_{\mathrm{b}}} = \mathbf{0.6667}}

\nota{Es más conservador usar el factor de distribución para momento, FDM.}

\begin{longtable}{@{}l p{0.52\linewidth} r l l@{}}
\toprule
\textbf{Dato} & \textbf{Descripción} & \textbf{Valor} & \textbf{Unidad} & \textbf{Ref.}\\
\midrule
$L_{\mathrm{2cd}}$ & Separación entre ejes del camión & \cellcolor{fondodato}4.268 & $\mathrm{m}$ & {\scriptsize } \\
\bottomrule
\end{longtable}

\eqitem{Deflexión por eje delantero (P/4) al centro}{}{\Delta_{P2} = \dfrac{0.25 \cdot P \cdot {L_{\mathrm{1}}}^{3}}{48 \cdot E_{\mathrm{cs}} \cdot I_{\mathrm{c}}} = \mathbf{0.2383}\,\mathrm{cm}}

\eqitem{Posición de los ejes traseros}{}{a_{CD1} = \dfrac{L_{\mathrm{1}} - L_{\mathrm{2cd}}}{2} = \mathbf{16.166}\,\mathrm{m}}

\eqitem{Deflexión por ejes traseros}{}{\Delta_{P1} = \dfrac{P \cdot a_{CD1}}{24 \cdot E_{\mathrm{cs}} \cdot I_{\mathrm{c}}} \cdot \left(3 \cdot {L_{\mathrm{1}}}^{2} - 4 \cdot {a_{CD1}}^{2}\right) = \mathbf{1.869}\,\mathrm{cm}}

\eqitem{Deflexión total del camión}{}{\Delta_{PT} = \Delta_{P1} + \Delta_{P2} = \mathbf{2.108}\,\mathrm{cm}}

\eqitem{Factor de impacto}{}{\mathit{FIM} = 1 + I = \mathbf{1.204}}

\eqitem{Deflexión por CV + impacto}{}{\Delta_{LT} = \Delta_{PT} \cdot \mathit{FIM} \cdot \mathit{FDM} = \mathbf{2.081}\,\mathrm{cm}}

\eqitem{Deflexión por carga viva $\le$ L/800}{}{\Delta_{LT} < \Delta_{LL,adm}\;\;\Rightarrow\;\;\text{\ok{OK}}}

\clearpage
\section{Viga interior preforzada - AASHTO LRFD}

\nota{Puente tipo losa-vigas preforzadas. Viga AASHTO-PCI Bulb-Tee BT-72. Diseño según especificaciones AASHTO LRFD (carga HL-93: camión HS20-44 + carril).}

\subsection{2. Materiales}

\begin{longtable}{@{}l p{0.52\linewidth} r l l@{}}
\toprule
\textbf{Dato} & \textbf{Descripción} & \textbf{Valor} & \textbf{Unidad} & \textbf{Ref.}\\
\midrule
$t_{\mathrm{s}}$ & Espesor de la losa interior & \cellcolor{fondodato}19 & $\mathrm{cm}$ & {\scriptsize } \\
$t_{\mathrm{ws}}$ & Sobre-espesor de desgaste & \cellcolor{fondodato}1 & $\mathrm{cm}$ & {\scriptsize } \\
\bottomrule
\end{longtable}

\eqitem{Espesor total}{}{t = t_{\mathrm{s}} + t_{\mathrm{ws}} = \mathbf{20}\,\mathrm{cm}}

\begin{longtable}{@{}l p{0.52\linewidth} r l l@{}}
\toprule
\textbf{Dato} & \textbf{Descripción} & \textbf{Valor} & \textbf{Unidad} & \textbf{Ref.}\\
\midrule
$t_o$ & Espesor de la losa en el voladizo & \cellcolor{fondodato}20 & $\mathrm{cm}$ & {\scriptsize } \\
$f'_{c,d1}$ & f'c de diseño de la losa & \cellcolor{fondodato}280 & $\mathrm{kgf/cm^{2}}$ & {\scriptsize } \\
$f'_{c,d2}$ & f'c de diseño de la viga & \cellcolor{fondodato}450 & $\mathrm{kgf/cm^{2}}$ & {\scriptsize } \\
$f'_{ci}$ & f'c de la viga en la transferencia & \cellcolor{fondodato}405 & $\mathrm{kgf/cm^{2}}$ & {\scriptsize } \\
$W_{\mathrm{c}}$ & Peso unitario del hormigón & \cellcolor{fondodato}2.400 & $\mathrm{tonf/m^{3}}$ & {\scriptsize } \\
$L_{\mathrm{2}}$ & Longitud total de la viga & \cellcolor{fondodato}36.900 & $\mathrm{m}$ & {\scriptsize } \\
$L_{\mathrm{1}}$ & Tramo de diseño & \cellcolor{fondodato}36.600 & $\mathrm{m}$ & {\scriptsize } \\
$d_t$ & Diámetro del torón (7 alambres, baja relajación) & \cellcolor{fondodato}0.5000 & $\mathrm{in}$ & {\scriptsize } \\
$f_{\mathrm{pu}}$ & Resistencia a la tensión del torón & \cellcolor{fondodato}19 & $\mathrm{tonf/cm^{2}}$ & {\scriptsize } \\
$A_{\mathrm{st}}$ & Área de un torón & \cellcolor{fondodato}98.700 & $\mathrm{mm^{2}}$ & {\scriptsize } \\
\bottomrule
\end{longtable}

\eqitem{Resistencia a la fluencia}{LRFD T. 5.4.4.1-1}{f_{\mathrm{py}} = 0.9 \cdot f_{\mathrm{pu}} = \mathbf{17.100}\,\mathrm{tonf/cm^{2}}}

\nota{Esfuerzos límites para los torones [LRFD T. 5.9.3-1].}

\eqitem{Antes de la transferencia (fpi $\le$ 0.75 fpu)}{}{f_{\mathrm{pi}} = 0.75 \cdot f_{\mathrm{pu}} = \mathbf{14.250}\,\mathrm{tonf/cm^{2}}}

\eqitem{Estado límite de servicio (después de pérdidas)}{}{f_{pe,lim} = 0.8 \cdot f_{\mathrm{py}} = \mathbf{13.680}\,\mathrm{tonf/cm^{2}}}

\begin{longtable}{@{}l p{0.52\linewidth} r l l@{}}
\toprule
\textbf{Dato} & \textbf{Descripción} & \textbf{Valor} & \textbf{Unidad} & \textbf{Ref.}\\
\midrule
$E_{\mathrm{p}}$ & Módulo de elasticidad del torón & \cellcolor{fondodato}2\,003.75 & $\mathrm{tonf/cm^{2}}$ & {\scriptsize LRFD T. 5.4.4.2} \\
$f_{\mathrm{y}}$ & Fluencia del acero de refuerzo & \cellcolor{fondodato}4.200 & $\mathrm{tonf/cm^{2}}$ & {\scriptsize } \\
$E_{\mathrm{s}}$ & Módulo de elasticidad del acero de refuerzo & \cellcolor{fondodato}2\,038.90 & $\mathrm{tonf/cm^{2}}$ & {\scriptsize LRFD T. 5.4.3.2} \\
$f_{\mathrm{ws}}$ & Espesor de la capa de rodadura & \cellcolor{fondodato}5 & $\mathrm{cm}$ & {\scriptsize } \\
$W_{\mathrm{b}}$ & Peso de la barrera New Jersey (por lado) & \cellcolor{fondodato}0.4465 & $\mathrm{tonf/m}$ & {\scriptsize } \\
\bottomrule
\end{longtable}

\subsection{3. Propiedades de la sección transversal de una viga interior}

\subsubsection{3.1 Sección no compuesta}

\begin{longtable}{@{}l p{0.52\linewidth} r l l@{}}
\toprule
\textbf{Dato} & \textbf{Descripción} & \textbf{Valor} & \textbf{Unidad} & \textbf{Ref.}\\
\midrule
$A_{\mathrm{g}}$ & Área de la viga & \cellcolor{fondodato}0.4948 & $\mathrm{m^{2}}$ & {\scriptsize } \\
$h_{\mathrm{g}}$ & Altura de la viga & \cellcolor{fondodato}182.88 & $\mathrm{cm}$ & {\scriptsize } \\
$I_{\mathrm{g}}$ & Inercia de la viga & \cellcolor{fondodato}0.2272 & $\mathrm{m^{4}}$ & {\scriptsize } \\
$y_{\mathrm{b}}$ & CG - fibra inferior & \cellcolor{fondodato}92.964 & $\mathrm{cm}$ & {\scriptsize } \\
\bottomrule
\end{longtable}

\eqitem{CG - fibra superior}{}{y_{\mathrm{t}} = h_{\mathrm{g}} - y_{\mathrm{b}} = \mathbf{89.916}\,\mathrm{cm}}

\eqitem{Módulo de sección, fibra inferior}{}{S_{\mathrm{b}} = \dfrac{I_{\mathrm{g}}}{y_{\mathrm{b}}} = \mathbf{0.2444}\,\mathrm{m^{3}}}

\eqitem{Módulo de sección, fibra superior}{}{S_{\mathrm{t}} = \dfrac{I_{\mathrm{g}}}{y_{\mathrm{t}}} = \mathbf{0.2527}\,\mathrm{m^{3}}}

\eqitem{Peso de la viga}{}{W_{\mathrm{g}} = A_{\mathrm{g}} \cdot W_{\mathrm{c}} = \mathbf{1.188}\,\mathrm{tonf/m}}

\eqitem{M.E. de la losa}{LRFD Eq. 5.4.2.4-1}{E_{\mathrm{c}} = 33000 \cdot {\left(\dfrac{W_{\mathrm{c}}}{\dfrac{\mathrm{kip}}{{\mathrm{ft}}^{3}}}\right)}^{1.5} \cdot \sqrt{\dfrac{f'_{c,d1}}{\mathrm{ksi}}} \cdot \mathrm{ksi} = \mathbf{268.52}\,\mathrm{tonf/cm^{2}}}

\eqitem{M.E. de la viga en la transferencia}{}{E_{\mathrm{ci}} = 33000 \cdot {\left(\dfrac{W_{\mathrm{c}}}{\dfrac{\mathrm{kip}}{{\mathrm{ft}}^{3}}}\right)}^{1.5} \cdot \sqrt{\dfrac{f'_{ci}}{\mathrm{ksi}}} \cdot \mathrm{ksi} = \mathbf{322.94}\,\mathrm{tonf/cm^{2}}}

\eqitem{M.E. de la viga (cargas de servicio)}{}{E_{\mathrm{cs}} = 33000 \cdot {\left(\dfrac{W_{\mathrm{c}}}{\dfrac{\mathrm{kip}}{{\mathrm{ft}}^{3}}}\right)}^{1.5} \cdot \sqrt{\dfrac{f'_{c,d2}}{\mathrm{ksi}}} \cdot \mathrm{ksi} = \mathbf{340.41}\,\mathrm{tonf/cm^{2}}}

\subsubsection{3.2 Sección compuesta. Ancho efectivo del ala [LRFD 4.6.2.6.1]}

\eqitem{1/4 del tramo de diseño}{}{b_{\mathrm{ea1}} = 0.25 \cdot L_{\mathrm{1}} = \mathbf{915}\,\mathrm{cm}}

\begin{longtable}{@{}l p{0.52\linewidth} r l l@{}}
\toprule
\textbf{Dato} & \textbf{Descripción} & \textbf{Valor} & \textbf{Unidad} & \textbf{Ref.}\\
\midrule
$S$ & Separación entre vigas (promedio) & \cellcolor{fondodato}275 & $\mathrm{cm}$ & {\scriptsize } \\
\bottomrule
\end{longtable}

\eqitem{Separación entre vigas}{}{b_{\mathrm{ea2}} = S = \mathbf{275}\,\mathrm{cm}}

\begin{longtable}{@{}l p{0.52\linewidth} r l l@{}}
\toprule
\textbf{Dato} & \textbf{Descripción} & \textbf{Valor} & \textbf{Unidad} & \textbf{Ref.}\\
\midrule
$t_{\mathrm{w}}$ & Espesor del alma & \cellcolor{fondodato}15.240 & $\mathrm{cm}$ & {\scriptsize } \\
$b_{\mathrm{ew1}}$ & Ancho del ala superior de la viga & \cellcolor{fondodato}106.68 & $\mathrm{cm}$ & {\scriptsize } \\
\bottomrule
\end{longtable}

\eqitem{Mayor de: tw ó 1/2 ancho del ala superior}{}{b_{\mathrm{ew}} = \max\left(0.5 \cdot b_{\mathrm{ew1}},\ t_{\mathrm{w}}\right) = \mathbf{53.340}\,\mathrm{cm}}

\eqitem{12 ts + bew}{}{b_{\mathrm{ea3}} = 12 \cdot t_{\mathrm{s}} + b_{\mathrm{ew}} = \mathbf{281.34}\,\mathrm{cm}}

\eqitem{Ancho de ala efectivo}{}{w_{\mathrm{fe}} = \min\left(b_{\mathrm{ea1}},\ b_{\mathrm{ea2}},\ b_{\mathrm{ea3}}\right) = \mathbf{275}\,\mathrm{cm}}

\eqitem{Razón modular losa/viga}{}{n_{dg} = \dfrac{E_{\mathrm{c}}}{E_{\mathrm{cs}}} = \mathbf{0.7888}}

\subsubsection{3.2.3 Propiedades de la sección transformada}

\eqitem{Ancho de ala transformado}{}{w_{\mathrm{ft}} = n_{dg} \cdot w_{\mathrm{fe}} = \mathbf{2.169}\,\mathrm{m}}

\eqitem{Área de ala transformada}{}{a_{\mathrm{ft}} = w_{\mathrm{ft}} \cdot t_{\mathrm{s}} = \mathbf{0.4122}\,\mathrm{m^{2}}}

\eqitem{Inercia propia del ala}{}{I_{\mathrm{ft}} = \dfrac{w_{\mathrm{ft}} \cdot {t_{\mathrm{s}}}^{3}}{12} = \mathbf{1.240\times10^{-3}}\,\mathrm{m^{4}}}

\begin{longtable}{@{}l p{0.52\linewidth} r l l@{}}
\toprule
\textbf{Dato} & \textbf{Descripción} & \textbf{Valor} & \textbf{Unidad} & \textbf{Ref.}\\
\midrule
$t_{\mathrm{h}}$ & Espesor mínimo de cartela (centro del tramo) & \cellcolor{fondodato}1 & $\mathrm{cm}$ & {\scriptsize } \\
$w_{\mathrm{fg}}$ & Ancho del ala superior (viga) & \cellcolor{fondodato}106.68 & $\mathrm{cm}$ & {\scriptsize } \\
\bottomrule
\end{longtable}

\eqitem{Ancho de cartela transformado}{}{w_{\mathrm{th}} = n_{dg} \cdot w_{\mathrm{fg}} = \mathbf{0.8415}\,\mathrm{m}}

\eqitem{Área de cartela transformada}{}{a_{\mathrm{th}} = w_{\mathrm{th}} \cdot t_{\mathrm{h}} = \mathbf{8.415\times10^{-3}}\,\mathrm{m^{2}}}

\eqitem{Inercia de cartela}{}{I_{\mathrm{at}} = \dfrac{w_{\mathrm{th}} \cdot {t_{\mathrm{h}}}^{3}}{12} = \mathbf{7.013}\,\mathrm{cm^{4}}}

\begin{center}
{\small\textbf{Propiedades de la sección compuesta (parte 1)}}\\[3pt]
\adjustbox{max width=\linewidth}{\begin{tabular}{ccccc}
\toprule
\begin{tabular}[b]{@{}c@{}}\textbf{Elemento}\end{tabular} & \begin{tabular}[b]{@{}c@{}}\textbf{Área}\\ {\scriptsize $\mathrm{m^{2}}$}\end{tabular} & \begin{tabular}[b]{@{}c@{}}\textbf{yb}\\ {\scriptsize $\mathrm{m}$}\end{tabular} & \begin{tabular}[b]{@{}c@{}}\textbf{A$\cdot$yb}\\ {\scriptsize $\mathrm{m^{3}}$}\end{tabular} & \begin{tabular}[b]{@{}c@{}}\textbf{I propia}\\ {\scriptsize $\mathrm{m^{4}}$}\end{tabular} \\
\midrule
Viga & $0.4948$ & $0.9296$ & $0.4600$ & $0.2272$ \\
Cartela & $8.415\times10^{-3}$ & $1.834$ & $0.0154$ & $7.013\times10^{-8}$ \\
Losa & $0.4122$ & $1.934$ & $0.7970$ & $1.240\times10^{-3}$ \\
\bottomrule
\end{tabular}}
\end{center}

\eqitem{Área total de la sección compuesta}{}{A_{\mathrm{c}} = A_{\mathrm{v}} + A_{\mathrm{h}} + A_{\mathrm{l}} = \mathbf{0.9154}\,\mathrm{m^{2}}}

\eqitem{CG de la sección compuesta - fibra inferior de la viga}{}{y_{\mathrm{bc}} = \dfrac{A_{\mathrm{y,v}} + A_{\mathrm{y,h}} + A_{\mathrm{y,l}}}{A_{\mathrm{c}}} = \mathbf{1.390}\,\mathrm{m}}

\begin{center}
{\small\textbf{Propiedades de la sección compuesta (parte 2: Steiner)}}\\[3pt]
\adjustbox{max width=\linewidth}{\begin{tabular}{ccc}
\toprule
\begin{tabular}[b]{@{}c@{}}\textbf{Elemento}\end{tabular} & \begin{tabular}[b]{@{}c@{}}\textbf{A(ybc - yb)$^{2}$}\\ {\scriptsize $\mathrm{m^{4}}$}\end{tabular} & \begin{tabular}[b]{@{}c@{}}\textbf{I + A(ybc - yb)$^{2}$}\\ {\scriptsize $\mathrm{m^{4}}$}\end{tabular} \\
\midrule
Viga & $0.1049$ & $0.3321$ \\
Cartela & $1.657\times10^{-3}$ & $1.657\times10^{-3}$ \\
Losa & $0.1218$ & $0.1231$ \\
\bottomrule
\end{tabular}}
\end{center}

\eqitem{Inercia de la sección compuesta}{}{I_{\mathrm{c}} = I_{\mathrm{t,v}} + I_{\mathrm{t,h}} + I_{\mathrm{t,l}} = \mathbf{0.4569}\,\mathrm{m^{4}}}

\eqitem{Altura total de la sección compuesta}{}{h_{\mathrm{sc}} = h_{\mathrm{g}} + t_{\mathrm{h}} + t_{\mathrm{s}} = \mathbf{2.029}\,\mathrm{m}}

\eqitem{CG - fibra superior de la viga}{}{y_{\mathrm{tg}} = h_{\mathrm{g}} - y_{\mathrm{bc}} = \mathbf{43.872}\,\mathrm{cm}}

\eqitem{CG - fibra superior de la losa}{}{y_{\mathrm{tc}} = h_{\mathrm{sc}} - y_{\mathrm{bc}} = \mathbf{63.872}\,\mathrm{cm}}

\eqitem{Módulo de sección compuesta, fibra inferior viga}{}{S_{\mathrm{bc}} = \dfrac{I_{\mathrm{c}}}{y_{\mathrm{bc}}} = \mathbf{0.3287}\,\mathrm{m^{3}}}

\eqitem{Módulo de sección compuesta, fibra superior viga}{}{S_{\mathrm{tg}} = \dfrac{I_{\mathrm{c}}}{y_{\mathrm{tg}}} = \mathbf{1.041}\,\mathrm{m^{3}}}

\eqitem{Módulo de sección compuesta, fibra superior losa}{}{S_{\mathrm{tc}} = \dfrac{1}{n_{dg}} \cdot \dfrac{I_{\mathrm{c}}}{y_{\mathrm{tc}}} = \mathbf{0.9068}\,\mathrm{m^{3}}}

\subsection{4. Fuerzas cortantes y momentos flectores}

\subsubsection{4.1 Cargas muertas [LRFD Art. 3.3.2]}

\eqitem{Peso de la losa}{}{W_{\mathrm{d}} = t \cdot S \cdot W_{\mathrm{c}} = \mathbf{1.320}\,\mathrm{tonf/m}}

\eqitem{Peso de la cartela}{}{W_{\mathrm{ha}} = t_{\mathrm{h}} \cdot w_{\mathrm{fg}} \cdot W_{\mathrm{c}} = \mathbf{0.0256}\,\mathrm{tonf/m}}

\eqitem{Peso de losa + cartela}{}{W_{\mathrm{s}} = W_{\mathrm{d}} + W_{\mathrm{ha}} = \mathbf{1.346}\,\mathrm{tonf/m}}

\nota{LRFD A. 4.6.2.2.1: las CM de barreras y carpeta se distribuyen uniformemente entre todas las vigas si: ancho de losa constante; Nb $\ge$ 4; voladizo $\le$ 3.0 ft; vigas paralelas de igual rigidez; curvatura en planta < 4$^{\circ}$; sección de la Tabla 4.6.2.2.1-1.}

\begin{longtable}{@{}l p{0.52\linewidth} r l l@{}}
\toprule
\textbf{Dato} & \textbf{Descripción} & \textbf{Valor} & \textbf{Unidad} & \textbf{Ref.}\\
\midrule
$N_{\mathrm{b}}$ & Número de vigas & \cellcolor{fondodato}6 & $$ & {\scriptsize } \\
\bottomrule
\end{longtable}

\eqitem{Número de vigas Nb > 4}{}{N_{\mathrm{b}} > 4\;\;\Rightarrow\;\;\text{\ok{OK}}}

\begin{longtable}{@{}l p{0.52\linewidth} r l l@{}}
\toprule
\textbf{Dato} & \textbf{Descripción} & \textbf{Valor} & \textbf{Unidad} & \textbf{Ref.}\\
\midrule
$L_o$ & Longitud del voladizo & \cellcolor{fondodato}91.440 & $\mathrm{cm}$ & {\scriptsize } \\
$L_{\mathrm{bar}}$ & Ancho de la barrera en la base & \cellcolor{fondodato}45.720 & $\mathrm{cm}$ & {\scriptsize } \\
\bottomrule
\end{longtable}

\eqitem{Parte de la calzada en el voladizo}{}{d_e = L_o - L_{\mathrm{bar}} = \mathbf{1.500}\,\mathrm{ft}}

\eqitem{Calzada en el voladizo $\le$ 3.0 ft}{}{d_e \le 3\,\mathrm{ft}\;\;\Rightarrow\;\;\text{\ok{OK}}}

\eqitem{Peso de las 2 barreras por viga}{}{\mathit{WDC}_{\mathrm{b}} = \dfrac{W_{\mathrm{b}} \cdot 2}{N_{\mathrm{b}}} = \mathbf{0.1488}\,\mathrm{tonf/m}}

\begin{longtable}{@{}l p{0.52\linewidth} r l l@{}}
\toprule
\textbf{Dato} & \textbf{Descripción} & \textbf{Valor} & \textbf{Unidad} & \textbf{Ref.}\\
\midrule
$w$ & Ancho de capa de rodadura & \cellcolor{fondodato}14.634 & $\mathrm{m}$ & {\scriptsize } \\
\bottomrule
\end{longtable}

\eqitem{Peso de carpeta asfáltica por viga}{}{\mathit{WDW} = \dfrac{W_{\mathrm{c}} \cdot f_{\mathrm{ws}} \cdot w}{N_{\mathrm{b}}} = \mathbf{0.2927}\,\mathrm{tonf/m}}

\begin{center}
{\small\textbf{Cortantes (tonf) y momentos (tonf$\cdot$m) no mayorados por CM - viga interior}}\\[3pt]
\adjustbox{max width=\linewidth}{\begin{tabular}{cccccccccc}
\toprule
\begin{tabular}[b]{@{}c@{}}\textbf{x/L}\end{tabular} & \begin{tabular}[b]{@{}c@{}}\textbf{x}\\ {\scriptsize $\mathrm{m}$}\end{tabular} & \begin{tabular}[b]{@{}c@{}}\textbf{Vg}\\ {\scriptsize $\mathrm{tonf}$}\end{tabular} & \begin{tabular}[b]{@{}c@{}}\textbf{Mg}\\ {\scriptsize $\mathrm{tonf\cdot m}$}\end{tabular} & \begin{tabular}[b]{@{}c@{}}\textbf{Vs}\\ {\scriptsize $\mathrm{tonf}$}\end{tabular} & \begin{tabular}[b]{@{}c@{}}\textbf{Ms}\\ {\scriptsize $\mathrm{tonf\cdot m}$}\end{tabular} & \begin{tabular}[b]{@{}c@{}}\textbf{Vb}\\ {\scriptsize $\mathrm{tonf}$}\end{tabular} & \begin{tabular}[b]{@{}c@{}}\textbf{Mb}\\ {\scriptsize $\mathrm{tonf\cdot m}$}\end{tabular} & \begin{tabular}[b]{@{}c@{}}\textbf{Vws}\\ {\scriptsize $\mathrm{tonf}$}\end{tabular} & \begin{tabular}[b]{@{}c@{}}\textbf{Mws}\\ {\scriptsize $\mathrm{tonf\cdot m}$}\end{tabular} \\
\midrule
\cellcolor{fondodato}$0$ & $0$ & $21.732$ & $0$ & $24.625$ & $0$ & $2.724$ & $0$ & $5.356$ & $0$ \\
\cellcolor{fondodato}$0.0600$ & $2.196$ & $19.124$ & $44.859$ & $21.670$ & $50.831$ & $2.397$ & $5.622$ & $4.713$ & $11.056$ \\
\cellcolor{fondodato}$0.1000$ & $3.660$ & $17.385$ & $71.584$ & $19.700$ & $81.113$ & $2.179$ & $8.972$ & $4.285$ & $17.643$ \\
\cellcolor{fondodato}$0.2000$ & $7.320$ & $13.039$ & $127.26$ & $14.775$ & $144.20$ & $1.634$ & $15.950$ & $3.214$ & $31.365$ \\
\cellcolor{fondodato}$0.3000$ & $10.980$ & $8.693$ & $167.03$ & $9.850$ & $189.26$ & $1.089$ & $20.934$ & $2.142$ & $41.167$ \\
\cellcolor{fondodato}$0.4000$ & $14.640$ & $4.346$ & $190.89$ & $4.925$ & $216.30$ & $0.5447$ & $23.925$ & $1.071$ & $47.047$ \\
\cellcolor{fondodato}$0.5000$ & $18.300$ & $0$ & $198.84$ & $0$ & $225.31$ & $0$ & $24.921$ & $0$ & $49.008$ \\
\bottomrule
\end{tabular}}
\end{center}

\subsubsection{4.2 Carga viva [LRFD A. 3.6.1.2]}

\nota{Factores de distribución de carga viva para una viga interior [LRFD A. 4.6.2.2]; sección transversal tipo (k) [LRFD T. 4.6.2.2.1-1].}

\eqitem{Número de carriles de diseño}{LRFD A. 3.6.1.1.1}{\mathit{NL} = \left\lfloor \dfrac{w}{12\,\mathrm{ft}}\right\rfloor = \mathbf{4}}

\eqitem{Separación entre vigas}{}{S = S = \mathbf{9.022}\,\mathrm{ft}}

\eqitem{(a) 3.5 $\le$ S $\le$ 16.0 ft}{}{S \ge 3.5\,\mathrm{ft} \;\wedge\; S \le 16\,\mathrm{ft}\;\;\Rightarrow\;\;\text{\ok{OK}}}

\eqitem{Espesor de losa}{}{t_s = t_{\mathrm{s}} = \mathbf{7.480}\,\mathrm{in}}

\eqitem{(b) 4.5 $\le$ ts $\le$ 12.0 in}{}{t_{\mathrm{s}} \ge 4.5\,\mathrm{in} \;\wedge\; t_{\mathrm{s}} \le 12\,\mathrm{in}\;\;\Rightarrow\;\;\text{\ok{OK}}}

\eqitem{Tramo}{}{L_1 = L_{\mathrm{1}} = \mathbf{120.08}\,\mathrm{ft}}

\eqitem{(c) 20 $\le$ L $\le$ 240 ft}{}{L_{\mathrm{1}} \ge 20\,\mathrm{ft} \;\wedge\; L_{\mathrm{1}} \le 240\,\mathrm{ft}\;\;\Rightarrow\;\;\text{\ok{OK}}}

\eqitem{(d) Nb $\ge$ 4}{}{N_{\mathrm{b}} \ge 4\;\;\Rightarrow\;\;\text{\ok{OK}}}

\eqitem{Razón modular viga/losa}{}{n_c = \dfrac{E_{\mathrm{cs}}}{E_{\mathrm{c}}} = \mathbf{1.268}}

\eqitem{Distancia entre CG de viga y losa}{}{e_{\mathrm{g}} = \dfrac{t_{\mathrm{s}}}{2} + t_{\mathrm{h}} + y_{\mathrm{t}} = \mathbf{39.534}\,\mathrm{in}}

\eqitem{Parámetro de rigidez longitudinal}{LRFD Eq. 4.6.2.2.1-1}{K_{\mathrm{g}} = n_c \cdot \left(I_{\mathrm{g}} + A_{\mathrm{g}} \cdot {e_{\mathrm{g}}}^{2}\right) = \mathbf{2\,211\,647.48}\,\mathrm{in^{4}}}

\eqitem{(e) 10 000 $\le$ Kg $\le$ 7 000 000 in$^{4}$}{}{K_{\mathrm{g}} \ge 10000 \cdot {\mathrm{in}}^{4} \;\wedge\; K_{\mathrm{g}} \le 7000000 \cdot {\mathrm{in}}^{4}\;\;\Rightarrow\;\;\text{\ok{OK}}}

\eqitem{Distribución para momento: 1 carril de diseño}{LRFD T. 4.6.2.2.2b-1}{\mathit{DFM}_{\mathrm{1}} = 0.06 + {\left(\dfrac{S}{14\,\mathrm{ft}}\right)}^{0.4} \cdot {\left(\dfrac{S}{L_{\mathrm{1}}}\right)}^{0.3} \cdot {\left(\dfrac{\dfrac{K_{\mathrm{g}}}{{\mathrm{in}}^{4}}}{12 \cdot \dfrac{L_{\mathrm{1}}}{\mathrm{ft}} \cdot {\left(\dfrac{t_{\mathrm{s}}}{\mathrm{in}}\right)}^{3}}\right)}^{0.1} = \mathbf{0.4994}}

\eqitem{Distribución para momento: 2 o más carriles}{LRFD T. 4.6.2.2.2b-1}{\mathit{DFM}_{\mathrm{2}} = 0.075 + {\left(\dfrac{S}{9.5\,\mathrm{ft}}\right)}^{0.6} \cdot {\left(\dfrac{S}{L_{\mathrm{1}}}\right)}^{0.2} \cdot {\left(\dfrac{\dfrac{K_{\mathrm{g}}}{{\mathrm{in}}^{4}}}{12 \cdot \dfrac{L_{\mathrm{1}}}{\mathrm{ft}} \cdot {\left(\dfrac{t_{\mathrm{s}}}{\mathrm{in}}\right)}^{3}}\right)}^{0.1} = \mathbf{0.7329}}

\eqitem{Factor de distribución para momento (controla)}{}{\mathit{FDM}_{\mathrm{imax}} = \max\left(\mathit{DFM}_{\mathrm{1}},\ \mathit{DFM}_{\mathrm{2}}\right) = \mathbf{0.7329}}

\begin{longtable}{@{}l p{0.52\linewidth} r l l@{}}
\toprule
\textbf{Dato} & \textbf{Descripción} & \textbf{Valor} & \textbf{Unidad} & \textbf{Ref.}\\
\midrule
$\mathit{FPM}$ & Factor de presencia múltiple & \cellcolor{fondodato}1.200 & $$ & {\scriptsize LRFD A. 3.6.1.1.2} \\
\bottomrule
\end{longtable}

\eqitem{Factor de distribución para fatiga}{}{\mathit{DFM}_{\mathrm{f}} = \dfrac{\mathit{DFM}_{\mathrm{1}}}{\mathit{FPM}} = \mathbf{0.4162}}

\eqitem{Distribución para cortante: 1 carril}{LRFD T. 4.6.2.2.3a-1}{\mathit{DFV}_{\mathrm{1}} = 0.36 + \dfrac{S}{25\,\mathrm{ft}} = \mathbf{0.7209}}

\eqitem{Distribución para cortante: 2 o más carriles}{}{\mathit{DFV}_{\mathrm{2}} = 0.2 + \dfrac{S}{12\,\mathrm{ft}} - {\left(\dfrac{S}{35\,\mathrm{ft}}\right)}^{2} = \mathbf{0.8854}}

\eqitem{Factor de distribución para cortante (controla)}{}{\mathit{FDV}_{\mathrm{imax}} = \max\left(\mathit{DFV}_{\mathrm{1}},\ \mathit{DFV}_{\mathrm{2}}\right) = \mathbf{0.8854}}

\nota{Amplificación dinámica IM: solo se aplica a la carga del camión [LRFD T. 3.6.2.1-1].}

\begin{longtable}{@{}l p{0.52\linewidth} r l l@{}}
\toprule
\textbf{Dato} & \textbf{Descripción} & \textbf{Valor} & \textbf{Unidad} & \textbf{Ref.}\\
\midrule
$IM_f$ & IM para el estado límite de fatiga & \cellcolor{fondodato}15 & $\%$ & {\scriptsize } \\
$\mathit{IM}$ & IM para los demás estados límite & \cellcolor{fondodato}33 & $\%$ & {\scriptsize } \\
$P$ & Carga por eje del camión HS20-44 & \cellcolor{fondodato}32 & $\mathrm{kip}$ & {\scriptsize } \\
$l_1$ & Separación ejes delantero-central & \cellcolor{fondodato}14 & $\mathrm{ft}$ & {\scriptsize } \\
$l_2$ & Separación ejes central-trasero (camión de diseño) & \cellcolor{fondodato}14 & $\mathrm{ft}$ & {\scriptsize } \\
$l_{2,f}$ & Separación ejes central-trasero (camión de fatiga) & \cellcolor{fondodato}30 & $\mathrm{ft}$ & {\scriptsize } \\
$\mathit{WLL}$ & Carga de carril de diseño & \cellcolor{fondodato}0.6400 & $\mathrm{kip/ft}$ & {\scriptsize } \\
\bottomrule
\end{longtable}

\nota{VLT = V camión$\cdot$DFV$\cdot$(1+IM); MLT = M camión$\cdot$DFM$\cdot$(1+IM); VLL = V carril$\cdot$DFV; MLL = M carril$\cdot$DFM; Mf = M camión de fatiga$\cdot$DFMf$\cdot$(1+IMf). Cortante máximo de carril: w(L-x)$^{2}$/(2L).}

\begin{center}
{\small\textbf{Cortantes (tonf) y momentos (tonf$\cdot$m) no mayorados por CV - viga interior}}\\[3pt]
\adjustbox{max width=\linewidth}{\begin{tabular}{ccccccc}
\toprule
\begin{tabular}[b]{@{}c@{}}\textbf{x/L}\end{tabular} & \begin{tabular}[b]{@{}c@{}}\textbf{x}\\ {\scriptsize $\mathrm{m}$}\end{tabular} & \begin{tabular}[b]{@{}c@{}}\textbf{VLT}\\ {\scriptsize $\mathrm{tonf}$}\end{tabular} & \begin{tabular}[b]{@{}c@{}}\textbf{MLT}\\ {\scriptsize $\mathrm{tonf\cdot m}$}\end{tabular} & \begin{tabular}[b]{@{}c@{}}\textbf{VLL}\\ {\scriptsize $\mathrm{tonf}$}\end{tabular} & \begin{tabular}[b]{@{}c@{}}\textbf{MLL}\\ {\scriptsize $\mathrm{tonf\cdot m}$}\end{tabular} & \begin{tabular}[b]{@{}c@{}}\textbf{Mf (fatiga)}\\ {\scriptsize $\mathrm{tonf\cdot m}$}\end{tabular} \\
\midrule
$0$ & $0$ & $35.469$ & $0$ & $15.432$ & $0$ & $0$ \\
$0.0600$ & $2.196$ & $33.162$ & $60.280$ & $13.636$ & $26.368$ & $27.056$ \\
$0.1000$ & $3.660$ & $31.623$ & $95.805$ & $12.500$ & $42.077$ & $42.805$ \\
$0.2000$ & $7.320$ & $27.778$ & $168.31$ & $9.877$ & $74.804$ & $74.168$ \\
$0.3000$ & $10.980$ & $23.932$ & $217.51$ & $7.562$ & $98.180$ & $95.890$ \\
$0.4000$ & $14.640$ & $20.086$ & $246.42$ & $5.556$ & $112.21$ & $107.44$ \\
$0.5000$ & $18.300$ & $16.240$ & $253.55$ & $3.858$ & $116.88$ & $107.55$ \\
\bottomrule
\end{tabular}}
\end{center}

\subsection{5. Estimación del pretensado}

\subsubsection{5.1 Esfuerzos por carga de servicio en el centro del tramo (Servicio III)}

\eqitem{Momento Mg en el centro del tramo}{}{M_{\mathrm{g}} = M_{\mathrm{g,6}} = \mathbf{198.84}\,\mathrm{tonf\cdot m}}

\eqitem{Momento Ms en el centro del tramo}{}{M_{\mathrm{s}} = M_{\mathrm{s,6}} = \mathbf{225.31}\,\mathrm{tonf\cdot m}}

\eqitem{Momento Mb en el centro del tramo}{}{M_{\mathrm{b}} = M_{\mathrm{b,6}} = \mathbf{24.921}\,\mathrm{tonf\cdot m}}

\eqitem{Momento Mws en el centro del tramo}{}{M_{\mathrm{ws}} = M_{\mathrm{ws,6}} = \mathbf{49.008}\,\mathrm{tonf\cdot m}}

\eqitem{Momento MLT en el centro del tramo}{}{\mathit{MLT} = \mathit{MLT}_{\mathrm{6}} = \mathbf{253.55}\,\mathrm{tonf\cdot m}}

\eqitem{Momento MLL en el centro del tramo}{}{\mathit{MLL} = \mathit{MLL}_{\mathrm{6}} = \mathbf{116.88}\,\mathrm{tonf\cdot m}}

\eqitem{Esfuerzo de tensión en la fibra inferior}{}{f_{\mathrm{b}} = \dfrac{M_{\mathrm{g}} + M_{\mathrm{s}}}{S_{\mathrm{b}}} + \dfrac{M_{\mathrm{b}} + M_{\mathrm{ws}} + 0.8 \cdot \left(\mathit{MLT} + \mathit{MLL}\right)}{S_{\mathrm{bc}}} = \mathbf{0.2862}\,\mathrm{tonf/cm^{2}}}

\subsubsection{5.2 Esfuerzos límites del hormigón}

\eqitem{Tensión límite en servicio}{LRFD A. 5.9.4.2b}{f_{\mathrm{ltcs}} = -0.19 \cdot \sqrt{\dfrac{f'_{c,d2}}{\mathrm{ksi}}} \cdot \mathrm{ksi} = \mathbf{-0.0338}\,\mathrm{tonf/cm^{2}}}

\subsubsection{5.3 Número de torones requeridos}

\eqitem{Precompresión requerida en la fibra inferior}{}{f_{\mathrm{pb}} = f_{\mathrm{b}} + f_{\mathrm{ltcs}} = \mathbf{0.2524}\,\mathrm{tonf/cm^{2}}}

\eqitem{Distancia estimada 0.05h}{}{y_{\mathrm{bs,e}} = 0.05 \cdot h_{\mathrm{g}} = \mathbf{9.144}\,\mathrm{cm}}

\begin{longtable}{@{}l p{0.52\linewidth} r l l@{}}
\toprule
\textbf{Dato} & \textbf{Descripción} & \textbf{Valor} & \textbf{Unidad} & \textbf{Ref.}\\
\midrule
$y_{\mathrm{bs,a}}$ & Distancia asumida CG torones - fibra inferior & \cellcolor{fondodato}10 & $\mathrm{cm}$ & {\scriptsize } \\
\bottomrule
\end{longtable}

\eqitem{Excentricidad en el centro del claro}{}{e_{\mathrm{c1}} = y_{\mathrm{b}} - y_{\mathrm{bs,a}} = \mathbf{82.964}\,\mathrm{cm}}

\eqitem{Fuerza de pretensado total después de pérdidas}{}{P_{\mathrm{pe}} = \dfrac{f_{\mathrm{pb}} \cdot A_{\mathrm{g}} \cdot S_{\mathrm{b}}}{S_{\mathrm{b}} + A_{\mathrm{g}} \cdot e_{\mathrm{c1}}} = \mathbf{466.09}\,\mathrm{tonf}}

\begin{longtable}{@{}l p{0.52\linewidth} r l l@{}}
\toprule
\textbf{Dato} & \textbf{Descripción} & \textbf{Valor} & \textbf{Unidad} & \textbf{Ref.}\\
\midrule
$p$ & Pérdida final asumida (\% de fpi) & \cellcolor{fondodato}25 & $\%$ & {\scriptsize } \\
\bottomrule
\end{longtable}

\eqitem{Fuerza final por torón}{}{F_{pe} = A_{\mathrm{st}} \cdot f_{\mathrm{pi}} \cdot \left(1 - p\right) = \mathbf{10.549}\,\mathrm{tonf}}

\eqitem{Número de torones requeridos}{}{N_{\mathrm{ti}} = \dfrac{P_{\mathrm{pe}}}{F_{pe}} = \mathbf{44.185}}

\begin{longtable}{@{}l p{0.52\linewidth} r l l@{}}
\toprule
\textbf{Dato} & \textbf{Descripción} & \textbf{Valor} & \textbf{Unidad} & \textbf{Ref.}\\
\midrule
$N_{\mathrm{t}}$ & Número de torones adoptado (2a iteración) & \cellcolor{fondodato}48 & $$ & {\scriptsize } \\
\bottomrule
\end{longtable}

\eqitem{Torones adoptados $\ge$ requeridos}{}{N_{\mathrm{t}} \ge N_{\mathrm{ti}}\;\;\Rightarrow\;\;\text{\ok{OK}}}

\eqitem{Área total de acero de preesfuerzo}{}{A_{\mathrm{ps}} = A_{\mathrm{st}} \cdot N_{\mathrm{t}} = \mathbf{47.376}\,\mathrm{cm^{2}}}

\begin{center}
{\small\textbf{Distribución de torones (centro del tramo)}}\\[3pt]
\adjustbox{max width=\linewidth}{\begin{tabular}{cccc}
\toprule
\begin{tabular}[b]{@{}c@{}}\textbf{Fila}\end{tabular} & \begin{tabular}[b]{@{}c@{}}\textbf{\# Torones (A)}\end{tabular} & \begin{tabular}[b]{@{}c@{}}\textbf{Dist. al fondo (B)}\\ {\scriptsize $\mathrm{m}$}\end{tabular} & \begin{tabular}[b]{@{}c@{}}\textbf{(A)$\cdot$(B)}\\ {\scriptsize $\mathrm{m}$}\end{tabular} \\
\midrule
y1 & \cellcolor{fondodato}$12$ & \cellcolor{fondodato}$0.0500$ & $0.6000$ \\
y2 & \cellcolor{fondodato}$12$ & \cellcolor{fondodato}$0.1000$ & $1.200$ \\
y3 & \cellcolor{fondodato}$8$ & \cellcolor{fondodato}$0.1500$ & $1.200$ \\
y4 & \cellcolor{fondodato}$4$ & \cellcolor{fondodato}$0.2000$ & $0.8000$ \\
y5 & \cellcolor{fondodato}$2$ & \cellcolor{fondodato}$0.2500$ & $0.5000$ \\
y6 & \cellcolor{fondodato}$2$ & \cellcolor{fondodato}$0.3000$ & $0.6000$ \\
y7 & \cellcolor{fondodato}$2$ & \cellcolor{fondodato}$0.3500$ & $0.7000$ \\
y8 & \cellcolor{fondodato}$2$ & \cellcolor{fondodato}$0.4000$ & $0.8000$ \\
y9 & \cellcolor{fondodato}$2$ & \cellcolor{fondodato}$0.4500$ & $0.9000$ \\
y10 & \cellcolor{fondodato}$2$ & \cellcolor{fondodato}$0.5000$ & $1$ \\
\bottomrule
\end{tabular}}
\end{center}

\eqitem{Total de torones en la tabla}{}{N_{\mathrm{t,tab}} = \sum A = \mathbf{48}}

\eqitem{La tabla de distribución contiene Nt torones}{}{N_{\mathrm{t,tab}} \ge N_{\mathrm{t}} \;\wedge\; N_{\mathrm{t,tab}} \le N_{\mathrm{t}}\;\;\Rightarrow\;\;\text{\ok{OK}}}

\eqitem{CG de torones - fibra inferior}{}{y_{\mathrm{bs}} = \dfrac{\sum A\cdot B}{\sum A} = \mathbf{17.292}\,\mathrm{cm}}

\eqitem{Excentricidad de los torones}{}{e_{\mathrm{c2}} = y_{\mathrm{b}} - y_{\mathrm{bs}} = \mathbf{75.672}\,\mathrm{cm}}

\eqitem{Excentricidad de los torones}{}{e_{\mathrm{pg}} = e_{\mathrm{c2}} = \mathbf{75.672}\,\mathrm{cm}}

\subsection{6. Pérdidas del pretensado}

\nota{$\Delta$fpT = $\Delta$fpES + $\Delta$fpSR + $\Delta$fpCR + $\Delta$fpR2  [LRFD Eq. 5.9.5.1-1]}

\subsubsection{6.1 Acortamiento elástico [LRFD Eq. 5.9.5.2.3a]}

\begin{longtable}{@{}l p{0.52\linewidth} r l l@{}}
\toprule
\textbf{Dato} & \textbf{Descripción} & \textbf{Valor} & \textbf{Unidad} & \textbf{Ref.}\\
\midrule
$P_{fpi}$ & Pérdida inicial de fpi empleada & \cellcolor{fondodato}9 & $\%$ & {\scriptsize } \\
\bottomrule
\end{longtable}

\eqitem{Fuerza por torón en la transferencia}{}{F_{\mathrm{stbt}} = A_{\mathrm{st}} \cdot f_{\mathrm{pi}} \cdot \left(1 - P_{fpi}\right) = \mathbf{12.799}\,\mathrm{tonf}}

\eqitem{Fuerza total de pretensado (antes de transferencia)}{}{P_{\mathrm{pi}} = N_{\mathrm{t}} \cdot F_{\mathrm{stbt}} = \mathbf{614.35}\,\mathrm{tonf}}

\eqitem{Esfuerzo en el CG de los torones}{}{f_{\mathrm{cgp}} = \dfrac{P_{\mathrm{pi}}}{A_{\mathrm{g}}} + \dfrac{P_{\mathrm{pi}} \cdot {e_{\mathrm{pg}}}^{2}}{I_{\mathrm{g}}} - \dfrac{M_{\mathrm{g}} \cdot e_{\mathrm{pg}}}{I_{\mathrm{g}}} = \mathbf{0.2128}\,\mathrm{tonf/cm^{2}}}

\eqitem{Pérdida por acortamiento elástico}{}{\Delta f_{pES} = \dfrac{E_{\mathrm{p}}}{E_{\mathrm{ci}}} \cdot f_{\mathrm{cgp}} = \mathbf{13\,201.26}\,\mathrm{tonf/m^{2}}}

\eqitem{Pérdida real por acortamiento elástico}{}{p_{\mathrm{act}} = \dfrac{\Delta f_{pES}}{f_{\mathrm{pi}}} = \mathbf{0.0926}}

\eqitem{Pérdida real $\approx$ asumida ($\pm$1 \%): no se requiere iterar}{}{\left|p_{\mathrm{act}} - P_{fpi}\right| \le 0.01\;\;\Rightarrow\;\;\text{\ok{OK}}}

\subsubsection{6.2 Contracción [LRFD Eq. 5.9.5.4.2-1]}

\begin{longtable}{@{}l p{0.52\linewidth} r l l@{}}
\toprule
\textbf{Dato} & \textbf{Descripción} & \textbf{Valor} & \textbf{Unidad} & \textbf{Ref.}\\
\midrule
$H$ & Humedad relativa (\%) & \cellcolor{fondodato}70 & $$ & {\scriptsize } \\
\bottomrule
\end{longtable}

\eqitem{Pérdida por contracción}{}{\Delta f_{pSR} = 17\,\mathrm{ksi} - 0.15\,\mathrm{ksi} \cdot H = \mathbf{4\,569.95}\,\mathrm{tonf/m^{2}}}

\subsubsection{6.3 Flujo plástico [LRFD Eq. 5.9.5.4.3-1]}

\eqitem{Cambio de esfuerzo por cargas permanentes}{}{\Delta f_{cdp} = \dfrac{M_{\mathrm{s}} \cdot e_{\mathrm{pg}}}{I_{\mathrm{g}}} + \dfrac{\left(M_{\mathrm{b}} + M_{\mathrm{ws}}\right) \cdot \left(y_{\mathrm{bc}} - y_{\mathrm{bs}}\right)}{I_{\mathrm{c}}} = \mathbf{947.34}\,\mathrm{tonf/m^{2}}}

\eqitem{Pérdida por flujo plástico}{}{\Delta f_{pCR} = 12 \cdot f_{\mathrm{cgp}} - 7 \cdot \Delta f_{cdp} = \mathbf{18\,900.33}\,\mathrm{tonf/m^{2}}}

\subsubsection{6.4 Relajación de los torones}

\begin{longtable}{@{}l p{0.52\linewidth} r l l@{}}
\toprule
\textbf{Dato} & \textbf{Descripción} & \textbf{Valor} & \textbf{Unidad} & \textbf{Ref.}\\
\midrule
$\Delta f_{pR1}$ & Relajación antes de la transferencia (se considera en fábrica) & \cellcolor{fondodato}0 & $\mathrm{ksi}$ & {\scriptsize } \\
\bottomrule
\end{longtable}

\eqitem{Relajación después de la transferencia}{LRFD Eq. 5.9.5.4.4c-1}{\Delta f_{pR2} = 0.3 \cdot \left(20\,\mathrm{ksi} - 0.4 \cdot \Delta f_{pES} - 0.2 \cdot \left(\Delta f_{pSR} + \Delta f_{pCR}\right)\right) = \mathbf{1\,226.05}\,\mathrm{tonf/m^{2}}}

\subsubsection{6.5 Pérdidas totales en la transferencia}

\eqitem{Pérdida inicial}{}{\Delta f_{pi} = \Delta f_{pES} = \mathbf{1.320}\,\mathrm{tonf/cm^{2}}}

\eqitem{Esfuerzo en torones después de la transferencia}{}{f_{\mathrm{pt}} = f_{\mathrm{pi}} - \Delta f_{pi} = \mathbf{12.930}\,\mathrm{tonf/cm^{2}}}

\eqitem{Fuerza en cada torón}{}{F_{\mathrm{t}} = f_{\mathrm{pt}} \cdot A_{\mathrm{st}} = \mathbf{12.762}\,\mathrm{tonf}}

\eqitem{\% de pérdida inicial}{}{p_i = \dfrac{\Delta f_{pi}}{f_{\mathrm{pi}}} \cdot 100 = \mathbf{9.264}}

\subsubsection{6.6 Pérdidas totales en cargas de servicio}

\eqitem{Pérdidas totales}{}{\Delta f_{pT} = \Delta f_{pES} + \Delta f_{pSR} + \Delta f_{pCR} + \Delta f_{pR2} = \mathbf{37\,897.59}\,\mathrm{tonf/m^{2}}}

\eqitem{Esfuerzo después de todas las pérdidas}{}{f_{\mathrm{pe}} = f_{\mathrm{pi}} - \Delta f_{pT} = \mathbf{10.460}\,\mathrm{tonf/cm^{2}}}

\eqitem{fpe $\le$ 0.8 fpy}{LRFD T. 5.9.3-1}{f_{\mathrm{pe}} \le 0.8 \cdot f_{\mathrm{py}}\;\;\Rightarrow\;\;\text{\ok{OK}}}

\eqitem{Fuerza por torón}{}{F_{tLL} = f_{\mathrm{pe}} \cdot A_{\mathrm{st}} = \mathbf{10.324}\,\mathrm{tonf}}

\eqitem{Fuerza total de pretensado después de pérdidas}{}{F_T = F_{tLL} \cdot N_{\mathrm{t}} = \mathbf{495.56}\,\mathrm{tonf}}

\eqitem{\% de pérdida final}{}{p_{\mathrm{t}} = \dfrac{f_{\mathrm{pi}} - f_{\mathrm{pe}}}{f_{\mathrm{pi}}} \cdot 100 = \mathbf{26.595}}

\subsection{7. Esfuerzos en el hormigón en la transferencia}

\eqitem{Fuerza total de pretensado después de la transferencia}{}{P_{pi,transf} = F_{\mathrm{t}} \cdot N_{\mathrm{t}} = \mathbf{612.57}\,\mathrm{tonf}}

\subsubsection{7.1 Esfuerzos límites [LRFD T. 5.9.4]}

\eqitem{Compresión en la transferencia}{}{f_{\mathrm{lc}} = 0.6 \cdot f'_{ci} = \mathbf{0.2430}\,\mathrm{tonf/cm^{2}}}

\eqitem{Tensión sin refuerzo auxiliar adherido}{}{f_{\mathrm{lt1}} = -0.0948 \cdot \sqrt{\dfrac{f'_{ci}}{\mathrm{ksi}}} \cdot \mathrm{ksi} = \mathbf{-0.2275}\,\mathrm{ksi}}

\eqitem{Tensión límite (gobierna el menor valor absoluto)}{}{f_{\mathrm{lt}} = \max\left(f_{\mathrm{lt1}},\ -0.2 \cdot \mathrm{ksi}\right) = \mathbf{-0.2000}\,\mathrm{ksi}}

\eqitem{Tensión con refuerzo auxiliar adherido}{}{f_{\mathrm{lt2}} = -0.22 \cdot \sqrt{\dfrac{f'_{ci}}{\mathrm{ksi}}} \cdot \mathrm{ksi} = \mathbf{-0.5280}\,\mathrm{ksi}}

\subsubsection{7.2 Esfuerzos en la sección de longitud de transferencia}

\eqitem{Longitud de transferencia (60 dt)}{LRFD Art. 5.8.2.3}{L_{\mathrm{transf}} = 60 \cdot d_t = \mathbf{0.7620}\,\mathrm{m}}

\eqitem{Sección de análisis}{}{x_{\mathrm{t}} = L_{\mathrm{transf}} = \mathbf{0.7620}\,\mathrm{m}}

\eqitem{Momento por peso de la viga}{}{M_{\mathrm{glt}} = 0.5 \cdot W_{\mathrm{g}} \cdot x_{\mathrm{t}} \cdot \left(L_{\mathrm{2}} - x_{\mathrm{t}}\right) = \mathbf{16.350}\,\mathrm{tonf\cdot m}}

\eqitem{Esfuerzo en la parte superior}{}{f_{\mathrm{t,lt}} = \dfrac{P_{pi,transf}}{A_{\mathrm{g}}} - \dfrac{P_{pi,transf} \cdot e_{\mathrm{pg}}}{S_{\mathrm{t}}} + \dfrac{M_{\mathrm{glt}}}{S_{\mathrm{t}}} = \mathbf{-0.0532}\,\mathrm{tonf/cm^{2}}}

\eqitem{Fibra superior (torones rectos)}{}{f_{\mathrm{t,lt}} \ge f_{\mathrm{lt2}}\;\;\Rightarrow\;\;\text{\nook{NO OK}}}

\eqitem{Esfuerzo en la parte inferior}{}{f_{\mathrm{b,lt}} = \dfrac{P_{pi,transf}}{A_{\mathrm{g}}} + \dfrac{P_{pi,transf} \cdot e_{\mathrm{pg}}}{S_{\mathrm{b}}} - \dfrac{M_{\mathrm{glt}}}{S_{\mathrm{b}}} = \mathbf{0.3068}\,\mathrm{tonf/cm^{2}}}

\eqitem{Fibra inferior (torones rectos)}{}{f_{\mathrm{b,lt}} \le f_{\mathrm{lc}}\;\;\Rightarrow\;\;\text{\nook{NO OK}}}

\nota{Los esfuerzos superan los límites: se desvían 12 torones en los puntos 0.4L.}

\begin{center}
{\small\textbf{Torones rectos en el extremo (desde el fondo)}}\\[3pt]
\adjustbox{max width=\linewidth}{\begin{tabular}{cccc}
\toprule
\begin{tabular}[b]{@{}c@{}}\textbf{Fila}\end{tabular} & \begin{tabular}[b]{@{}c@{}}\textbf{\# Torones (A)}\end{tabular} & \begin{tabular}[b]{@{}c@{}}\textbf{Dist. al fondo (B)}\\ {\scriptsize $\mathrm{m}$}\end{tabular} & \begin{tabular}[b]{@{}c@{}}\textbf{(A)$\cdot$(B)}\\ {\scriptsize $\mathrm{m}$}\end{tabular} \\
\midrule
y1 & \cellcolor{fondodato}$12$ & \cellcolor{fondodato}$0.0500$ & $0.6000$ \\
y2 & \cellcolor{fondodato}$12$ & \cellcolor{fondodato}$0.1000$ & $1.200$ \\
y3 & \cellcolor{fondodato}$8$ & \cellcolor{fondodato}$0.1500$ & $1.200$ \\
y4 & \cellcolor{fondodato}$4$ & \cellcolor{fondodato}$0.2000$ & $0.8000$ \\
\bottomrule
\end{tabular}}
\end{center}

\eqitem{Torones inferiores rectos}{}{N_{blt} = \sum A = \mathbf{36}}

\eqitem{CG de torones rectos - fibra inferior}{}{y_{\mathrm{be}} = \dfrac{\mathit{AD}_{\mathrm{e,1}} + \mathit{AD}_{\mathrm{e,2}} + \mathit{AD}_{\mathrm{e,3}} + \mathit{AD}_{\mathrm{e,4}}}{N_{blt}} = \mathbf{10.556}\,\mathrm{cm}}

\begin{center}
{\small\textbf{Torones desviados en el extremo (desde arriba)}}\\[3pt]
\adjustbox{max width=\linewidth}{\begin{tabular}{cccc}
\toprule
\begin{tabular}[b]{@{}c@{}}\textbf{Fila}\end{tabular} & \begin{tabular}[b]{@{}c@{}}\textbf{\# Torones (A)}\end{tabular} & \begin{tabular}[b]{@{}c@{}}\textbf{Dist. desde arriba (B)}\\ {\scriptsize $\mathrm{m}$}\end{tabular} & \begin{tabular}[b]{@{}c@{}}\textbf{(A)$\cdot$(B)}\\ {\scriptsize $\mathrm{m}$}\end{tabular} \\
\midrule
y5 & \cellcolor{fondodato}$2$ & \cellcolor{fondodato}$0.3000$ & $0.6000$ \\
y6 & \cellcolor{fondodato}$2$ & \cellcolor{fondodato}$0.2500$ & $0.5000$ \\
y7 & \cellcolor{fondodato}$2$ & \cellcolor{fondodato}$0.2000$ & $0.4000$ \\
y8 & \cellcolor{fondodato}$2$ & \cellcolor{fondodato}$0.1500$ & $0.3000$ \\
y9 & \cellcolor{fondodato}$2$ & \cellcolor{fondodato}$0.1000$ & $0.2000$ \\
y10 & \cellcolor{fondodato}$2$ & \cellcolor{fondodato}$0.0500$ & $0.1000$ \\
\bottomrule
\end{tabular}}
\end{center}

\eqitem{Torones desviados}{}{N_{tlt} = \sum A = \mathbf{12}}

\eqitem{CG de torones desviados - fibra superior (extremo)}{}{y_{\mathrm{te}} = \dfrac{\mathit{AD}_{\mathrm{t,5}} + \mathit{AD}_{\mathrm{t,6}} + \mathit{AD}_{\mathrm{t,7}} + \mathit{AD}_{\mathrm{t,8}} + \mathit{AD}_{\mathrm{t,9}} + \mathit{AD}_{\mathrm{t,10}}}{N_{tlt}} = \mathbf{17.500}\,\mathrm{cm}}

\eqitem{Rectos + desviados = Nt}{}{N_{blt} + N_{tlt} \ge N_{\mathrm{t}} \;\wedge\; N_{blt} + N_{tlt} \le N_{\mathrm{t}}\;\;\Rightarrow\;\;\text{\ok{OK}}}

\begin{center}
{\small\textbf{Torones desviados en el punto 'harp' (desde el fondo)}}\\[3pt]
\adjustbox{max width=\linewidth}{\begin{tabular}{cccc}
\toprule
\begin{tabular}[b]{@{}c@{}}\textbf{Fila}\end{tabular} & \begin{tabular}[b]{@{}c@{}}\textbf{\# Torones (A)}\end{tabular} & \begin{tabular}[b]{@{}c@{}}\textbf{Dist. al fondo (B)}\\ {\scriptsize $\mathrm{m}$}\end{tabular} & \begin{tabular}[b]{@{}c@{}}\textbf{(A)$\cdot$(B)}\\ {\scriptsize $\mathrm{m}$}\end{tabular} \\
\midrule
y5 & \cellcolor{fondodato}$2$ & \cellcolor{fondodato}$0.2500$ & $0.5000$ \\
y6 & \cellcolor{fondodato}$2$ & \cellcolor{fondodato}$0.3000$ & $0.6000$ \\
y7 & \cellcolor{fondodato}$2$ & \cellcolor{fondodato}$0.3500$ & $0.7000$ \\
y8 & \cellcolor{fondodato}$2$ & \cellcolor{fondodato}$0.4000$ & $0.8000$ \\
y9 & \cellcolor{fondodato}$2$ & \cellcolor{fondodato}$0.4500$ & $0.9000$ \\
y10 & \cellcolor{fondodato}$2$ & \cellcolor{fondodato}$0.5000$ & $1$ \\
\bottomrule
\end{tabular}}
\end{center}

\eqitem{Torones desviados en el punto harp}{}{N_{thp} = \sum A = \mathbf{12}}

\eqitem{CG de torones desviados - fibra inferior (punto harp)}{}{y_{\mathrm{ph}} = \dfrac{\mathit{AD}_{\mathrm{h,5}} + \mathit{AD}_{\mathrm{h,6}} + \mathit{AD}_{\mathrm{h,7}} + \mathit{AD}_{\mathrm{h,8}} + \mathit{AD}_{\mathrm{h,9}} + \mathit{AD}_{\mathrm{h,10}}}{N_{thp}} = \mathbf{37.500}\,\mathrm{cm}}

\eqitem{Punto harp (0.4 L2)}{}{l_{\mathrm{lt}} = 0.4 \cdot L_{\mathrm{2}} = \mathbf{14.760}\,\mathrm{m}}

\eqitem{CG torones desviados - fibra superior (long. transf.)}{}{y_{\mathrm{lt}} = y_{\mathrm{te}} + \dfrac{h_{\mathrm{g}} - y_{\mathrm{ph}} - y_{\mathrm{te}}}{l_{\mathrm{lt}}} \cdot L_{\mathrm{transf}} = \mathbf{24.102}\,\mathrm{cm}}

\begin{longtable}{@{}l p{0.52\linewidth} r l l@{}}
\toprule
\textbf{Dato} & \textbf{Descripción} & \textbf{Valor} & \textbf{Unidad} & \textbf{Ref.}\\
\midrule
$d_{\mathrm{str}}$ & Distancia entre torones & \cellcolor{fondodato}0.0500 & $\mathrm{m}$ & {\scriptsize } \\
\bottomrule
\end{longtable}

\eqitem{CG de todos los torones (long. transf.)}{}{y_{\mathrm{lt}} = \dfrac{N_{blt} \cdot y_{\mathrm{be}} + N_{tlt} \cdot \left(h_{\mathrm{g}} - y_{\mathrm{lt}}\right)}{N_{\mathrm{t}}} = \mathbf{47.611}\,\mathrm{cm}}

\begin{center}
{\small\textbf{Patrón de torones en el extremo (desde el fondo)}}\\[3pt]
\adjustbox{max width=\linewidth}{\begin{tabular}{cccc}
\toprule
\begin{tabular}[b]{@{}c@{}}\textbf{Fila}\end{tabular} & \begin{tabular}[b]{@{}c@{}}\textbf{\# Torones (A)}\end{tabular} & \begin{tabular}[b]{@{}c@{}}\textbf{Dist. al fondo (B)}\\ {\scriptsize $\mathrm{m}$}\end{tabular} & \begin{tabular}[b]{@{}c@{}}\textbf{(A)$\cdot$(B)}\\ {\scriptsize $\mathrm{m}$}\end{tabular} \\
\midrule
y1 & $12$ & $0.0500$ & $0.6000$ \\
y2 & $12$ & $0.1000$ & $1.200$ \\
y3 & $8$ & $0.1500$ & $1.200$ \\
y4 & $4$ & $0.2000$ & $0.8000$ \\
y5 & $2$ & $1.529$ & $3.058$ \\
y6 & $2$ & $1.579$ & $3.158$ \\
y7 & $2$ & $1.629$ & $3.258$ \\
y8 & $2$ & $1.679$ & $3.358$ \\
y9 & $2$ & $1.729$ & $3.458$ \\
y10 & $2$ & $1.779$ & $3.558$ \\
\bottomrule
\end{tabular}}
\end{center}

\eqitem{CG de todos los torones - fondo (extremo)}{}{y_{\mathrm{bs,ext}} = \dfrac{\mathit{AE}_{\mathrm{x,1}} + \mathit{AE}_{\mathrm{x,2}} + \mathit{AE}_{\mathrm{x,3}} + \mathit{AE}_{\mathrm{x,4}} + \mathit{AE}_{\mathrm{x,5}} + \mathit{AE}_{\mathrm{x,6}} + \mathit{AE}_{\mathrm{x,7}} + \mathit{AE}_{\mathrm{x,8}} + \mathit{AE}_{\mathrm{x,9}} + \mathit{AE}_{\mathrm{x,10}}}{N_{\mathrm{t}}} = \mathbf{49.262}\,\mathrm{cm}}

\eqitem{Excentricidad en la longitud de transferencia}{}{e_{\mathrm{xc}} = y_{\mathrm{b}} - y_{\mathrm{lt}} = \mathbf{45.353}\,\mathrm{cm}}

\eqitem{Esfuerzo en la parte superior}{}{f_{\mathrm{t,ext}} = \dfrac{P_{pi,transf}}{A_{\mathrm{g}}} - \dfrac{P_{pi,transf} \cdot e_{\mathrm{xc}}}{S_{\mathrm{t}}} + \dfrac{M_{\mathrm{glt}}}{S_{\mathrm{t}}} = \mathbf{0.0203}\,\mathrm{tonf/cm^{2}}}

\eqitem{Fibra superior (con torones desviados)}{}{f_{\mathrm{t,ext}} \le f_{\mathrm{lc}}\;\;\Rightarrow\;\;\text{\ok{OK}}}

\eqitem{Esfuerzo en la parte inferior}{}{f_{\mathrm{b,ext}} = \dfrac{P_{pi,transf}}{A_{\mathrm{g}}} + \dfrac{P_{pi,transf} \cdot e_{\mathrm{xc}}}{S_{\mathrm{b}}} - \dfrac{M_{\mathrm{glt}}}{S_{\mathrm{b}}} = \mathbf{0.2308}\,\mathrm{tonf/cm^{2}}}

\eqitem{Fibra inferior (con torones desviados)}{}{f_{\mathrm{b,ext}} \le f_{\mathrm{lc}}\;\;\Rightarrow\;\;\text{\ok{OK}}}

\subsubsection{7.3 Esfuerzos en los puntos 'harp'}

\eqitem{Ubicación del punto harp}{}{x_{\mathrm{ph}} = l_{\mathrm{lt}} = \mathbf{14.760}\,\mathrm{m}}

\eqitem{Momento por peso de la viga}{}{M_{\mathrm{gph}} = 0.5 \cdot W_{\mathrm{g}} \cdot x_{\mathrm{ph}} \cdot \left(L_{\mathrm{2}} - x_{\mathrm{ph}}\right) = \mathbf{194.03}\,\mathrm{tonf\cdot m}}

\eqitem{Esfuerzo en la parte superior}{}{f_{\mathrm{t,ph}} = \dfrac{P_{pi,transf}}{A_{\mathrm{g}}} - \dfrac{P_{pi,transf} \cdot e_{\mathrm{pg}}}{S_{\mathrm{t}}} + \dfrac{M_{\mathrm{gph}}}{S_{\mathrm{t}}} = \mathbf{0.0171}\,\mathrm{tonf/cm^{2}}}

\eqitem{Fibra superior en punto harp}{}{f_{\mathrm{t,ph}} \le f_{\mathrm{lc}}\;\;\Rightarrow\;\;\text{\ok{OK}}}

\eqitem{Esfuerzo en la parte inferior}{}{f_{\mathrm{b,ph}} = \dfrac{P_{pi,transf}}{A_{\mathrm{g}}} + \dfrac{P_{pi,transf} \cdot e_{\mathrm{pg}}}{S_{\mathrm{b}}} - \dfrac{M_{\mathrm{gph}}}{S_{\mathrm{b}}} = \mathbf{0.2341}\,\mathrm{tonf/cm^{2}}}

\eqitem{Fibra inferior en punto harp}{}{f_{\mathrm{b,ph}} \le f_{\mathrm{lc}}\;\;\Rightarrow\;\;\text{\ok{OK}}}

\subsubsection{7.4 Esfuerzos en el centro del tramo}

\eqitem{Centro del tramo}{}{x_{\mathrm{cl}} = 0.5 \cdot L_{\mathrm{2}} = \mathbf{18.450}\,\mathrm{m}}

\eqitem{Momento por peso de la viga}{}{M_{\mathrm{gcl}} = 0.5 \cdot W_{\mathrm{g}} \cdot x_{\mathrm{cl}} \cdot \left(L_{\mathrm{2}} - x_{\mathrm{cl}}\right) = \mathbf{202.12}\,\mathrm{tonf\cdot m}}

\eqitem{Esfuerzo en la parte superior}{}{f_{\mathrm{t,cl}} = \dfrac{P_{pi,transf}}{A_{\mathrm{g}}} - \dfrac{P_{pi,transf} \cdot e_{\mathrm{pg}}}{S_{\mathrm{t}}} + \dfrac{M_{\mathrm{gcl}}}{S_{\mathrm{t}}} = \mathbf{0.0203}\,\mathrm{tonf/cm^{2}}}

\eqitem{Fibra superior en centro del tramo}{}{f_{\mathrm{t,cl}} \le f_{\mathrm{lc}}\;\;\Rightarrow\;\;\text{\ok{OK}}}

\eqitem{Esfuerzo en la parte inferior}{}{f_{\mathrm{b,cl}} = \dfrac{P_{pi,transf}}{A_{\mathrm{g}}} + \dfrac{P_{pi,transf} \cdot e_{\mathrm{pg}}}{S_{\mathrm{b}}} - \dfrac{M_{\mathrm{gcl}}}{S_{\mathrm{b}}} = \mathbf{0.2308}\,\mathrm{tonf/cm^{2}}}

\eqitem{Fibra inferior en centro del tramo}{}{f_{\mathrm{b,cl}} \le f_{\mathrm{lc}}\;\;\Rightarrow\;\;\text{\ok{OK}}}

\subsubsection{7.5 Fuerzas de sujeción}

\eqitem{Esfuerzo en el torón al tensar}{}{f_s = 0.8 \cdot f_{\mathrm{pu}} = \mathbf{15.200}\,\mathrm{tonf/cm^{2}}}

\eqitem{Fuerza por torón}{}{F_s = A_{\mathrm{st}} \cdot f_s = \mathbf{15.002}\,\mathrm{tonf}}

\eqitem{Ángulo 'harp'}{}{\psi = \arctan\left(\dfrac{h_{\mathrm{g}} - y_{\mathrm{ph}} - y_{\mathrm{te}}}{l_{\mathrm{lt}}}\right) = \mathbf{4.952}^{\circ}}

\begin{longtable}{@{}l p{0.52\linewidth} r l l@{}}
\toprule
\textbf{Dato} & \textbf{Descripción} & \textbf{Valor} & \textbf{Unidad} & \textbf{Ref.}\\
\midrule
$f_f$ & Factor de fricción & \cellcolor{fondodato}1.050 & $$ & {\scriptsize } \\
\bottomrule
\end{longtable}

\eqitem{Fuerza de sujeción por torón}{}{F_{\mathrm{suj}} = f_f \cdot F_s \cdot \sin\left(\psi\right) = \mathbf{1.360}\,\mathrm{tonf}}

\eqitem{Fuerza de sujeción total}{}{F_{\mathrm{sujtot}} = N_{tlt} \cdot F_{\mathrm{suj}} = \mathbf{16.316}\,\mathrm{tonf}}

\subsection{8. Esfuerzos en el hormigón debido a las cargas de servicio}

\eqitem{Pretensado efectivo (pérdidas refinadas)}{}{P_{pe,c.s} = F_T = \mathbf{495.56}\,\mathrm{tonf}}

\subsubsection{8.1 Esfuerzos límites [LRFD Art. 5.9.4.2]}

\eqitem{Compresión viga - cargas permanentes (Servicio I)}{}{f_{\mathrm{pcsl2}} = 0.45 \cdot f'_{c,d2} = \mathbf{0.2025}\,\mathrm{tonf/cm^{2}}}

\eqitem{Compresión losa - cargas permanentes}{}{f_{\mathrm{pcsl1}} = 0.45 \cdot f'_{c,d1} = \mathbf{0.1260}\,\mathrm{tonf/cm^{2}}}

\eqitem{Compresión viga - todas las cargas}{}{f_{\mathrm{ptcsl2}} = 0.6 \cdot f'_{c,d2} = \mathbf{0.2700}\,\mathrm{tonf/cm^{2}}}

\eqitem{Compresión losa - todas las cargas}{}{f_{\mathrm{ptcsl1}} = 0.6 \cdot f'_{c,d1} = \mathbf{0.1680}\,\mathrm{tonf/cm^{2}}}

\eqitem{Tensión en viga (Servicio III)}{}{f_{\mathrm{pttsl2}} = -0.19 \cdot \sqrt{\dfrac{f'_{c,d2}}{\mathrm{ksi}}} \cdot \mathrm{ksi} = \mathbf{-0.0338}\,\mathrm{tonf/cm^{2}}}

\subsubsection{8.2 Esfuerzos en el centro del tramo}

\eqitem{Fibra superior viga - cargas permanentes}{}{f_{\mathrm{tg1}} = \dfrac{P_{pe,c.s}}{A_{\mathrm{g}}} - \dfrac{P_{pe,c.s} \cdot e_{\mathrm{pg}}}{S_{\mathrm{t}}} + \dfrac{M_{\mathrm{g}} + M_{\mathrm{s}}}{S_{\mathrm{t}}} + \dfrac{M_{\mathrm{ws}} + M_{\mathrm{b}}}{S_{\mathrm{tg}}} = \mathbf{0.1267}\,\mathrm{tonf/cm^{2}}}

\eqitem{Fibra superior viga - permanentes}{}{f_{\mathrm{tg1}} \le f_{\mathrm{pcsl2}}\;\;\Rightarrow\;\;\text{\ok{OK}}}

\eqitem{Fibra superior viga - permanentes + transitorias}{}{f_{\mathrm{tg2}} = f_{\mathrm{tg1}} + \dfrac{\mathit{MLT} + \mathit{MLL}}{S_{\mathrm{tg}}} = \mathbf{0.1623}\,\mathrm{tonf/cm^{2}}}

\eqitem{Fibra superior viga - todas las cargas}{}{f_{\mathrm{tg2}} \le f_{\mathrm{ptcsl2}}\;\;\Rightarrow\;\;\text{\ok{OK}}}

\eqitem{Fibra superior losa - permanentes}{}{f_{\mathrm{tg1t}} = \dfrac{M_{\mathrm{ws}} + M_{\mathrm{b}}}{S_{\mathrm{tc}}} = \mathbf{8.153\times10^{-3}}\,\mathrm{tonf/cm^{2}}}

\eqitem{Losa - permanentes}{}{f_{\mathrm{tg1t}} \le f_{\mathrm{pcsl1}}\;\;\Rightarrow\;\;\text{\ok{OK}}}

\eqitem{Fibra superior losa - permanentes + transitorias}{}{f_{\mathrm{tg2t}} = f_{\mathrm{tg1t}} + \dfrac{\mathit{MLT} + \mathit{MLL}}{S_{\mathrm{tc}}} = \mathbf{0.0490}\,\mathrm{tonf/cm^{2}}}

\eqitem{Losa - todas las cargas}{}{f_{\mathrm{tg2t}} \le f_{\mathrm{ptcsl1}}\;\;\Rightarrow\;\;\text{\ok{OK}}}

\eqitem{Fibra inferior viga - Servicio III}{}{f_b = \dfrac{P_{pe,c.s}}{A_{\mathrm{g}}} + \dfrac{P_{pe,c.s} \cdot e_{\mathrm{pg}}}{S_{\mathrm{b}}} - \dfrac{M_{\mathrm{g}} + M_{\mathrm{s}}}{S_{\mathrm{b}}} - \dfrac{M_{\mathrm{ws}} + M_{\mathrm{b}} + 0.8 \cdot \left(\mathit{MLT} + \mathit{MLL}\right)}{S_{\mathrm{bc}}} = \mathbf{-0.0326}\,\mathrm{tonf/cm^{2}}}

\eqitem{Fibra inferior viga - tensión Servicio III}{}{f_b \ge f_{\mathrm{pttsl2}}\;\;\Rightarrow\;\;\text{\ok{OK}}}

\subsubsection{8.3 Esfuerzo límite de fatiga [LRFD A. 5.5.3.1]}

\eqitem{Compresión inferior por permanentes y pretensado}{}{f_{\mathrm{tg2lf}} = \dfrac{P_{pe,c.s}}{A_{\mathrm{g}}} + \dfrac{P_{pe,c.s} \cdot e_{\mathrm{pg}}}{S_{\mathrm{b}}} - \dfrac{M_{\mathrm{g}} + M_{\mathrm{s}}}{S_{\mathrm{b}}} - \dfrac{M_{\mathrm{ws}} + M_{\mathrm{b}}}{S_{\mathrm{bc}}} = \mathbf{0.0575}\,\mathrm{tonf/cm^{2}}}

\eqitem{Momento de fatiga (centro del tramo)}{}{M_{\mathrm{f}} = M_{\mathrm{f,6}} = \mathbf{107.55}\,\mathrm{tonf\cdot m}}

\eqitem{Fatiga: 2(0.75 Mf/Sbc) < f compresión}{}{2 \cdot \dfrac{0.75 \cdot M_{\mathrm{f}}}{S_{\mathrm{bc}}} < f_{\mathrm{tg2lf}}\;\;\Rightarrow\;\;\text{\ok{OK}}}

\subsection{9. Estado límite de resistencia}

\begin{longtable}{@{}l p{0.52\linewidth} r l l@{}}
\toprule
\textbf{Dato} & \textbf{Descripción} & \textbf{Valor} & \textbf{Unidad} & \textbf{Ref.}\\
\midrule
$\eta_D$ & Factor de ductilidad & \cellcolor{fondodato}1 & $$ & {\scriptsize } \\
$\eta_R$ & Factor de redundancia & \cellcolor{fondodato}1 & $$ & {\scriptsize } \\
$\eta_I$ & Factor de importancia operativa & \cellcolor{fondodato}1 & $$ & {\scriptsize } \\
\bottomrule
\end{longtable}

\eqitem{Modificador de cargas}{}{\eta_i = \eta_D \cdot \eta_R \cdot \eta_I = \mathbf{1}}

\begin{longtable}{@{}l p{0.52\linewidth} r l l@{}}
\toprule
\textbf{Dato} & \textbf{Descripción} & \textbf{Valor} & \textbf{Unidad} & \textbf{Ref.}\\
\midrule
$\gamma_{DC}$ & Factor de carga DC & \cellcolor{fondodato}1.250 & $$ & {\scriptsize } \\
$\gamma_{DW}$ & Factor de carga DW & \cellcolor{fondodato}1.500 & $$ & {\scriptsize } \\
$\gamma_{LL}$ & Factor de carga LL & \cellcolor{fondodato}1.750 & $$ & {\scriptsize } \\
\bottomrule
\end{longtable}

\eqitem{Factor de impacto}{}{\mathit{FIM} = 1 + \mathit{IM} = \mathbf{1.330}}

\eqitem{Momento de diseño (Resistencia I)}{}{M_{\mathrm{u}} = \eta_i \cdot \left(\gamma_{DC} \cdot \left(M_{\mathrm{g}} + M_{\mathrm{b}} + M_{\mathrm{s}}\right) + \gamma_{DW} \cdot M_{\mathrm{ws}} + \gamma_{LL} \cdot \left(\mathit{MLT} + \mathit{MLL}\right)\right) = \mathbf{1\,283.11}\,\mathrm{tonf\cdot m}}

\nota{fps = fpu(1 - k$\cdot$c/d\_p)  [LRFD Eq. 5.7.3.1.1-1], válido si fpe $\ge$ 0.5 fpu.}

\eqitem{Factor k (torones de baja relajación)}{LRFD Eq. 5.7.3.1.1-2}{k = 2 \cdot \left(1.04 - \dfrac{f_{\mathrm{py}}}{f_{\mathrm{pu}}}\right) = \mathbf{0.2800}}

\eqitem{Fibra comprimida - CG de los torones}{}{d_{\mathrm{p}} = h_{\mathrm{sc}} - y_{\mathrm{bs}} = \mathbf{1.856}\,\mathrm{m}}

\begin{longtable}{@{}l p{0.52\linewidth} r l l@{}}
\toprule
\textbf{Dato} & \textbf{Descripción} & \textbf{Valor} & \textbf{Unidad} & \textbf{Ref.}\\
\midrule
$A_{\mathrm{s}}$ & Refuerzo a tensión no pretensado & \cellcolor{fondodato}0 & $\mathrm{cm^{2}}$ & {\scriptsize } \\
$A'_s$ & Refuerzo a compresión & \cellcolor{fondodato}0 & $\mathrm{cm^{2}}$ & {\scriptsize } \\
$f_y$ & fy del refuerzo a tensión & \cellcolor{fondodato}4\,200 & $\mathrm{kgf/cm^{2}}$ & {\scriptsize } \\
$f'_y$ & fy del refuerzo a compresión & \cellcolor{fondodato}4\,200 & $\mathrm{kgf/cm^{2}}$ & {\scriptsize } \\
\bottomrule
\end{longtable}

\eqitem{Ancho efectivo del ala comprimida}{}{b = w_{\mathrm{fe}} = \mathbf{2.750}\,\mathrm{m}}

\eqitem{Factor $\beta$1}{LRFD Art. 5.7.2.2}{\beta_{\mathrm{1}} = \begin{cases}0.85 & \text{si } \operatorname{red}\left(\dfrac{f'_{c,d1}}{\dfrac{\mathrm{kgf}}{{\mathrm{cm}}^{2}}}\right) \le 280\\ \begin{cases}0.65 & \text{si } \operatorname{red}\left(\dfrac{f'_{c,d1}}{\dfrac{\mathrm{kgf}}{{\mathrm{cm}}^{2}}}\right) > 550\\ 0.85 - \dfrac{0.05 \cdot \left(f'_{c,d1} - \dfrac{280\,\mathrm{kgf}}{{\mathrm{cm}}^{2}}\right)}{\dfrac{70\,\mathrm{kgf}}{{\mathrm{cm}}^{2}}} & \text{en otro caso}\end{cases} & \text{en otro caso}\end{cases} = \mathbf{0.8500}}

\eqitem{Profundidad del eje neutro}{LRFD Eq. 5.7.3.1.1-4}{c = \dfrac{A_{\mathrm{ps}} \cdot f_{\mathrm{pu}} + A_{\mathrm{s}} \cdot f_y - A'_s \cdot f'_y}{0.85 \cdot f'_{c,d1} \cdot \beta_{\mathrm{1}} \cdot b + \dfrac{k \cdot A_{\mathrm{ps}} \cdot f_{\mathrm{pu}}}{d_{\mathrm{p}}}} = \mathbf{15.795}\,\mathrm{cm}}

\eqitem{Comportamiento de sección rectangular (c < ts)}{}{c < t_{\mathrm{s}}\;\;\Rightarrow\;\;\text{\ok{OK}}}

\eqitem{fpe $\ge$ 0.5 fpu}{}{f_{\mathrm{pe}} \ge 0.5 \cdot f_{\mathrm{pu}}\;\;\Rightarrow\;\;\text{\ok{OK}}}

\eqitem{Esfuerzo promedio en el torón}{}{f_{\mathrm{ps}} = f_{\mathrm{pu}} \cdot \left(1 - \dfrac{k \cdot c}{d_{\mathrm{p}}}\right) = \mathbf{18.547}\,\mathrm{tonf/cm^{2}}}

\eqitem{Profundidad del bloque de compresión}{}{a = \beta_{\mathrm{1}} \cdot c = \mathbf{13.425}\,\mathrm{cm}}

\eqitem{Resistencia nominal a flexión}{LRFD Eq. 5.7.3.2.2-1}{M_{\mathrm{n}} = A_{\mathrm{ps}} \cdot f_{\mathrm{ps}} \cdot \left(d_{\mathrm{p}} - 0.5 \cdot a\right) = \mathbf{1\,571.77}\,\mathrm{tonf\cdot m}}

\begin{longtable}{@{}l p{0.52\linewidth} r l l@{}}
\toprule
\textbf{Dato} & \textbf{Descripción} & \textbf{Valor} & \textbf{Unidad} & \textbf{Ref.}\\
\midrule
$\phi$ & Factor de resistencia (tensión controlada) & \cellcolor{fondodato}1 & $$ & {\scriptsize LRFD 5.5.4.2.1} \\
\bottomrule
\end{longtable}

\eqitem{Resistencia a flexión mayorada}{}{M_{\mathrm{r}} = \phi \cdot M_{\mathrm{n}} = \mathbf{1\,571.77}\,\mathrm{tonf\cdot m}}

\eqitem{Resistencia a flexión Mr > Mu}{}{M_{\mathrm{r}} > M_{\mathrm{u}}\;\;\Rightarrow\;\;\text{\ok{OK}}}

\subsection{10. Límites del refuerzo [LRFD Art. 5.7.3.3]}

\eqitem{Peralte efectivo (As = 0)}{}{d_{\mathrm{e}} = d_{\mathrm{p}} = \mathbf{185.59}\,\mathrm{cm}}

\eqitem{Máximo refuerzo: c/de $\le$ 0.42}{}{\dfrac{c}{d_{\mathrm{e}}} \le 0.42\;\;\Rightarrow\;\;\text{\ok{OK}}}

\eqitem{Módulo de ruptura}{LRFD Art. 5.4.2.6}{f_{\mathrm{r}} = 0.24 \cdot \sqrt{\dfrac{f'_{c,d2}}{\mathrm{ksi}}} \cdot \mathrm{ksi} = \mathbf{0.0427}\,\mathrm{tonf/cm^{2}}}

\eqitem{Compresión por pretensado efectivo}{}{f_{pb} = \dfrac{P_{pe,c.s}}{A_{\mathrm{g}}} + \dfrac{P_{pe,c.s} \cdot e_{\mathrm{pg}}}{S_{\mathrm{b}}} = \mathbf{0.2536}\,\mathrm{tonf/cm^{2}}}

\eqitem{Momento por CM (sección no compuesta)}{}{M_{\mathrm{dnc}} = M_{\mathrm{g}} + M_{\mathrm{s}} = \mathbf{424.16}\,\mathrm{tonf\cdot m}}

\eqitem{Momento de fisuración}{}{M_{\mathrm{cr}} = S_{\mathrm{bc}} \cdot \left(f_{\mathrm{r}} + f_{pb}\right) - M_{\mathrm{dnc}} \cdot \left(\dfrac{S_{\mathrm{bc}}}{S_{\mathrm{b}}} - 1\right) = \mathbf{827.53}\,\mathrm{tonf\cdot m}}

\eqitem{1.2 Mcr}{}{1.2M_{cr} = 1.2 \cdot M_{\mathrm{cr}} = \mathbf{993.04}\,\mathrm{tonf\cdot m}}

\eqitem{1.33 Mu}{}{1.33M_u = 1.33 \cdot M_{\mathrm{u}} = \mathbf{1\,706.54}\,\mathrm{tonf\cdot m}}

\eqitem{Menor valor}{}{M_{\mathrm{menor}} = \min\left(1.2M_{cr},\ 1.33M_u\right) = \mathbf{993.04}\,\mathrm{tonf\cdot m}}

\eqitem{Refuerzo mínimo}{}{M_{\mathrm{r}} > M_{\mathrm{menor}}\;\;\Rightarrow\;\;\text{\ok{OK}}}

\subsection{11. Diseño por corte}

\nota{Se requiere refuerzo transversal cuando Vu > 0.5$\phi$(Vc + Vp)  [LRFD Eq. 5.8.2.4-1].}

\begin{longtable}{@{}l p{0.52\linewidth} r l l@{}}
\toprule
\textbf{Dato} & \textbf{Descripción} & \textbf{Valor} & \textbf{Unidad} & \textbf{Ref.}\\
\midrule
$\phi_c$ & Factor de resistencia para cortante & \cellcolor{fondodato}0.9000 & $$ & {\scriptsize LRFD 5.5.4.2.1} \\
\bottomrule
\end{longtable}

\subsubsection{11.1 Sección crítica [LRFD Art. 5.8.3.2]}

\eqitem{Altura total de la sección compuesta}{}{h_{\mathrm{c}} = h_{\mathrm{sc}} = \mathbf{202.88}\,\mathrm{cm}}

\begin{longtable}{@{}l p{0.52\linewidth} r l l@{}}
\toprule
\textbf{Dato} & \textbf{Descripción} & \textbf{Valor} & \textbf{Unidad} & \textbf{Ref.}\\
\midrule
$\theta$ & Ángulo de inclinación de los esfuerzos diagonales & \cellcolor{fondodato}23 & $^{\circ}$ & {\scriptsize } \\
\bottomrule
\end{longtable}

\eqitem{CG del grupo de torones rectos}{}{y_{\mathrm{bshp}} = y_{\mathrm{be}} = \mathbf{10.556}\,\mathrm{cm}}

\eqitem{Peralte efectivo en la sección crítica}{}{d_e = h_{\mathrm{c}} - y_{\mathrm{bshp}} = \mathbf{192.32}\,\mathrm{cm}}

\eqitem{de - a/2}{}{d_{\mathrm{v0}} = d_e - 0.5 \cdot a = \mathbf{185.61}\,\mathrm{cm}}

\eqitem{dv $\ge$ 0.9 de}{}{d_{\mathrm{v0}} \ge 0.9 \cdot d_e\;\;\Rightarrow\;\;\text{\ok{OK}}}

\eqitem{dv $\ge$ 0.72 h}{}{d_{\mathrm{v0}} \ge 0.72 \cdot h_{\mathrm{c}}\;\;\Rightarrow\;\;\text{\ok{OK}}}

\eqitem{Peralte de cortante adoptado}{}{d_{\mathrm{v}} = \max\left(d_{\mathrm{v0}},\ 0.9 \cdot d_e,\ 0.72 \cdot h_{\mathrm{c}}\right) = \mathbf{185.61}\,\mathrm{cm}}

\eqitem{dv}{}{s_{\mathrm{c1}} = d_{\mathrm{v}} = \mathbf{185.61}\,\mathrm{cm}}

\eqitem{0.5 dv cot $\theta$}{}{s_{\mathrm{c2}} = \dfrac{0.5 \cdot d_{\mathrm{v}}}{\tan\left(\theta\right)} = \mathbf{218.64}\,\mathrm{cm}}

\eqitem{Sección crítica desde el eje del apoyo}{}{s_{\mathrm{cmax}} = \max\left(s_{\mathrm{c1}},\ s_{\mathrm{c2}}\right) = \mathbf{218.64}\,\mathrm{cm}}

\eqitem{Ubicación relativa de la sección crítica}{}{x_{\mathrm{sc}} = \dfrac{s_{\mathrm{cmax}}}{L_{\mathrm{1}}} = \mathbf{0.0597}}

\nota{Las fuerzas en la sección crítica se toman de la fila x/L = 0.06 de las tablas.}

\subsubsection{11.2 Contribución del hormigón al corte [LRFD Eq. 5.8.3.3-3]}

\eqitem{Momento mayorado (sección 0.06L)}{}{M_{\mathrm{u1}} = \eta_i \cdot \left(\gamma_{DC} \cdot \left(M_{\mathrm{g,1}} + M_{\mathrm{b,1}} + M_{\mathrm{s,1}}\right) + \gamma_{DW} \cdot M_{\mathrm{ws,1}} + \gamma_{LL} \cdot \left(\mathit{MLT}_{\mathrm{1}} + \mathit{MLL}_{\mathrm{1}}\right)\right) = \mathbf{294.86}\,\mathrm{tonf\cdot m}}

\eqitem{Cortante mayorado (sección 0.06L)}{}{V_{\mathrm{u1}} = \eta_i \cdot \left(\gamma_{DC} \cdot \left(V_{\mathrm{g,1}} + V_{\mathrm{b,1}} + V_{\mathrm{s,1}}\right) + \gamma_{DW} \cdot V_{\mathrm{ws,1}} + \gamma_{LL} \cdot \left(\mathit{VLT}_{\mathrm{1}} + \mathit{VLL}_{\mathrm{1}}\right)\right) = \mathbf{142.95}\,\mathrm{tonf}}

\begin{longtable}{@{}l p{0.52\linewidth} r l l@{}}
\toprule
\textbf{Dato} & \textbf{Descripción} & \textbf{Valor} & \textbf{Unidad} & \textbf{Ref.}\\
\midrule
$N_{\mathrm{u1}}$ & Fuerza axial mayorada & \cellcolor{fondodato}0 & $\mathrm{kip}$ & {\scriptsize } \\
\bottomrule
\end{longtable}

\eqitem{Parámetro fpo}{LRFD Art. 5.8.3.4.2}{f_{\mathrm{po}} = 0.75 \cdot f_{\mathrm{pu}} = \mathbf{14.250}\,\mathrm{tonf/cm^{2}}}

\eqitem{Componente vertical del pretensado}{}{V_{\mathrm{p}} = F_{tLL} \cdot N_{thp} \cdot \sin\left(\psi\right) = \mathbf{10.694}\,\mathrm{tonf}}

\eqitem{Área de torones en el lado de tensión}{}{A_{ps,b} = N_{blt} \cdot A_{\mathrm{st}} = \mathbf{35.532}\,\mathrm{cm^{2}}}

\eqitem{Deformación en el refuerzo a tensión $\varepsilon$x}{LRFD Eq. 5.8.3.4.2-2}{\varepsilon_s = \dfrac{\dfrac{\left|M_{\mathrm{u1}}\right|}{d_{\mathrm{p}}} + 0.5 \cdot N_{\mathrm{u1}} + \dfrac{0.5 \cdot \left|V_{\mathrm{u1}} - V_{\mathrm{p}}\right|}{\tan\left(\theta\right)} - A_{ps,b} \cdot f_{\mathrm{po}}}{2 \cdot \left(E_{\mathrm{s}} \cdot A_{\mathrm{s}} + E_{\mathrm{p}} \cdot A_{ps,b}\right)} = \mathbf{-1.346\times10^{-3}}}

\begin{longtable}{@{}l p{0.52\linewidth} r l l@{}}
\toprule
\textbf{Dato} & \textbf{Descripción} & \textbf{Valor} & \textbf{Unidad} & \textbf{Ref.}\\
\midrule
$A_{ct}$ & Área de hormigón en el lado a tensión & \cellcolor{fondodato}405 & $\mathrm{in^{2}}$ & {\scriptsize } \\
\bottomrule
\end{longtable}

\eqitem{$\varepsilon$x corregida (si es negativa)}{LRFD Eq. 5.8.3.4.2-3}{\varepsilon_x = \begin{cases}\dfrac{\dfrac{\left|M_{\mathrm{u1}}\right|}{d_{\mathrm{p}}} + 0.5 \cdot N_{\mathrm{u1}} + \dfrac{0.5 \cdot \left|V_{\mathrm{u1}} - V_{\mathrm{p}}\right|}{\tan\left(\theta\right)} - A_{ps,b} \cdot f_{\mathrm{po}}}{2 \cdot \left(E_{\mathrm{cs}} \cdot A_{ct} + E_{\mathrm{s}} \cdot A_{\mathrm{s}} + E_{\mathrm{p}} \cdot A_{ps,b}\right)} & \text{si } \varepsilon_s < 0\\ \varepsilon_s & \text{en otro caso}\end{cases} = \mathbf{-9.975\times10^{-5}}}

\begin{longtable}{@{}l p{0.52\linewidth} r l l@{}}
\toprule
\textbf{Dato} & \textbf{Descripción} & \textbf{Valor} & \textbf{Unidad} & \textbf{Ref.}\\
\midrule
$b_{\mathrm{v}}$ & Ancho efectivo del alma & \cellcolor{fondodato}15.240 & $\mathrm{cm}$ & {\scriptsize } \\
\bottomrule
\end{longtable}

\eqitem{Esfuerzo por corte}{LRFD Eq. 5.8.3.4.2-1}{v_{\mathrm{u}} = \dfrac{\left|V_{\mathrm{u1}} - \phi_c \cdot V_{\mathrm{p}}\right|}{\phi_c \cdot b_{\mathrm{v}} \cdot d_{\mathrm{v}}} = \mathbf{0.0524}\,\mathrm{tonf/cm^{2}}}

\eqitem{Relación vu/f'c}{}{R_{\mathrm{el1}} = \dfrac{v_{\mathrm{u}}}{f'_{c,d2}} = \mathbf{0.1164}}

\begin{longtable}{@{}l p{0.52\linewidth} r l l@{}}
\toprule
\textbf{Dato} & \textbf{Descripción} & \textbf{Valor} & \textbf{Unidad} & \textbf{Ref.}\\
\midrule
$\theta_g$ & $\theta$ de la tabla LRFD 5.8.3.4.2-1 & \cellcolor{fondodato}22.800 & $^{\circ}$ & {\scriptsize } \\
$\beta_g$ & $\beta$ de la tabla LRFD 5.8.3.4.2-1 & \cellcolor{fondodato}2.940 & $$ & {\scriptsize } \\
\bottomrule
\end{longtable}

\eqitem{Contribución del hormigón}{}{V_{\mathrm{c}} = 0.0316 \cdot \beta_g \cdot \sqrt{\dfrac{f'_{c,d2}}{\mathrm{ksi}}} \cdot \mathrm{ksi} \cdot b_{\mathrm{v}} \cdot d_{\mathrm{v}} = \mathbf{46.744}\,\mathrm{tonf}}

\subsubsection{11.3 Contribución del refuerzo al corte}

\eqitem{0.5 $\phi$ (Vc + Vp)}{}{V = 0.5 \cdot \phi_c \cdot \left(V_{\mathrm{c}} + V_{\mathrm{p}}\right) = \mathbf{25.847}\,\mathrm{tonf}}

\eqitem{Requisito de refuerzo transversal}{LRFD Eq. 5.8.2.4-1}{V_{\mathrm{u1}} > V\;\;\Rightarrow\;\;\text{\info{Req. Ref. transversal}}}

\eqitem{Resistencia requerida del refuerzo}{}{V_{\mathrm{s1}} = \dfrac{V_{\mathrm{u1}}}{\phi_c} - V_{\mathrm{c}} - V_{\mathrm{p}} = \mathbf{101.40}\,\mathrm{tonf}}

\begin{longtable}{@{}l p{0.52\linewidth} r l l@{}}
\toprule
\textbf{Dato} & \textbf{Descripción} & \textbf{Valor} & \textbf{Unidad} & \textbf{Ref.}\\
\midrule
$\alpha$ & Inclinación del refuerzo transversal (estribos verticales) & \cellcolor{fondodato}90 & $^{\circ}$ & {\scriptsize } \\
$s_{\mathrm{p}}$ & Separación de estribos & \cellcolor{fondodato}0.3000 & $\mathrm{m}$ & {\scriptsize } \\
\bottomrule
\end{longtable}

\eqitem{Fluencia del refuerzo a cortante}{}{f_{\mathrm{yh}} = f_{\mathrm{y}} = \mathbf{42\,000}\,\mathrm{tonf/m^{2}}}

\eqitem{Área de refuerzo requerida}{}{A_{\mathrm{v}} = \dfrac{s_{\mathrm{p}} \cdot V_{\mathrm{s1}}}{\dfrac{f_{\mathrm{yh}} \cdot d_{\mathrm{v}}}{\tan\left(\theta\right)}} = \mathbf{1.656}\,\mathrm{cm^{2}}}

\eqitem{vu < 0.125 f'c}{LRFD Eq. 5.8.2.7-1}{v_{\mathrm{u}} < 0.125 \cdot f'_{c,d2}\;\;\Rightarrow\;\;\text{\ok{OK}}}

\begin{longtable}{@{}l p{0.52\linewidth} r l l@{}}
\toprule
\textbf{Dato} & \textbf{Descripción} & \textbf{Valor} & \textbf{Unidad} & \textbf{Ref.}\\
\midrule
$s_{\mathrm{maxc}}$ & Espaciamiento máx. 24 in & \cellcolor{fondodato}0.6100 & $\mathrm{m}$ & {\scriptsize } \\
\bottomrule
\end{longtable}

\eqitem{0.8 dv}{}{s_{\mathrm{c}} = 0.8 \cdot d_{\mathrm{v}} = \mathbf{148.49}\,\mathrm{cm}}

\eqitem{Espaciamiento máximo}{}{S_{\mathrm{max}} = \min\left(s_{\mathrm{maxc}},\ s_{\mathrm{c}}\right) = \mathbf{0.6100}\,\mathrm{m}}

\eqitem{Separación $\le$ máxima}{}{S_{\mathrm{max}} > s_{\mathrm{p}}\;\;\Rightarrow\;\;\text{\ok{OK}}}

\begin{longtable}{@{}l p{0.52\linewidth} r l l@{}}
\toprule
\textbf{Dato} & \textbf{Descripción} & \textbf{Valor} & \textbf{Unidad} & \textbf{Ref.}\\
\midrule
$d_{\mathrm{b}}$ & Diámetro del estribo & \cellcolor{fondodato}1.200 & $\mathrm{cm}$ & {\scriptsize } \\
\bottomrule
\end{longtable}

\eqitem{Área de una barra}{}{A_{\mathrm{b}} = 0.25 \cdot \pi \cdot {d_{\mathrm{b}}}^{2} = \mathbf{1.131}\,\mathrm{cm^{2}}}

\eqitem{Área del estribo (2 ramas)}{}{A_{\mathrm{bt}} = 2 \cdot A_{\mathrm{b}} = \mathbf{2.262}\,\mathrm{cm^{2}}}

\eqitem{Estribo proporcionado $\ge$ requerido}{}{A_{\mathrm{bt}} > A_{\mathrm{v}}\;\;\Rightarrow\;\;\text{\ok{OK}}}

\eqitem{Resistencia del refuerzo proporcionado}{LRFD Eq. 5.8.3.3-4}{V_{\mathrm{s2}} = \dfrac{\dfrac{A_{\mathrm{bt}} \cdot f_{\mathrm{yh}} \cdot d_{\mathrm{v}}}{\tan\left(\theta\right)}}{s_{\mathrm{p}}} = \mathbf{138.47}\,\mathrm{tonf}}

\eqitem{Refuerzo transversal mínimo}{LRFD Eq. 5.8.2.5-1}{A_{\mathrm{min}} = \dfrac{0.0316 \cdot \sqrt{\dfrac{f'_{c,d2}}{\mathrm{ksi}}} \cdot \mathrm{ksi} \cdot b_{\mathrm{v}} \cdot s_{\mathrm{p}}}{f_{\mathrm{yh}}} = \mathbf{0.6119}\,\mathrm{cm^{2}}}

\eqitem{Refuerzo $\ge$ mínimo}{}{A_{\mathrm{bt}} > A_{\mathrm{min}}\;\;\Rightarrow\;\;\text{\ok{OK}}}

\subsubsection{11.4 Máxima resistencia nominal al corte}

\eqitem{Vc + Vs}{}{V_{\mathrm{n2}} = V_{\mathrm{c}} + V_{\mathrm{s2}} = \mathbf{185.22}\,\mathrm{tonf}}

\eqitem{0.25 f'c bv dv}{LRFD Eq. 5.8.3.3-2}{V_{\mathrm{n3}} = 0.25 \cdot f'_{c,d2} \cdot b_{\mathrm{v}} \cdot d_{\mathrm{v}} = \mathbf{318.23}\,\mathrm{tonf}}

\eqitem{Vc + Vs $\le$ 0.25 f'c bv dv}{}{V_{\mathrm{n2}} < V_{\mathrm{n3}}\;\;\Rightarrow\;\;\text{\ok{OK}}}

\subsection{12. Transferencia de corte en la interfaz [LRFD A. 5.8.4]}

\eqitem{Distancia acero a tensión - mitad de la losa}{}{d_{\mathrm{v1}} = d_e - 0.5 \cdot t_{\mathrm{s}} = \mathbf{1.828}\,\mathrm{m}}

\eqitem{Cortante mayorado por cargas sobreimpuestas (0.06L)}{}{V_{u1} = \eta_i \cdot \left(\gamma_{DC} \cdot V_{\mathrm{b,1}} + \gamma_{DW} \cdot V_{\mathrm{ws,1}} + \gamma_{LL} \cdot \left(\mathit{VLT}_{\mathrm{1}} + \mathit{VLL}_{\mathrm{1}}\right)\right) = \mathbf{91.962}\,\mathrm{tonf}}

\eqitem{Cortante horizontal mayorado}{LRFD Eq. C5.8.4.1-1}{V_{\mathrm{hi}} = \dfrac{V_{u1}}{d_{\mathrm{v1}}} = \mathbf{50.301}\,\mathrm{tonf/m}}

\eqitem{Resistencia nominal requerida}{}{V_{\mathrm{ni}} = \dfrac{V_{\mathrm{hi}}}{\phi_c} = \mathbf{55.890}\,\mathrm{tonf/m}}

\begin{longtable}{@{}l p{0.52\linewidth} r l l@{}}
\toprule
\textbf{Dato} & \textbf{Descripción} & \textbf{Valor} & \textbf{Unidad} & \textbf{Ref.}\\
\midrule
$c$ & Factor de cohesión & \cellcolor{fondodato}0.0750 & $\mathrm{ksi}$ & {\scriptsize LRFD 5.8.4.2} \\
$\lambda$ & Hormigón de peso normal & \cellcolor{fondodato}1 & $$ & {\scriptsize } \\
\bottomrule
\end{longtable}

\eqitem{Coeficiente de fricción}{}{\mu = 0.6 \cdot \lambda = \mathbf{0.6000}}

\eqitem{Ancho de contacto losa-viga}{}{b_v = w_{\mathrm{fg}} = \mathbf{106.68}\,\mathrm{cm}}

\begin{longtable}{@{}l p{0.52\linewidth} r l l@{}}
\toprule
\textbf{Dato} & \textbf{Descripción} & \textbf{Valor} & \textbf{Unidad} & \textbf{Ref.}\\
\midrule
$L_{\mathrm{vi}}$ & Longitud unitaria de interfaz & \cellcolor{fondodato}2.540 & $\mathrm{cm}$ & {\scriptsize } \\
\bottomrule
\end{longtable}

\eqitem{Área que resiste la transferencia de corte}{}{A_{\mathrm{cv}} = b_v \cdot L_{\mathrm{vi}} = \mathbf{270.97}\,\mathrm{cm^{2}}}

\begin{longtable}{@{}l p{0.52\linewidth} r l l@{}}
\toprule
\textbf{Dato} & \textbf{Descripción} & \textbf{Valor} & \textbf{Unidad} & \textbf{Ref.}\\
\midrule
$P_{\mathrm{c}}$ & Fuerza de compresión normal al plano & \cellcolor{fondodato}0 & $\mathrm{tonf}$ & {\scriptsize } \\
\bottomrule
\end{longtable}

\eqitem{Refuerzo requerido en la interfaz (como en el PDF)}{}{A_{\mathrm{vf}} = \dfrac{\dfrac{V_{\mathrm{ni}} \cdot \mathrm{cm} - c \cdot A_{\mathrm{cv}}}{\mu} - P_{\mathrm{c}}}{f_{\mathrm{yh}}} = \mathbf{-0.3452}\,\mathrm{cm^{2}}}

\nota{Avf < 0: la cohesión supera la fuerza aplicada; se proporciona el refuerzo mínimo.}

\eqitem{Refuerzo proporcionado (estribo $\varnothing$12 mm @ 30 cm)}{}{A_{\mathrm{vfmin}} = \dfrac{A_{\mathrm{bt}}}{s_{\mathrm{p}}} = \mathbf{7.540}\,\mathrm{cm^{2}/m}}

\eqitem{Mínimo refuerzo en la interfaz 0.05 bv/fy}{LRFD Eq. 5.8.4.1-4}{A_{\mathrm{vf2}} = \dfrac{0.05\,\mathrm{ksi} \cdot b_v}{f_{\mathrm{yh}}} = \mathbf{8.929}\,\mathrm{cm^{2}/m}}

\eqitem{Refuerzo proporcionado $\ge$ mínimo}{}{A_{\mathrm{vfmin}} > A_{\mathrm{vf2}}\;\;\Rightarrow\;\;\text{\info{NO OK}}}

\eqitem{Vn/Acv}{}{R_{\mathrm{el2}} = \dfrac{V_{\mathrm{ni}}}{b_v} = \mathbf{0.0745}\,\mathrm{ksi}}

\eqitem{Vn/Acv < 0.1 ksi: se puede despreciar el mínimo}{}{R_{\mathrm{el2}} < 0.1\,\mathrm{ksi}\;\;\Rightarrow\;\;\text{\ok{OK}}}

\eqitem{Resistencia nominal proporcionada}{}{V_{\mathrm{nip}} = c \cdot b_v + \mu \cdot \dfrac{A_{\mathrm{bt}}}{12\,\mathrm{in}} \cdot f_{\mathrm{yh}} = \mathbf{4.197}\,\mathrm{kip/in}}

\begin{longtable}{@{}l p{0.52\linewidth} r l l@{}}
\toprule
\textbf{Dato} & \textbf{Descripción} & \textbf{Valor} & \textbf{Unidad} & \textbf{Ref.}\\
\midrule
$K_1$ & Factor K1 & \cellcolor{fondodato}0.2000 & $$ & {\scriptsize } \\
\bottomrule
\end{longtable}

\eqitem{0.2 f'c Acv}{LRFD Eq. 5.8.4.1-2}{k_{\mathrm{1fcbv}} = K_1 \cdot f'_{c,d1} \cdot b_v = \mathbf{33.453}\,\mathrm{kip/in}}

\begin{longtable}{@{}l p{0.52\linewidth} r l l@{}}
\toprule
\textbf{Dato} & \textbf{Descripción} & \textbf{Valor} & \textbf{Unidad} & \textbf{Ref.}\\
\midrule
$K_2$ & Factor K2 & \cellcolor{fondodato}0.8000 & $\mathrm{kip/in^{2}}$ & {\scriptsize } \\
\bottomrule
\end{longtable}

\eqitem{0.8 Acv}{LRFD Eq. 5.8.4.1-3}{k_{\mathrm{2Acv}} = K_2 \cdot b_v = \mathbf{33.600}\,\mathrm{kip/in}}

\eqitem{Vni $\le$ K1 f'c Acv}{}{V_{\mathrm{nip}} < k_{\mathrm{1fcbv}}\;\;\Rightarrow\;\;\text{\ok{OK}}}

\eqitem{Vni $\le$ K2 Acv}{}{V_{\mathrm{nip}} < k_{\mathrm{2Acv}}\;\;\Rightarrow\;\;\text{\ok{OK}}}

\subsection{13. Refuerzo longitudinal mínimo [LRFD Art. 5.8.3.5]}

\nota{Aps$\cdot$fps + As$\cdot$fy $\ge$ Mu/(dv$\cdot$$\phi$f) + 0.5Nu/$\phi$c + (Vu/$\phi$v - Vp - 0.5Vs)$\cdot$cot $\theta$  [LRFD Eq. 5.8.3.5-1]}

\subsubsection{13.1 Refuerzo requerido en la cara del apoyo}

\eqitem{Eje del apoyo al extremo de la viga}{}{d_{\mathrm{ap}} = 0.5 \cdot \left(L_{\mathrm{2}} - L_{\mathrm{1}}\right) = \mathbf{15}\,\mathrm{cm}}

\eqitem{Cortante mayorado en el apoyo}{}{V_{\mathrm{uapoyo}} = \eta_i \cdot \left(\gamma_{DC} \cdot \left(V_{\mathrm{g,0}} + V_{\mathrm{b,0}} + V_{\mathrm{s,0}}\right) + \gamma_{DW} \cdot V_{\mathrm{ws,0}} + \gamma_{LL} \cdot \left(\mathit{VLT}_{\mathrm{0}} + \mathit{VLL}_{\mathrm{0}}\right)\right) = \mathbf{158.46}\,\mathrm{tonf}}

\begin{longtable}{@{}l p{0.52\linewidth} r l l@{}}
\toprule
\textbf{Dato} & \textbf{Descripción} & \textbf{Valor} & \textbf{Unidad} & \textbf{Ref.}\\
\midrule
$M_u$ & Momento mayorado en el apoyo & \cellcolor{fondodato}0 & $\mathrm{tonf\cdot m}$ & {\scriptsize } \\
$N_u$ & Fuerza axial en el apoyo & \cellcolor{fondodato}0 & $\mathrm{tonf}$ & {\scriptsize } \\
\bottomrule
\end{longtable}

\eqitem{Tensión requerida en el refuerzo longitudinal}{}{V_{\mathrm{ca}} = \dfrac{M_u}{d_{\mathrm{v}} \cdot \phi_c} + \dfrac{0.5 \cdot N_u}{\phi_c} + \dfrac{\dfrac{V_{\mathrm{uapoyo}}}{\phi_c} - 0.5 \cdot V_{\mathrm{s2}} - V_{\mathrm{p}}}{\tan\left(\theta\right)} = \mathbf{226.49}\,\mathrm{tonf}}

\eqitem{Cruce de la fisura con el CG de torones rectos}{}{d_{\mathrm{apd}} = d_{\mathrm{ap}} + \dfrac{y_{\mathrm{be}}}{\tan\left(\theta\right)} = \mathbf{39.867}\,\mathrm{cm}}

\eqitem{Longitud de transferencia}{}{L_{\mathrm{transf2}} = 60 \cdot d_t = \mathbf{76.200}\,\mathrm{cm}}

\eqitem{Pretensado disponible}{}{f_{\mathrm{psb1}} = \dfrac{\left(f_{\mathrm{pi}} - \Delta f_{pT}\right) \cdot \left(d_{\mathrm{ap}} + \dfrac{y_{\mathrm{be}}}{\tan\left(\theta\right)}\right)}{L_{\mathrm{transf2}}} = \mathbf{5.473}\,\mathrm{tonf/cm^{2}}}

\eqitem{Fuerza de pretensado disponible}{}{F_{\mathrm{pa1}} = A_{ps,b} \cdot f_{\mathrm{psb1}} + A_{\mathrm{s}} \cdot f_y = \mathbf{194.46}\,\mathrm{tonf}}

\eqitem{Fuerza disponible $\ge$ requerida}{}{F_{\mathrm{pa1}} > V_{\mathrm{ca}}\;\;\Rightarrow\;\;\text{\info{NO OK}}}

\eqitem{Acero adicional requerido}{}{A_{\mathrm{sca}} = \dfrac{V_{\mathrm{ca}} - F_{\mathrm{pa1}}}{f_{\mathrm{yh}}} = \mathbf{7.626}\,\mathrm{cm^{2}}}

\begin{longtable}{@{}l p{0.52\linewidth} r l l@{}}
\toprule
\textbf{Dato} & \textbf{Descripción} & \textbf{Valor} & \textbf{Unidad} & \textbf{Ref.}\\
\midrule
$d_{\mathrm{ca}}$ & Diámetro del refuerzo adicional & \cellcolor{fondodato}1.600 & $\mathrm{cm}$ & {\scriptsize } \\
\bottomrule
\end{longtable}

\eqitem{Área de una barra}{}{A_{\mathrm{bca}} = 0.25 \cdot \pi \cdot {d_{\mathrm{ca}}}^{2} = \mathbf{2.011}\,\mathrm{cm^{2}}}

\eqitem{Número de barras}{}{N_{\mathrm{bca}} = \left\lceil \dfrac{A_{\mathrm{sca}}}{A_{\mathrm{bca}}}\right\rceil = \mathbf{4}}

\begin{longtable}{@{}l p{0.52\linewidth} r l l@{}}
\toprule
\textbf{Dato} & \textbf{Descripción} & \textbf{Valor} & \textbf{Unidad} & \textbf{Ref.}\\
\midrule
$w_{\mathrm{ap}}$ & Ancho de apoyo asumido & \cellcolor{fondodato}6 & $\mathrm{cm}$ & {\scriptsize } \\
\bottomrule
\end{longtable}

\eqitem{Cruce de la fisura desde el extremo}{}{d_{\mathrm{ap}} = 0.5 \cdot w_{\mathrm{ap}} + d_{\mathrm{apd}} = \mathbf{42.867}\,\mathrm{cm}}

\eqitem{Pretensado disponible}{}{f_{\mathrm{psb2}} = \dfrac{\left(f_{\mathrm{pi}} - \Delta f_{pT}\right) \cdot \left(d_{\mathrm{ap}} + 0.5 \cdot w_{\mathrm{ap}} + \dfrac{y_{\mathrm{be}}}{\tan\left(\theta\right)}\right)}{L_{\mathrm{transf2}}} = \mathbf{83.698}\,\mathrm{ksi}}

\eqitem{Fuerza disponible}{}{F_{\mathrm{pa2}} = A_{ps,b} \cdot f_{\mathrm{psb2}} + A_{\mathrm{s}} \cdot f_y = \mathbf{460.96}\,\mathrm{kip}}

\eqitem{Fuerza disponible > V (criterio del PDF)}{}{F_{\mathrm{pa2}} > V\;\;\Rightarrow\;\;\text{\ok{CUMPLE}}}

\subsection{14. Zona de anclaje pretensada [Art. 5.10.10]}

\eqitem{Fuerza en torones antes de la transferencia}{}{P_{pi} = N_{\mathrm{t}} \cdot A_{\mathrm{st}} \cdot f_{\mathrm{pi}} = \mathbf{675.11}\,\mathrm{tonf}}

\eqitem{Resistencia a la fractura (4\% Ppi)}{}{P_{\mathrm{r}} = 0.04 \cdot P_{pi} = \mathbf{27.004}\,\mathrm{tonf}}

\begin{longtable}{@{}l p{0.52\linewidth} r l l@{}}
\toprule
\textbf{Dato} & \textbf{Descripción} & \textbf{Valor} & \textbf{Unidad} & \textbf{Ref.}\\
\midrule
$f_s$ & Esfuerzo admisible en el acero & \cellcolor{fondodato}20 & $\mathrm{ksi}$ & {\scriptsize } \\
\bottomrule
\end{longtable}

\eqitem{Área de acero requerida}{}{A_{\mathrm{sza}} = \dfrac{P_{\mathrm{r}}}{f_s} = \mathbf{19.205}\,\mathrm{cm^{2}}}

\eqitem{Distancia h/4 desde el extremo}{}{d_{\mathrm{025h}} = 0.25 \cdot h_{\mathrm{g}} = \mathbf{45.720}\,\mathrm{cm}}

\begin{longtable}{@{}l p{0.52\linewidth} r l l@{}}
\toprule
\textbf{Dato} & \textbf{Descripción} & \textbf{Valor} & \textbf{Unidad} & \textbf{Ref.}\\
\midrule
$N_{\mathrm{bza}}$ & Número de barras (doble ramal) & \cellcolor{fondodato}5 & $$ & {\scriptsize } \\
$d_{\mathrm{bza}}$ & Diámetro de barra & \cellcolor{fondodato}1.600 & $\mathrm{cm}$ & {\scriptsize } \\
\bottomrule
\end{longtable}

\eqitem{Área de una barra}{}{A_{\mathrm{bza}} = \dfrac{\pi}{4} \cdot {d_{\mathrm{bza}}}^{2} = \mathbf{2.011}\,\mathrm{cm^{2}}}

\eqitem{Área total}{}{A_{\mathrm{btza}} = N_{\mathrm{bza}} \cdot 2 \cdot A_{\mathrm{bza}} = \mathbf{20.106}\,\mathrm{cm^{2}}}

\eqitem{Refuerzo en zona de anclaje}{}{A_{\mathrm{btza}} > A_{\mathrm{sza}}\;\;\Rightarrow\;\;\text{\ok{OK}}}

\begin{longtable}{@{}l p{0.52\linewidth} r l l@{}}
\toprule
\textbf{Dato} & \textbf{Descripción} & \textbf{Valor} & \textbf{Unidad} & \textbf{Ref.}\\
\midrule
$s_{\mathrm{ep,za}}$ & Espaciamiento & \cellcolor{fondodato}10 & $\mathrm{cm}$ & {\scriptsize } \\
$d_{\mathrm{za}}$ & Primera barra desde el extremo & \cellcolor{fondodato}5 & $\mathrm{cm}$ & {\scriptsize } \\
\bottomrule
\end{longtable}

\subsubsection{14.2 Refuerzo de confinamiento [LRFD Art. 5.10.10.2]}

\eqitem{Longitud confinada desde el extremo}{}{d_{\mathrm{cf}} = 1.5 \cdot h_{\mathrm{g}} = \mathbf{108}\,\mathrm{in}}

\begin{longtable}{@{}l p{0.52\linewidth} r l l@{}}
\toprule
\textbf{Dato} & \textbf{Descripción} & \textbf{Valor} & \textbf{Unidad} & \textbf{Ref.}\\
\midrule
$d_{\mathrm{rc}}$ & Espaciamiento máximo del refuerzo de confinamiento & \cellcolor{fondodato}15 & $\mathrm{cm}$ & {\scriptsize } \\
\bottomrule
\end{longtable}

\subsection{15. Deflexión y contraflecha [LRFD Art. 5.7.3.6.2]}

\eqitem{Fuerza total de pretensado}{}{P_{\mathrm{pt}} = N_{\mathrm{t}} \cdot F_{\mathrm{t}} = \mathbf{612.57}\,\mathrm{tonf}}

\eqitem{Diferencia de excentricidad centro - extremo}{}{e' = e_{\mathrm{c2}} - \left(y_{\mathrm{b}} - y_{\mathrm{bs,ext}}\right) = \mathbf{0.3197}\,\mathrm{m}}

\eqitem{Extremo de la viga - punto harp}{}{a = l_{\mathrm{lt}} = \mathbf{14.760}\,\mathrm{m}}

\eqitem{Contraflecha por pretensado ($\uparrow$)}{}{\Delta_p = \dfrac{P_{\mathrm{pt}}}{E_{\mathrm{ci}} \cdot I_{\mathrm{g}}} \cdot \left(\dfrac{e_{\mathrm{c2}} \cdot {L_{\mathrm{2}}}^{2}}{8} - \dfrac{e' \cdot {a}^{2}}{6}\right) = \mathbf{9.783}\,\mathrm{cm}}

\eqitem{Peso propio después de la transferencia ($\downarrow$)}{}{\Delta_g = \dfrac{5 \cdot W_{\mathrm{g}} \cdot {L_{\mathrm{2}}}^{4}}{384 \cdot E_{\mathrm{ci}} \cdot I_{\mathrm{g}}} = \mathbf{3.907}\,\mathrm{cm}}

\eqitem{Peso propio en el izaje ($\downarrow$)}{}{\Delta_{gi} = \dfrac{5 \cdot W_{\mathrm{g}} \cdot {L_{\mathrm{1}}}^{4}}{384 \cdot E_{\mathrm{ci}} \cdot I_{\mathrm{g}}} = \mathbf{3.781}\,\mathrm{cm}}

\eqitem{Losa y cartela ($\downarrow$)}{}{\Delta_s = \dfrac{5 \cdot W_{\mathrm{s}} \cdot {L_{\mathrm{1}}}^{4}}{384 \cdot E_{\mathrm{cs}} \cdot I_{\mathrm{g}}} = \mathbf{4.065}\,\mathrm{cm}}

\eqitem{Barreras y capa de rodadura ($\downarrow$)}{}{\Delta_{b+ws} = \dfrac{5 \cdot \left(\mathit{WDC}_{\mathrm{b}} + \mathit{WDW}\right) \cdot {L_{\mathrm{1}}}^{4}}{384 \cdot E_{\mathrm{cs}} \cdot I_{\mathrm{c}}} = \mathbf{0.6633}\,\mathrm{cm}}

\eqitem{Después de la transferencia ($\uparrow$)}{}{\Delta_{pg} = \Delta_p - \Delta_g = \mathbf{5.876}\,\mathrm{cm}}

\eqitem{Total en el izaje, factores PCI ($\uparrow$)}{}{\Delta_{pg,izaje} = 1.8 \cdot \Delta_p - 1.85 \cdot \Delta_{gi} = \mathbf{10.614}\,\mathrm{cm}}

\subsubsection{Deflexión por carga viva e impacto [LRFD Art. 2.5.2.6.2]}

\eqitem{Límite L/800}{}{\Delta_{LL,adm} = \dfrac{L_{\mathrm{1}}}{800} = \mathbf{4.575}\,\mathrm{cm}}

\eqitem{Factor de distribución para deflexión}{}{\mathit{FDF} = \dfrac{\mathit{NL}}{N_{\mathrm{b}}} = \mathbf{0.6667}}

\eqitem{Carga de carril por viga (con DFM)}{}{w = \mathit{WLL} \cdot \mathit{FDM}_{\mathrm{imax}} = \mathbf{0.6980}\,\mathrm{tonf/m}}

\eqitem{Deflexión por carga de carril ($\downarrow$)}{}{\Delta_{LL} = \dfrac{5 \cdot w \cdot {L_{\mathrm{1}}}^{4}}{384 \cdot E_{\mathrm{cs}} \cdot I_{\mathrm{c}}} = \mathbf{1.049}\,\mathrm{cm}}

\begin{longtable}{@{}l p{0.52\linewidth} r l l@{}}
\toprule
\textbf{Dato} & \textbf{Descripción} & \textbf{Valor} & \textbf{Unidad} & \textbf{Ref.}\\
\midrule
$P_{\mathrm{1}}$ & Carga por eje trasero & \cellcolor{fondodato}32 & $\mathrm{kip}$ & {\scriptsize } \\
$L_{\mathrm{2cd}}$ & Separación de ejes & \cellcolor{fondodato}168 & $\mathrm{in}$ & {\scriptsize } \\
$b_{\mathrm{d}}$ & Posición del eje delantero & \cellcolor{fondodato}11.895 & $\mathrm{m}$ & {\scriptsize } \\
\bottomrule
\end{longtable}

\eqitem{Centro del claro}{}{x_{\mathrm{d}} = 0.5 \cdot L_{\mathrm{1}} = \mathbf{18.300}\,\mathrm{m}}

\eqitem{Posición de los ejes traseros}{}{a_{CD1} = 0.5 \cdot \left(L_{\mathrm{1}} - L_{\mathrm{2cd}}\right) = \mathbf{53.039}\,\mathrm{ft}}

\eqitem{Deflexión por ejes traseros}{}{\Delta_{P1} = \dfrac{P_{\mathrm{1}} \cdot a_{CD1}}{24 \cdot E_{\mathrm{cs}} \cdot I_{\mathrm{c}}} \cdot \left(3 \cdot {L_{\mathrm{1}}}^{2} - 4 \cdot {a_{CD1}}^{2}\right) = \mathbf{1.869}\,\mathrm{cm}}

\eqitem{Deflexión por eje delantero}{}{\Delta_{P2} = \dfrac{0.25 \cdot P_{\mathrm{1}} \cdot b_{\mathrm{d}} \cdot x_{\mathrm{d}}}{6 \cdot E_{\mathrm{cs}} \cdot I_{\mathrm{c}} \cdot L_{\mathrm{1}}} \cdot \left({L_{\mathrm{1}}}^{2} - {b_{\mathrm{d}}}^{2} - {x_{\mathrm{d}}}^{2}\right) = \mathbf{0.1996}\,\mathrm{cm}}

\eqitem{Deflexión del camión}{}{\Delta_{PT} = \Delta_{P1} + \Delta_{P2} = \mathbf{2.069}\,\mathrm{cm}}

\eqitem{Camión con impacto ($\downarrow$)}{}{\Delta_{LT} = \Delta_{PT} \cdot \mathit{FIM} \cdot \mathit{FDM}_{\mathrm{imax}} = \mathbf{2.017}\,\mathrm{cm}}

\eqitem{25\% camión + carril}{}{\Delta_{LT,c2} = 0.25 \cdot \Delta_{LT} + \Delta_{LL} = \mathbf{1.553}\,\mathrm{cm}}

\eqitem{Deflexión por carga viva}{}{\Delta_{mayor} = \max\left(\Delta_{LT},\ \Delta_{LT,c2}\right) = \mathbf{2.017}\,\mathrm{cm}}

\eqitem{Deflexión por carga viva $\le$ L/800}{}{\Delta_{mayor} < \Delta_{LL,adm}\;\;\Rightarrow\;\;\text{\ok{OK}}}

\clearpage
\section{Descarga CM sobre una pila - vigas de acero}

\nota{Descarga por carga muerta de la superestructura (vigas de acero) sobre una pila.}

\subsection{Datos}

\begin{longtable}{@{}l p{0.52\linewidth} r l l@{}}
\toprule
\textbf{Dato} & \textbf{Descripción} & \textbf{Valor} & \textbf{Unidad} & \textbf{Ref.}\\
\midrule
$L_{\mathrm{p}}$ & Longitud del puente & \cellcolor{fondodato}120 & $\mathrm{m}$ & {\scriptsize } \\
$L_{\mathrm{t}}$ & Longitud del tramo & \cellcolor{fondodato}30 & $\mathrm{m}$ & {\scriptsize } \\
$A_{\mathrm{tp}}$ & Ancho del tablero & \cellcolor{fondodato}10 & $\mathrm{m}$ & {\scriptsize } \\
\bottomrule
\end{longtable}

\subsection{Altura y separación de vigas}

\begin{center}
{\small\textbf{Distribución transversal}}\\[3pt]
\adjustbox{max width=\linewidth}{\begin{tabular}{cccc}
\toprule
\begin{tabular}[b]{@{}c@{}}\textbf{Descripción}\end{tabular} & \begin{tabular}[b]{@{}c@{}}\textbf{Unidad}\\ {\scriptsize $\mathrm{m}$}\end{tabular} & \begin{tabular}[b]{@{}c@{}}\textbf{Cantidad}\end{tabular} & \begin{tabular}[b]{@{}c@{}}\textbf{Total}\\ {\scriptsize $\mathrm{m}$}\end{tabular} \\
\midrule
Separación vigas & \cellcolor{fondodato}$2.500$ & \cellcolor{fondodato}$3$ & $7.500$ \\
Longitud voladizo & \cellcolor{fondodato}$1.250$ & \cellcolor{fondodato}$2$ & $2.500$ \\
Ancho total puente &  &  & $10$ \\
\bottomrule
\end{tabular}}
\end{center}

\eqitem{Distribución transversal = ancho del tablero}{}{a_{\mathrm{n,t}} \ge A_{\mathrm{tp}} - 0.001 \;\wedge\; a_{\mathrm{n,t}} \le A_{\mathrm{tp}} + 0.001\;\;\Rightarrow\;\;\text{\ok{OK}}}

\subsection{Peso de la losa}

\begin{longtable}{@{}l p{0.52\linewidth} r l l@{}}
\toprule
\textbf{Dato} & \textbf{Descripción} & \textbf{Valor} & \textbf{Unidad} & \textbf{Ref.}\\
\midrule
$e_l$ & Espesor mínimo de losa & \cellcolor{fondodato}7 & $\mathrm{in}$ & {\scriptsize } \\
$e_a$ & Espesor de losa asumido & \cellcolor{fondodato}0.2000 & $\mathrm{m}$ & {\scriptsize } \\
\bottomrule
\end{longtable}

\eqitem{Espesor asumido $\ge$ mínimo}{}{e_a \ge e_l\;\;\Rightarrow\;\;\text{\ok{OK}}}

\begin{longtable}{@{}l p{0.52\linewidth} r l l@{}}
\toprule
\textbf{Dato} & \textbf{Descripción} & \textbf{Valor} & \textbf{Unidad} & \textbf{Ref.}\\
\midrule
$\gamma_{ha}$ & Peso unitario del hormigón armado & \cellcolor{fondodato}2.400 & $\mathrm{tonf/m^{3}}$ & {\scriptsize } \\
\bottomrule
\end{longtable}

\eqitem{Peso total de la losa}{}{W_{\mathrm{l}} = e_a \cdot A_{\mathrm{tp}} \cdot L_{\mathrm{t}} \cdot \gamma_{ha} = \mathbf{144}\,\mathrm{tonf}}

\subsection{Peso de las vigas de acero}

\eqitem{Altura de la viga (0.04 L, AASHTO T. 2.5.2.6.3-1)}{}{h_{\mathrm{vm}} = 0.04 \cdot L_{\mathrm{t}} = \mathbf{1.200}\,\mathrm{m}}

\begin{longtable}{@{}l p{0.52\linewidth} r l l@{}}
\toprule
\textbf{Dato} & \textbf{Descripción} & \textbf{Valor} & \textbf{Unidad} & \textbf{Ref.}\\
\midrule
$h_{\mathrm{vma}}$ & Altura de viga adoptada & \cellcolor{fondodato}1.200 & $\mathrm{m}$ & {\scriptsize } \\
$\gamma_a$ & Peso unitario del acero & \cellcolor{fondodato}7.850 & $\mathrm{tonf/m^{3}}$ & {\scriptsize } \\
\bottomrule
\end{longtable}

\begin{center}
{\small\textbf{Sección de la viga}}\\[3pt]
\adjustbox{max width=\linewidth}{\begin{tabular}{cccc}
\toprule
\begin{tabular}[b]{@{}c@{}}\textbf{Sección}\end{tabular} & \begin{tabular}[b]{@{}c@{}}\textbf{D1}\\ {\scriptsize $\mathrm{m}$}\end{tabular} & \begin{tabular}[b]{@{}c@{}}\textbf{D2}\\ {\scriptsize $\mathrm{m}$}\end{tabular} & \begin{tabular}[b]{@{}c@{}}\textbf{Área}\\ {\scriptsize $\mathrm{m^{2}}$}\end{tabular} \\
\midrule
Patín sup. & \cellcolor{fondodato}$0.3500$ & \cellcolor{fondodato}$0.0254$ & $8.890\times10^{-3}$ \\
Alma & \cellcolor{fondodato}$1.350$ & \cellcolor{fondodato}$0.0125$ & $0.0169$ \\
Patín inf. & \cellcolor{fondodato}$0.3500$ & \cellcolor{fondodato}$0.0254$ & $8.890\times10^{-3}$ \\
\bottomrule
\end{tabular}}
\end{center}

\eqitem{Sección de la viga (redondeada)}{}{S_{\mathrm{vm}} = \operatorname{red}\left(S_{\mathrm{a,1}} + S_{\mathrm{a,2}} + S_{\mathrm{a,3}}\right) = \mathbf{0.0347}\,\mathrm{m^{2}}}

\eqitem{Peso de una viga}{}{w_{\mathrm{vm}} = \gamma_a \cdot S_{\mathrm{vm}} \cdot L_{\mathrm{t}} = \mathbf{8.172}\,\mathrm{tonf}}

\begin{longtable}{@{}l p{0.52\linewidth} r l l@{}}
\toprule
\textbf{Dato} & \textbf{Descripción} & \textbf{Valor} & \textbf{Unidad} & \textbf{Ref.}\\
\midrule
$N_{\mathrm{vm}}$ & Número de vigas & \cellcolor{fondodato}4 & $$ & {\scriptsize } \\
\bottomrule
\end{longtable}

\eqitem{Peso total de las vigas}{}{W_{\mathrm{tvm}} = w_{\mathrm{vm}} \cdot N_{\mathrm{vm}} = \mathbf{32.687}\,\mathrm{tonf}}

\eqitem{Peso del diafragma metálico (5 \% de las vigas)}{}{W_{\mathrm{d}} = 0.05 \cdot W_{\mathrm{tvm}} = \mathbf{1.634}\,\mathrm{tonf}}

\subsection{Peso de barrera}

\nota{Sección de barrera descompuesta en rectángulos y triángulos; brazo medido desde el borde exterior.}

\begin{center}
{\small\textbf{Sección de la barrera}}\\[3pt]
\adjustbox{max width=\linewidth}{\begin{tabular}{cccccccc}
\toprule
\begin{tabular}[b]{@{}c@{}}\textbf{Elemento}\end{tabular} & \begin{tabular}[b]{@{}c@{}}\textbf{Base}\\ {\scriptsize $\mathrm{m}$}\end{tabular} & \begin{tabular}[b]{@{}c@{}}\textbf{Altura}\\ {\scriptsize $\mathrm{m}$}\end{tabular} & \begin{tabular}[b]{@{}c@{}}\textbf{Factor forma}\end{tabular} & \begin{tabular}[b]{@{}c@{}}\textbf{Área}\\ {\scriptsize $\mathrm{m^{2}}$}\end{tabular} & \begin{tabular}[b]{@{}c@{}}\textbf{Peso}\\ {\scriptsize $\mathrm{tonf/m}$}\end{tabular} & \begin{tabular}[b]{@{}c@{}}\textbf{C.G.}\\ {\scriptsize $\mathrm{m}$}\end{tabular} & \begin{tabular}[b]{@{}c@{}}\textbf{Momento}\\ {\scriptsize $\mathrm{tonf\cdot m/m}$}\end{tabular} \\
\midrule
A1 & \cellcolor{fondodato}$0.1500$ & \cellcolor{fondodato}$0.5250$ & \cellcolor{fondodato}$1$ & $0.0788$ & $0.1890$ & $0.0750$ & $0.0142$ \\
A2 & \cellcolor{fondodato}$0.0500$ & \cellcolor{fondodato}$0.5250$ & \cellcolor{fondodato}$0.5000$ & $0.0131$ & $0.0315$ & $0.1667$ & $5.250\times10^{-3}$ \\
A3 & \cellcolor{fondodato}$0.2000$ & \cellcolor{fondodato}$0.2500$ & \cellcolor{fondodato}$1$ & $0.0500$ & $0.1200$ & $0.1000$ & $0.0120$ \\
A4 & \cellcolor{fondodato}$0.1750$ & \cellcolor{fondodato}$0.2500$ & \cellcolor{fondodato}$0.5000$ & $0.0219$ & $0.0525$ & $0.2583$ & $0.0136$ \\
A5 & \cellcolor{fondodato}$0.3750$ & \cellcolor{fondodato}$0.0750$ & \cellcolor{fondodato}$1$ & $0.0281$ & $0.0675$ & $0.1875$ & $0.0127$ \\
\bottomrule
\end{tabular}}
\end{center}

\eqitem{Área de la barrera}{}{A_{\mathrm{bar}} = \sum A_i = \mathbf{0.1919}\,\mathrm{m^{2}}}

\eqitem{Peso de la barrera}{}{w_b = \sum w_i = \mathbf{0.4605}\,\mathrm{tonf/m}}

\eqitem{Momento respecto al borde}{}{M_{\mathrm{bar}} = \sum w_i x_i = \mathbf{0.0576}\,\mathrm{tonf\cdot m/m}}

\eqitem{C.G. de la barrera}{}{x_{\mathrm{cg,b}} = \dfrac{M_{\mathrm{bar}}}{w_b} = \mathbf{0.1252}\,\mathrm{m}}

\eqitem{Peso total de las barreras}{}{W_{\mathrm{b}} = w_b \cdot 2 \cdot L_{\mathrm{t}} = \mathbf{27.630}\,\mathrm{tonf}}

\subsection{Peso de carpeta asfáltica}

\eqitem{Ancho de la barrera}{}{A_{\mathrm{b}} = b_{\mathrm{5}} = \mathbf{0.3750}\,\mathrm{m}}

\eqitem{Ancho de carpeta asfáltica}{}{A_{\mathrm{ca}} = A_{\mathrm{tp}} - 2 \cdot A_{\mathrm{b}} = \mathbf{9.250}\,\mathrm{m}}

\begin{longtable}{@{}l p{0.52\linewidth} r l l@{}}
\toprule
\textbf{Dato} & \textbf{Descripción} & \textbf{Valor} & \textbf{Unidad} & \textbf{Ref.}\\
\midrule
$e_{\mathrm{ca}}$ & Espesor de carpeta asfáltica & \cellcolor{fondodato}0.0500 & $\mathrm{m}$ & {\scriptsize } \\
$\gamma_{asf}$ & Peso unitario del hormigón asfáltico & \cellcolor{fondodato}2.300 & $\mathrm{tonf/m^{3}}$ & {\scriptsize } \\
\bottomrule
\end{longtable}

\eqitem{Peso de carpeta asfáltica}{}{W_{\mathrm{ca}} = e_{\mathrm{ca}} \cdot A_{\mathrm{ca}} \cdot \gamma_{asf} \cdot L_{\mathrm{t}} = \mathbf{31.912}\,\mathrm{tonf}}

\subsection{Peso de instalaciones}

\begin{longtable}{@{}l p{0.52\linewidth} r l l@{}}
\toprule
\textbf{Dato} & \textbf{Descripción} & \textbf{Valor} & \textbf{Unidad} & \textbf{Ref.}\\
\midrule
$w_{inst}$ & Postes de alumbrado y demás & \cellcolor{fondodato}0.0300 & $\mathrm{tonf/m^{2}}$ & {\scriptsize } \\
\bottomrule
\end{longtable}

\eqitem{Peso total de instalaciones}{}{W_{inst,total} = w_{inst} \cdot L_{\mathrm{t}} \cdot A_{\mathrm{tp}} = \mathbf{9}\,\mathrm{tonf}}

\subsection{Peso total de los componentes}

\eqitem{Peso total}{}{W_{\mathrm{total}} = W_{\mathrm{tvm}} + W_{\mathrm{d}} + W_{\mathrm{b}} + W_{\mathrm{ca}} + W_{\mathrm{l}} + W_{inst,total} = \mathbf{246.86}\,\mathrm{tonf}}

\eqitem{Peso aferente a la pila}{}{W = \dfrac{W_{\mathrm{total}}}{L_{\mathrm{t}} \cdot A_{\mathrm{tp}}} = \mathbf{0.8229}\,\mathrm{tonf/m^{2}}}

\clearpage
\section{Descarga CM sobre una pila - vigas preforzadas}

\nota{Descarga por carga muerta de la superestructura (vigas AASHTO IV) sobre una pila.}

\subsection{Datos}

\begin{longtable}{@{}l p{0.52\linewidth} r l l@{}}
\toprule
\textbf{Dato} & \textbf{Descripción} & \textbf{Valor} & \textbf{Unidad} & \textbf{Ref.}\\
\midrule
$L_{\mathrm{p}}$ & Longitud del puente & \cellcolor{fondodato}120 & $\mathrm{m}$ & {\scriptsize } \\
$L_{\mathrm{t}}$ & Longitud del tramo & \cellcolor{fondodato}30 & $\mathrm{m}$ & {\scriptsize } \\
$A_{\mathrm{tp}}$ & Ancho del tablero & \cellcolor{fondodato}10 & $\mathrm{m}$ & {\scriptsize } \\
\bottomrule
\end{longtable}

\subsection{Altura y separación de vigas}

\begin{center}
{\small\textbf{Distribución transversal}}\\[3pt]
\adjustbox{max width=\linewidth}{\begin{tabular}{cccc}
\toprule
\begin{tabular}[b]{@{}c@{}}\textbf{Descripción}\end{tabular} & \begin{tabular}[b]{@{}c@{}}\textbf{Unidad}\\ {\scriptsize $\mathrm{m}$}\end{tabular} & \begin{tabular}[b]{@{}c@{}}\textbf{Cantidad}\end{tabular} & \begin{tabular}[b]{@{}c@{}}\textbf{Total}\\ {\scriptsize $\mathrm{m}$}\end{tabular} \\
\midrule
Separación vigas & \cellcolor{fondodato}$2.500$ & \cellcolor{fondodato}$3$ & $7.500$ \\
Longitud voladizo & \cellcolor{fondodato}$1.250$ & \cellcolor{fondodato}$2$ & $2.500$ \\
Ancho total puente &  &  & $10$ \\
\bottomrule
\end{tabular}}
\end{center}

\eqitem{Distribución transversal = ancho del tablero}{}{a_{\mathrm{n,t}} \ge A_{\mathrm{tp}} - 0.001 \;\wedge\; a_{\mathrm{n,t}} \le A_{\mathrm{tp}} + 0.001\;\;\Rightarrow\;\;\text{\ok{OK}}}

\subsection{Peso de la losa}

\begin{longtable}{@{}l p{0.52\linewidth} r l l@{}}
\toprule
\textbf{Dato} & \textbf{Descripción} & \textbf{Valor} & \textbf{Unidad} & \textbf{Ref.}\\
\midrule
$e_l$ & Espesor mínimo de losa & \cellcolor{fondodato}7 & $\mathrm{in}$ & {\scriptsize } \\
$e_a$ & Espesor de losa asumido & \cellcolor{fondodato}0.2000 & $\mathrm{m}$ & {\scriptsize } \\
\bottomrule
\end{longtable}

\eqitem{Espesor asumido $\ge$ mínimo}{}{e_a \ge e_l\;\;\Rightarrow\;\;\text{\ok{OK}}}

\begin{longtable}{@{}l p{0.52\linewidth} r l l@{}}
\toprule
\textbf{Dato} & \textbf{Descripción} & \textbf{Valor} & \textbf{Unidad} & \textbf{Ref.}\\
\midrule
$\gamma_{ha}$ & Peso unitario del hormigón armado & \cellcolor{fondodato}2.400 & $\mathrm{tonf/m^{3}}$ & {\scriptsize } \\
\bottomrule
\end{longtable}

\eqitem{Peso total de la losa}{}{W_{\mathrm{l}} = e_a \cdot A_{\mathrm{tp}} \cdot L_{\mathrm{t}} \cdot \gamma_{ha} = \mathbf{144}\,\mathrm{tonf}}

\subsection{Peso de las vigas preforzadas (AASHTO tipo IV)}

\begin{longtable}{@{}l p{0.52\linewidth} r l l@{}}
\toprule
\textbf{Dato} & \textbf{Descripción} & \textbf{Valor} & \textbf{Unidad} & \textbf{Ref.}\\
\midrule
$w_{\mathrm{vp}}$ & Peso por metro de la viga & \cellcolor{fondodato}0.8220 & $\mathrm{kip/ft}$ & {\scriptsize } \\
\bottomrule
\end{longtable}

\eqitem{Peso de una viga}{}{w_{\mathrm{1vp}} = w_{\mathrm{vp}} \cdot L_{\mathrm{t}} = \mathbf{36.698}\,\mathrm{tonf}}

\begin{longtable}{@{}l p{0.52\linewidth} r l l@{}}
\toprule
\textbf{Dato} & \textbf{Descripción} & \textbf{Valor} & \textbf{Unidad} & \textbf{Ref.}\\
\midrule
$N_{\mathrm{vm}}$ & Número de vigas & \cellcolor{fondodato}4 & $$ & {\scriptsize } \\
\bottomrule
\end{longtable}

\eqitem{Peso total de las vigas}{}{W_{\mathrm{tvm}} = w_{\mathrm{1vp}} \cdot N_{\mathrm{vm}} = \mathbf{146.79}\,\mathrm{tonf}}

\subsection{Peso de diafragmas}

\begin{longtable}{@{}l p{0.52\linewidth} r l l@{}}
\toprule
\textbf{Dato} & \textbf{Descripción} & \textbf{Valor} & \textbf{Unidad} & \textbf{Ref.}\\
\midrule
$S_{\mathrm{v}}$ & Separación de vigas & \cellcolor{fondodato}2.500 & $\mathrm{m}$ & {\scriptsize } \\
$h_{\mathrm{d}}$ & Altura del diafragma & \cellcolor{fondodato}54 & $\mathrm{in}$ & {\scriptsize } \\
$A_{\mathrm{d}}$ & Ancho del diafragma & \cellcolor{fondodato}0.2500 & $\mathrm{m}$ & {\scriptsize } \\
$L_{\mathrm{d}}$ & Longitud del diafragma & \cellcolor{fondodato}2.300 & $\mathrm{m}$ & {\scriptsize } \\
\bottomrule
\end{longtable}

\eqitem{Peso de un diafragma}{}{w_{\mathrm{d}} = h_{\mathrm{d}} \cdot A_{\mathrm{d}} \cdot L_{\mathrm{d}} \cdot \gamma_{ha} = \mathbf{1.893}\,\mathrm{tonf}}

\begin{longtable}{@{}l p{0.52\linewidth} r l l@{}}
\toprule
\textbf{Dato} & \textbf{Descripción} & \textbf{Valor} & \textbf{Unidad} & \textbf{Ref.}\\
\midrule
$N_{\mathrm{d}}$ & Número de diafragmas & \cellcolor{fondodato}12 & $$ & {\scriptsize } \\
\bottomrule
\end{longtable}

\eqitem{Peso total de diafragmas}{}{W_{\mathrm{dha}} = N_{\mathrm{d}} \cdot w_{\mathrm{d}} = \mathbf{22.714}\,\mathrm{tonf}}

\subsection{Peso de barrera}

\nota{Sección de barrera descompuesta en rectángulos y triángulos; brazo medido desde el borde exterior.}

\begin{center}
{\small\textbf{Sección de la barrera}}\\[3pt]
\adjustbox{max width=\linewidth}{\begin{tabular}{cccccccc}
\toprule
\begin{tabular}[b]{@{}c@{}}\textbf{Elemento}\end{tabular} & \begin{tabular}[b]{@{}c@{}}\textbf{Base}\\ {\scriptsize $\mathrm{m}$}\end{tabular} & \begin{tabular}[b]{@{}c@{}}\textbf{Altura}\\ {\scriptsize $\mathrm{m}$}\end{tabular} & \begin{tabular}[b]{@{}c@{}}\textbf{Factor forma}\end{tabular} & \begin{tabular}[b]{@{}c@{}}\textbf{Área}\\ {\scriptsize $\mathrm{m^{2}}$}\end{tabular} & \begin{tabular}[b]{@{}c@{}}\textbf{Peso}\\ {\scriptsize $\mathrm{tonf/m}$}\end{tabular} & \begin{tabular}[b]{@{}c@{}}\textbf{C.G.}\\ {\scriptsize $\mathrm{m}$}\end{tabular} & \begin{tabular}[b]{@{}c@{}}\textbf{Momento}\\ {\scriptsize $\mathrm{tonf\cdot m/m}$}\end{tabular} \\
\midrule
A1 & \cellcolor{fondodato}$0.1500$ & \cellcolor{fondodato}$0.5250$ & \cellcolor{fondodato}$1$ & $0.0788$ & $0.1890$ & $0.0750$ & $0.0142$ \\
A2 & \cellcolor{fondodato}$0.0500$ & \cellcolor{fondodato}$0.5250$ & \cellcolor{fondodato}$0.5000$ & $0.0131$ & $0.0315$ & $0.1667$ & $5.250\times10^{-3}$ \\
A3 & \cellcolor{fondodato}$0.2000$ & \cellcolor{fondodato}$0.2500$ & \cellcolor{fondodato}$1$ & $0.0500$ & $0.1200$ & $0.1000$ & $0.0120$ \\
A4 & \cellcolor{fondodato}$0.1750$ & \cellcolor{fondodato}$0.2500$ & \cellcolor{fondodato}$0.5000$ & $0.0219$ & $0.0525$ & $0.2583$ & $0.0136$ \\
A5 & \cellcolor{fondodato}$0.3750$ & \cellcolor{fondodato}$0.0750$ & \cellcolor{fondodato}$1$ & $0.0281$ & $0.0675$ & $0.1875$ & $0.0127$ \\
\bottomrule
\end{tabular}}
\end{center}

\eqitem{Área de la barrera}{}{A_{\mathrm{bar}} = \sum A_i = \mathbf{0.1919}\,\mathrm{m^{2}}}

\eqitem{Peso de la barrera}{}{w_b = \sum w_i = \mathbf{0.4605}\,\mathrm{tonf/m}}

\eqitem{Momento respecto al borde}{}{M_{\mathrm{bar}} = \sum w_i x_i = \mathbf{0.0576}\,\mathrm{tonf\cdot m/m}}

\eqitem{C.G. de la barrera}{}{x_{\mathrm{cg,b}} = \dfrac{M_{\mathrm{bar}}}{w_b} = \mathbf{0.1252}\,\mathrm{m}}

\eqitem{Peso total de las barreras}{}{W_{\mathrm{b}} = w_b \cdot 2 \cdot L_{\mathrm{t}} = \mathbf{27.630}\,\mathrm{tonf}}

\subsection{Peso de carpeta asfáltica}

\eqitem{Ancho de la barrera}{}{A_{\mathrm{b}} = b_{\mathrm{5}} = \mathbf{0.3750}\,\mathrm{m}}

\eqitem{Ancho de carpeta asfáltica}{}{A_{\mathrm{ca}} = A_{\mathrm{tp}} - 2 \cdot A_{\mathrm{b}} = \mathbf{9.250}\,\mathrm{m}}

\begin{longtable}{@{}l p{0.52\linewidth} r l l@{}}
\toprule
\textbf{Dato} & \textbf{Descripción} & \textbf{Valor} & \textbf{Unidad} & \textbf{Ref.}\\
\midrule
$e_{\mathrm{ca}}$ & Espesor de carpeta asfáltica & \cellcolor{fondodato}0.0500 & $\mathrm{m}$ & {\scriptsize } \\
$\gamma_{asf}$ & Peso unitario del hormigón asfáltico & \cellcolor{fondodato}2.300 & $\mathrm{tonf/m^{3}}$ & {\scriptsize } \\
\bottomrule
\end{longtable}

\eqitem{Peso de carpeta asfáltica}{}{W_{\mathrm{ca}} = e_{\mathrm{ca}} \cdot A_{\mathrm{ca}} \cdot \gamma_{asf} \cdot L_{\mathrm{t}} = \mathbf{31.912}\,\mathrm{tonf}}

\subsection{Peso de instalaciones}

\begin{longtable}{@{}l p{0.52\linewidth} r l l@{}}
\toprule
\textbf{Dato} & \textbf{Descripción} & \textbf{Valor} & \textbf{Unidad} & \textbf{Ref.}\\
\midrule
$w_{inst}$ & Postes de alumbrado y demás & \cellcolor{fondodato}0.0300 & $\mathrm{tonf/m^{2}}$ & {\scriptsize } \\
\bottomrule
\end{longtable}

\eqitem{Peso total de instalaciones}{}{W_{inst,total} = w_{inst} \cdot L_{\mathrm{t}} \cdot A_{\mathrm{tp}} = \mathbf{9}\,\mathrm{tonf}}

\subsection{Peso total de los componentes}

\eqitem{Peso total}{}{W_{\mathrm{total}} = W_{\mathrm{tvm}} + W_{\mathrm{dha}} + W_{\mathrm{b}} + W_{\mathrm{ca}} + W_{\mathrm{l}} + W_{inst,total} = \mathbf{382.05}\,\mathrm{tonf}}

\eqitem{Peso aferente a la pila}{}{W = \dfrac{W_{\mathrm{total}}}{L_{\mathrm{t}} \cdot A_{\mathrm{tp}}} = \mathbf{1.273}\,\mathrm{tonf/m^{2}}}

\clearpage
\section{Diseño de viga cabezal (método ASD)}

\nota{Viga cabezal de pila en voladizo. Carga viva camión HS20. AASHTO Estándar, método ASD.}

\subsection{Datos}

\begin{longtable}{@{}l p{0.52\linewidth} r l l@{}}
\toprule
\textbf{Dato} & \textbf{Descripción} & \textbf{Valor} & \textbf{Unidad} & \textbf{Ref.}\\
\midrule
$\gamma_{ha}$ & Peso específico del hormigón armado & \cellcolor{fondodato}2\,400 & $\mathrm{kgf/m^{3}}$ & {\scriptsize } \\
$f'_c$ & f'c del hormigón & \cellcolor{fondodato}210 & $\mathrm{kgf/cm^{2}}$ & {\scriptsize } \\
\bottomrule
\end{longtable}

\nota{Carga viva: camión HS20.}

\begin{longtable}{@{}l p{0.52\linewidth} r l l@{}}
\toprule
\textbf{Dato} & \textbf{Descripción} & \textbf{Valor} & \textbf{Unidad} & \textbf{Ref.}\\
\midrule
$P_{\mathrm{1}}$ & Carga del eje delantero & \cellcolor{fondodato}8 & $\mathrm{kip}$ & {\scriptsize } \\
$P_{\mathrm{2}}$ & Carga del eje central & \cellcolor{fondodato}32 & $\mathrm{kip}$ & {\scriptsize } \\
$P_{\mathrm{3}}$ & Carga del eje trasero & \cellcolor{fondodato}32 & $\mathrm{kip}$ & {\scriptsize } \\
$L_1$ & Distancia entre ejes P1 y P2 & \cellcolor{fondodato}14 & $\mathrm{ft}$ & {\scriptsize } \\
$L_2$ & Distancia entre ejes P2 y P3 & \cellcolor{fondodato}14 & $\mathrm{ft}$ & {\scriptsize } \\
$\mathit{LT}_{\mathrm{1}}$ & Longitud del tramo 1 & \cellcolor{fondodato}18 & $\mathrm{m}$ & {\scriptsize } \\
$\mathit{LT}_{\mathrm{2}}$ & Longitud del tramo 2 & \cellcolor{fondodato}18 & $\mathrm{m}$ & {\scriptsize } \\
$V_{\mathrm{ext}}$ & Carga muerta de viga exterior & \cellcolor{fondodato}1\,637 & $\mathrm{kgf/m}$ & {\scriptsize } \\
$V_{\mathrm{int}}$ & Carga muerta de viga interior & \cellcolor{fondodato}1\,786 & $\mathrm{kgf/m}$ & {\scriptsize } \\
\bottomrule
\end{longtable}

\subsection{Cargas por línea de ruedas}

\eqitem{Eje P1}{}{P_{\mathrm{1r}} = P_{\mathrm{1}} \cdot 0.5 = \mathbf{1\,814.37}\,\mathrm{kgf}}

\eqitem{Eje P2}{}{P_{\mathrm{2r}} = P_{\mathrm{2}} \cdot 0.5 = \mathbf{7\,257.48}\,\mathrm{kgf}}

\eqitem{Eje P3}{}{P_{\mathrm{3r}} = P_{\mathrm{3}} \cdot 0.5 = \mathbf{7\,257.48}\,\mathrm{kgf}}

\subsection{Reacción de la pila debido a carga viva}

\eqitem{Debido a P1}{}{\mathit{RP}_{\mathrm{1}} = \dfrac{P_{\mathrm{1r}} \cdot \left(\mathit{LT}_{\mathrm{1}} - L_1\right)}{\mathit{LT}_{\mathrm{1}}} = \mathbf{1\,384.24}\,\mathrm{kgf}}

\eqitem{Debido a P2 (sobre la pila)}{}{\mathit{RP}_{\mathrm{2}} = \dfrac{P_{\mathrm{2r}} \cdot \mathit{LT}_{\mathrm{1}}}{\mathit{LT}_{\mathrm{1}}} = \mathbf{7\,257.48}\,\mathrm{kgf}}

\eqitem{Debido a P3}{}{\mathit{RP}_{\mathrm{3}} = \dfrac{P_{\mathrm{3r}} \cdot \left(\mathit{LT}_{\mathrm{2}} - L_2\right)}{\mathit{LT}_{\mathrm{2}}} = \mathbf{5\,536.97}\,\mathrm{kgf}}

\eqitem{Reacción de la pila}{}{\mathit{RP} = \mathit{RP}_{\mathrm{1}} + \mathit{RP}_{\mathrm{2}} + \mathit{RP}_{\mathrm{3}} = \mathbf{14\,178.69}\,\mathrm{kgf}}

\eqitem{Carga viva por rueda}{}{w = \mathit{RP} = \mathbf{14\,178.69}\,\mathrm{kgf}}

\subsection{Reacciones en la pila debido a la carga muerta}

\eqitem{Viga exterior, tramo 1}{}{V_{\mathrm{ext1cm}} = 0.5 \cdot V_{\mathrm{ext}} \cdot \mathit{LT}_{\mathrm{1}} = \mathbf{14\,733}\,\mathrm{kgf}}

\eqitem{Viga exterior, tramo 2}{}{V_{\mathrm{ext2cm}} = 0.5 \cdot V_{\mathrm{ext}} \cdot \mathit{LT}_{\mathrm{2}} = \mathbf{14\,733}\,\mathrm{kgf}}

\eqitem{Viga exterior (2 tramos), V1}{}{V_{\mathrm{1cm}} = V_{\mathrm{ext1cm}} + V_{\mathrm{ext2cm}} = \mathbf{29\,466}\,\mathrm{kgf}}

\eqitem{Viga interior, tramo 1}{}{V_{\mathrm{int1cm}} = 0.5 \cdot V_{\mathrm{int}} \cdot \mathit{LT}_{\mathrm{1}} = \mathbf{16\,074}\,\mathrm{kgf}}

\eqitem{Viga interior, tramo 2}{}{V_{\mathrm{int2cm}} = 0.5 \cdot V_{\mathrm{int}} \cdot \mathit{LT}_{\mathrm{2}} = \mathbf{16\,074}\,\mathrm{kgf}}

\eqitem{Viga interior (2 tramos), V2}{}{V_{\mathrm{2cm}} = V_{\mathrm{int1cm}} + V_{\mathrm{int2cm}} = \mathbf{32\,148}\,\mathrm{kgf}}

\subsection{Geometría de la viga cabezal}

\begin{longtable}{@{}l p{0.52\linewidth} r l l@{}}
\toprule
\textbf{Dato} & \textbf{Descripción} & \textbf{Valor} & \textbf{Unidad} & \textbf{Ref.}\\
\midrule
$\mathit{LT}$ & Longitud del tramo (valor numérico en m) & \cellcolor{fondodato}18 & $$ & {\scriptsize } \\
$H_{\mathrm{p}}$ & Altura de la pila (valor numérico en m) & \cellcolor{fondodato}9.800 & $$ & {\scriptsize } \\
$s$ & Esviaje (grados) & \cellcolor{fondodato}0 & $$ & {\scriptsize } \\
\bottomrule
\end{longtable}

\eqitem{Longitud de apoyo N (mm)}{}{N_{\mathrm{ap1}} = \left(305 + 2.5 \cdot \mathit{LT} + 10 \cdot H_{\mathrm{p}}\right) \cdot \left(1 + 0.000125 \cdot {s}^{2}\right) = \mathbf{448}}

\eqitem{Apoyo 1 (mm)}{}{N_{\mathrm{as1}} = N_{\mathrm{ap1}} = \mathbf{448}}

\eqitem{Apoyo 2 (mm)}{}{N_{\mathrm{as2}} = N_{\mathrm{ap1}} = \mathbf{448}}

\eqitem{Ancho mínimo de viga cabezal (mm)}{}{b_{\mathrm{vc}} = N_{\mathrm{as1}} + N_{\mathrm{as2}} + 75 = \mathbf{971}}

\begin{longtable}{@{}l p{0.52\linewidth} r l l@{}}
\toprule
\textbf{Dato} & \textbf{Descripción} & \textbf{Valor} & \textbf{Unidad} & \textbf{Ref.}\\
\midrule
$b$ & Ancho de viga cabezal asumido & \cellcolor{fondodato}1.200 & $\mathrm{m}$ & {\scriptsize } \\
\bottomrule
\end{longtable}

\eqitem{Ancho asumido $\ge$ ancho mínimo}{}{b \ge b_{\mathrm{vc}} \cdot \mathrm{mm}\;\;\Rightarrow\;\;\text{\ok{OK}}}

\begin{longtable}{@{}l p{0.52\linewidth} r l l@{}}
\toprule
\textbf{Dato} & \textbf{Descripción} & \textbf{Valor} & \textbf{Unidad} & \textbf{Ref.}\\
\midrule
$\mathit{HVC}_{\mathrm{ex}}$ & Altura de viga cabezal en el extremo & \cellcolor{fondodato}0.9000 & $\mathrm{m}$ & {\scriptsize } \\
$\mathit{HVC}_{\mathrm{in}}$ & Altura de viga cabezal en la cara de la columna & \cellcolor{fondodato}2.100 & $\mathrm{m}$ & {\scriptsize } \\
$\mathit{VVC}$ & Volado de la viga cabezal & \cellcolor{fondodato}3.500 & $\mathrm{m}$ & {\scriptsize } \\
$\mathit{AD}_{\mathrm{p}}$ & Ancho / diámetro de pila & \cellcolor{fondodato}2.400 & $\mathrm{m}$ & {\scriptsize } \\
$S_{\mathrm{v}}$ & Separación entre vigas & \cellcolor{fondodato}2 & $\mathrm{m}$ & {\scriptsize } \\
$\mathit{DV}_{\mathrm{1vol}}$ & Viga 1 (exterior) al extremo de la viga cabezal & \cellcolor{fondodato}0.7000 & $\mathrm{m}$ & {\scriptsize } \\
$D_{\mathrm{ca}}$ & Distancia entre ejes de ruedas del camión & \cellcolor{fondodato}1.800 & $\mathrm{m}$ & {\scriptsize } \\
$D_{\mathrm{erint}}$ & Eje de viga cabezal a línea de ruedas interior & \cellcolor{fondodato}1.500 & $\mathrm{m}$ & {\scriptsize } \\
\bottomrule
\end{longtable}

\subsection{Distribución de carga viva}

\eqitem{Línea de ruedas exterior: viga interior}{}{V_{\mathrm{2ext}} = \dfrac{w \cdot \left(2 \cdot S_{\mathrm{v}} - \left(D_{\mathrm{erint}} + D_{\mathrm{ca}}\right)\right)}{S_{\mathrm{v}}} = \mathbf{4\,962.54}\,\mathrm{kgf}}

\eqitem{Línea de ruedas exterior: viga exterior}{}{V_{\mathrm{1ext}} = w - V_{\mathrm{2ext}} = \mathbf{9\,216.15}\,\mathrm{kgf}}

\eqitem{Línea de ruedas interior: viga interior 3}{}{V_{\mathrm{3int}} = \dfrac{w \cdot \left(S_{\mathrm{v}} - D_{\mathrm{erint}}\right)}{S_{\mathrm{v}}} = \mathbf{3\,544.67}\,\mathrm{kgf}}

\eqitem{Línea de ruedas interior: viga interior 2}{}{V_{\mathrm{2int}} = w - V_{\mathrm{3int}} = \mathbf{10\,634.02}\,\mathrm{kgf}}

\subsection{Fuerza y esfuerzo de corte [STD A 8.22 / A 8.15.5.1 / A 8.15.5.2]}

\begin{longtable}{@{}l p{0.52\linewidth} r l l@{}}
\toprule
\textbf{Dato} & \textbf{Descripción} & \textbf{Valor} & \textbf{Unidad} & \textbf{Ref.}\\
\midrule
$r_{\mathrm{ec}}$ & Recubrimiento & \cellcolor{fondodato}5 & $\mathrm{cm}$ & {\scriptsize } \\
$R_{\mathrm{efT}}$ & Diámetro del refuerzo transversal & \cellcolor{fondodato}1.200 & $\mathrm{cm}$ & {\scriptsize } \\
$R_{\mathrm{efL}}$ & Diámetro del refuerzo longitudinal & \cellcolor{fondodato}2.200 & $\mathrm{cm}$ & {\scriptsize } \\
\bottomrule
\end{longtable}

\eqitem{Esfuerzo cortante resistido por el hormigón}{}{v_{\mathrm{h}} = 0.079 \cdot \sqrt{\dfrac{f'_c}{\mathrm{MPa}}} \cdot \mathrm{MPa} = \mathbf{3.656}\,\mathrm{kgf/cm^{2}}}

\subsubsection{Sección a-a}

\eqitem{Altura de la sección a-a}{}{H_{\mathrm{sa}} = \mathit{HVC}_{\mathrm{ex}} + \dfrac{\mathit{DV}_{\mathrm{1vol}} \cdot \left(\mathit{HVC}_{\mathrm{in}} - \mathit{HVC}_{\mathrm{ex}}\right)}{\mathit{VVC}} = \mathbf{1.140}\,\mathrm{m}}

\eqitem{Peso de viga cabezal}{}{W_{\mathrm{a}} = 0.5 \cdot \left(\mathit{HVC}_{\mathrm{ex}} + H_{\mathrm{sa}}\right) \cdot \mathit{DV}_{\mathrm{1vol}} \cdot b \cdot \gamma_{ha} = \mathbf{2\,056.32}\,\mathrm{kgf}}

\eqitem{Fuerza cortante}{}{V_{\mathrm{a}} = V_{\mathrm{1cm}} + V_{\mathrm{1ext}} + W_{\mathrm{a}} = \mathbf{40\,738.47}\,\mathrm{kgf}}

\eqitem{Peralte efectivo}{}{d_{\mathrm{a}} = H_{\mathrm{sa}} - \left(r_{\mathrm{ec}} + R_{\mathrm{efT}} + 0.5 \cdot R_{\mathrm{efL}}\right) = \mathbf{106.70}\,\mathrm{cm}}

\eqitem{Esfuerzo cortante}{}{v_{\mathrm{a}} = \dfrac{V_{\mathrm{a}}}{b \cdot d_{\mathrm{a}}} = \mathbf{3.182}\,\mathrm{kgf/cm^{2}}}

\eqitem{Sección a-a}{}{v_{\mathrm{h}} < v_{\mathrm{a}}\;\;\Rightarrow\;\;\text{\info{No Requiere Ref. Transversal}}}

\subsubsection{Sección b-b}

\eqitem{Altura de la sección b-b}{}{H_{\mathrm{sb}} = \mathit{HVC}_{\mathrm{ex}} + \dfrac{\left(\mathit{DV}_{\mathrm{1vol}} + S_{\mathrm{v}}\right) \cdot \left(\mathit{HVC}_{\mathrm{in}} - \mathit{HVC}_{\mathrm{ex}}\right)}{\mathit{VVC}} = \mathbf{1.826}\,\mathrm{m}}

\eqitem{Peso de viga cabezal}{}{W_b = 0.5 \cdot \left(\mathit{HVC}_{\mathrm{ex}} + H_{\mathrm{sb}}\right) \cdot \left(\mathit{DV}_{\mathrm{1vol}} + S_{\mathrm{v}}\right) \cdot b \cdot \gamma_{ha} = \mathbf{10\,597.58}\,\mathrm{kgf}}

\eqitem{Fuerza cortante}{}{V_{\mathrm{b}} = V_{\mathrm{1cm}} + V_{\mathrm{2cm}} + V_{\mathrm{1ext}} + V_{\mathrm{2ext}} + V_{\mathrm{2int}} + W_b = \mathbf{97\,024.29}\,\mathrm{kgf}}

\eqitem{Peralte efectivo}{}{d_{\mathrm{b}} = H_{\mathrm{sb}} - \left(r_{\mathrm{ec}} + R_{\mathrm{efT}} + 0.5 \cdot R_{\mathrm{efL}}\right) = \mathbf{175.27}\,\mathrm{cm}}

\eqitem{Esfuerzo cortante}{}{v_{\mathrm{b}} = \dfrac{V_{\mathrm{b}}}{b \cdot d_{\mathrm{b}}} = \mathbf{4.613}\,\mathrm{kgf/cm^{2}}}

\eqitem{Sección b-b}{}{v_{\mathrm{h}} < v_{\mathrm{b}}\;\;\Rightarrow\;\;\text{\info{Requiere Ref. Transversal}}}

\subsubsection{Sección c-c}

\eqitem{Altura de la sección c-c}{}{H_{\mathrm{sc}} = \mathit{HVC}_{\mathrm{ex}} + \mathit{HVC}_{\mathrm{in}} - \mathit{HVC}_{\mathrm{ex}} = \mathbf{2.100}\,\mathrm{m}}

\eqitem{Peso de viga cabezal}{}{W_c = 0.5 \cdot \left(\mathit{HVC}_{\mathrm{ex}} + \mathit{HVC}_{\mathrm{in}}\right) \cdot \mathit{VVC} \cdot b \cdot \gamma_{ha} = \mathbf{15\,120}\,\mathrm{kgf}}

\eqitem{Fuerza cortante}{}{V_{\mathrm{c}} = V_{\mathrm{1cm}} + V_{\mathrm{2cm}} + V_{\mathrm{1ext}} + V_{\mathrm{2ext}} + V_{\mathrm{2int}} + W_c = \mathbf{101\,546.71}\,\mathrm{kgf}}

\eqitem{Peralte efectivo}{}{d_{\mathrm{c}} = H_{\mathrm{sc}} - \left(r_{\mathrm{ec}} + R_{\mathrm{efT}} + 0.5 \cdot R_{\mathrm{efL}}\right) = \mathbf{202.70}\,\mathrm{cm}}

\eqitem{Esfuerzo cortante}{}{v_{\mathrm{c}} = \dfrac{V_{\mathrm{c}}}{b \cdot d_{\mathrm{c}}} = \mathbf{4.175}\,\mathrm{kgf/cm^{2}}}

\eqitem{Sección c-c}{}{v_{\mathrm{h}} < v_{\mathrm{c}}\;\;\Rightarrow\;\;\text{\info{Requiere Ref. Transversal}}}

\subsection{Separación del refuerzo transversal [STD A 8.15.5.3 / A 8.15.2.2]}

\nota{Av = (v - vh)$\cdot$bw$\cdot$s / fs}

\begin{longtable}{@{}l p{0.52\linewidth} r l l@{}}
\toprule
\textbf{Dato} & \textbf{Descripción} & \textbf{Valor} & \textbf{Unidad} & \textbf{Ref.}\\
\midrule
$f_{\mathrm{y}}$ & Esfuerzo de fluencia & \cellcolor{fondodato}420 & $\mathrm{MPa}$ & {\scriptsize } \\
$f_{\mathrm{s}}$ & Esfuerzo admisible del acero & \cellcolor{fondodato}24\,000 & $\mathrm{psi}$ & {\scriptsize } \\
\bottomrule
\end{longtable}

\eqitem{Área de una barra}{}{A_{\mathrm{b}} = 0.25 \cdot \pi \cdot {R_{\mathrm{efT}}}^{2} = \mathbf{1.131}\,\mathrm{cm^{2}}}

\eqitem{Esfuerzo de corte máximo}{}{E_{\mathrm{sfMax}} = \max\left(v_{\mathrm{a}},\ v_{\mathrm{b}},\ v_{\mathrm{c}}\right) = \mathbf{4.613}\,\mathrm{kgf/cm^{2}}}

\eqitem{Diferencia de esfuerzos (v - vh)}{}{v_{\mathrm{max}} = E_{\mathrm{sfMax}} - v_{\mathrm{h}} = \mathbf{0.9573}\,\mathrm{kgf/cm^{2}}}

\eqitem{Límite 1}{}{v_{\mathrm{1}} = 0.17 \cdot \sqrt{\dfrac{f'_c}{\mathrm{MPa}}} \cdot \mathrm{MPa} = \mathbf{7.867}\,\mathrm{kgf/cm^{2}}}

\eqitem{Límite 2}{}{v_{\mathrm{2}} = 0.33 \cdot \sqrt{\dfrac{f'_c}{\mathrm{MPa}}} \cdot \mathrm{MPa} = \mathbf{15.271}\,\mathrm{kgf/cm^{2}}}

\eqitem{(v - vh) $\le$ límite 1}{}{v_{\mathrm{max}} \le v_{\mathrm{1}}\;\;\Rightarrow\;\;\text{\ok{CUMPLE}}}

\eqitem{(v - vh) $\le$ límite 2}{}{v_{\mathrm{max}} \le v_{\mathrm{2}}\;\;\Rightarrow\;\;\text{\ok{CUMPLE}}}

\eqitem{Separación requerida (2 ramales)}{}{S = \dfrac{2 \cdot A_{\mathrm{b}} \cdot f_{\mathrm{s}}}{v_{\mathrm{max}} \cdot b} = \mathbf{33.225}\,\mathrm{cm}}

\begin{longtable}{@{}l p{0.52\linewidth} r l l@{}}
\toprule
\textbf{Dato} & \textbf{Descripción} & \textbf{Valor} & \textbf{Unidad} & \textbf{Ref.}\\
\midrule
$S_{\mathrm{esc}}$ & Separación escogida del refuerzo transversal & \cellcolor{fondodato}30 & $\mathrm{cm}$ & {\scriptsize } \\
\bottomrule
\end{longtable}

\eqitem{Separación escogida $\le$ requerida}{}{S_{\mathrm{esc}} \le S\;\;\Rightarrow\;\;\text{\ok{CUMPLE}}}

\subsection{Refuerzo por flexión (sección c-c)}

\eqitem{Peso propio de la sección rectangular}{}{P_{\mathrm{p1}} = \mathit{HVC}_{\mathrm{ex}} \cdot \mathit{VVC} \cdot b \cdot \gamma_{ha} = \mathbf{9\,072}\,\mathrm{kgf}}

\eqitem{Peso propio de la sección triangular}{}{P_{\mathrm{p2}} = 0.5 \cdot \mathit{VVC} \cdot \left(\mathit{HVC}_{\mathrm{in}} - \mathit{HVC}_{\mathrm{ex}}\right) \cdot b \cdot \gamma_{ha} = \mathbf{6\,048}\,\mathrm{kgf}}

\eqitem{Momento sección rectangular}{}{M_{\mathrm{p1}} = P_{\mathrm{p1}} \cdot \mathit{VVC} \cdot 0.5 = \mathbf{15\,876}\,\mathrm{kgf\cdot m}}

\eqitem{Momento sección triangular}{}{M_{\mathrm{p2}} = \dfrac{P_{\mathrm{p2}} \cdot \mathit{VVC}}{3} = \mathbf{7\,056}\,\mathrm{kgf\cdot m}}

\eqitem{Carga muerta V1}{}{M_{\mathrm{v1cm}} = V_{\mathrm{1cm}} \cdot \left(\mathit{VVC} - \mathit{DV}_{\mathrm{1vol}}\right) = \mathbf{82\,504.80}\,\mathrm{kgf\cdot m}}

\eqitem{Carga muerta V2}{}{M_{\mathrm{v2cm}} = V_{\mathrm{2cm}} \cdot \left(\mathit{VVC} - \mathit{DV}_{\mathrm{1vol}} - S_{\mathrm{v}}\right) = \mathbf{25\,718.40}\,\mathrm{kgf\cdot m}}

\eqitem{Carga viva V1}{}{M_{\mathrm{v1cv}} = V_{\mathrm{1ext}} \cdot \left(\mathit{VVC} - \mathit{DV}_{\mathrm{1vol}}\right) = \mathbf{25.805}\,\mathrm{tonf\cdot m}}

\eqitem{Carga viva V2}{}{M_{\mathrm{v2cv}} = \left(V_{\mathrm{2ext}} + V_{\mathrm{2int}}\right) \cdot \left(\mathit{VVC} - \mathit{DV}_{\mathrm{1vol}} - S_{\mathrm{v}}\right) = \mathbf{12.477}\,\mathrm{tonf\cdot m}}

\eqitem{Momento total}{}{M_{\mathrm{t}} = M_{\mathrm{v1cv}} + M_{\mathrm{v2cv}} + M_{\mathrm{v1cm}} + M_{\mathrm{v2cm}} + M_{\mathrm{p1}} + M_{\mathrm{p2}} = \mathbf{169\,437.67}\,\mathrm{kgf\cdot m}}

\eqitem{Área de una varilla longitudinal}{}{\mathit{AR}_{\mathrm{efL}} = 0.25 \cdot \pi \cdot {R_{\mathrm{efL}}}^{2} = \mathbf{3.801}\,\mathrm{cm^{2}}}

\eqitem{Acero requerido}{}{A_{\mathrm{s}} = \dfrac{M_{\mathrm{t}}}{0.875 \cdot f_{\mathrm{s}} \cdot d_{\mathrm{c}}} = \mathbf{56.616}\,\mathrm{cm^{2}}}

\eqitem{Número de varillas requerido}{}{N_{\mathrm{var}} = \dfrac{A_{\mathrm{s}}}{\mathit{AR}_{\mathrm{efL}}} = \mathbf{14.894}}

\eqitem{Número de varillas a emplear}{}{N_{\mathrm{var,c}} = \left\lceil N_{\mathrm{var}}\right\rceil = \mathbf{15}}

\eqitem{Acero total}{}{A_{\mathrm{sT}} = N_{\mathrm{var,c}} \cdot \mathit{AR}_{\mathrm{efL}} = \mathbf{57.020}\,\mathrm{cm^{2}}}

\eqitem{Acero proporcionado $\ge$ requerido}{}{A_{\mathrm{sT}} \ge A_{\mathrm{s}}\;\;\Rightarrow\;\;\text{\ok{CUMPLE}}}

\subsubsection{Comprobación del mínimo ancho de la viga}

\eqitem{1.5 diámetros}{}{c_{\mathrm{av1}} = 1.5 \cdot R_{\mathrm{efL}} = \mathbf{3.300}\,\mathrm{cm}}

\begin{longtable}{@{}l p{0.52\linewidth} r l l@{}}
\toprule
\textbf{Dato} & \textbf{Descripción} & \textbf{Valor} & \textbf{Unidad} & \textbf{Ref.}\\
\midrule
$d_{\mathrm{ag}}$ & Tamaño máximo del agregado grueso (3/4") & \cellcolor{fondodato}1.905 & $\mathrm{cm}$ & {\scriptsize } \\
\bottomrule
\end{longtable}

\eqitem{1.5 veces el agregado}{}{c_{\mathrm{av2}} = 1.5 \cdot d_{\mathrm{ag}} = \mathbf{2.858}\,\mathrm{cm}}

\begin{longtable}{@{}l p{0.52\linewidth} r l l@{}}
\toprule
\textbf{Dato} & \textbf{Descripción} & \textbf{Valor} & \textbf{Unidad} & \textbf{Ref.}\\
\midrule
$c_{\mathrm{av3}}$ & Separación mínima absoluta & \cellcolor{fondodato}4.500 & $\mathrm{cm}$ & {\scriptsize } \\
\bottomrule
\end{longtable}

\eqitem{Mayor separación}{}{S_{\mathrm{mx}} = \max\left(c_{\mathrm{av1}},\ c_{\mathrm{av2}},\ c_{\mathrm{av3}}\right) = \mathbf{4.500}\,\mathrm{cm}}

\eqitem{Ancho mínimo requerido}{}{M_{\mathrm{inav}} = N_{\mathrm{var,c}} \cdot R_{\mathrm{efL}} + \left(N_{\mathrm{var,c}} - 1\right) \cdot S_{\mathrm{mx}} + 2 \cdot R_{\mathrm{efT}} + 2 \cdot r_{\mathrm{ec}} = \mathbf{108.40}\,\mathrm{cm}}

\eqitem{Ancho mínimo < ancho de viga}{}{M_{\mathrm{inav}} < b\;\;\Rightarrow\;\;\text{\ok{CUMPLE}}}

\subsection{Refuerzo de contracción y temperatura [STD A 8.20]}

\eqitem{Cuantía mínima}{}{A_{\mathrm{sct}} = \dfrac{0.125 \cdot {\mathrm{in}}^{2}}{\mathrm{ft}} = \mathbf{2.646}\,\mathrm{cm^{2}/m}}

\eqitem{Acero por lado}{}{A_{\mathrm{sl}} = A_{\mathrm{sct}} \cdot \mathit{HVC}_{\mathrm{in}} = \mathbf{5.556}\,\mathrm{cm^{2}}}

\begin{longtable}{@{}l p{0.52\linewidth} r l l@{}}
\toprule
\textbf{Dato} & \textbf{Descripción} & \textbf{Valor} & \textbf{Unidad} & \textbf{Ref.}\\
\midrule
$d_{\mathrm{basct}}$ & Diámetro del hierro & \cellcolor{fondodato}1.600 & $\mathrm{cm}$ & {\scriptsize } \\
\bottomrule
\end{longtable}

\eqitem{Área de hierro}{}{A_{\mathrm{basct}} = 0.25 \cdot \pi \cdot {d_{\mathrm{basct}}}^{2} = \mathbf{2.011}\,\mathrm{cm^{2}}}

\begin{longtable}{@{}l p{0.52\linewidth} r l l@{}}
\toprule
\textbf{Dato} & \textbf{Descripción} & \textbf{Valor} & \textbf{Unidad} & \textbf{Ref.}\\
\midrule
$n_{\mathrm{var1}}$ & Número de varillas por lado & \cellcolor{fondodato}3 & $$ & {\scriptsize } \\
\bottomrule
\end{longtable}

\eqitem{Área total por lado}{}{A_{\mathrm{sctlado}} = A_{\mathrm{basct}} \cdot n_{\mathrm{var1}} = \mathbf{6.032}\,\mathrm{cm^{2}}}

\eqitem{Cuantía proporcionada}{}{A_{\mathrm{sctc}} = \dfrac{A_{\mathrm{sctlado}}}{\mathit{HVC}_{\mathrm{in}}} = \mathbf{2.872}\,\mathrm{cm^{2}/m}}

\eqitem{Refuerzo de temperatura}{}{A_{\mathrm{sctc}} > A_{\mathrm{sct}}\;\;\Rightarrow\;\;\text{\ok{CUMPLE}}}

\eqitem{Separación máxima 3b}{}{s_{\mathrm{max1ct}} = 3 \cdot b = \mathbf{360}\,\mathrm{cm}}

\begin{longtable}{@{}l p{0.52\linewidth} r l l@{}}
\toprule
\textbf{Dato} & \textbf{Descripción} & \textbf{Valor} & \textbf{Unidad} & \textbf{Ref.}\\
\midrule
$s_{\mathrm{max2ct}}$ & Separación máxima 45 cm & \cellcolor{fondodato}45 & $\mathrm{cm}$ & {\scriptsize } \\
\bottomrule
\end{longtable}

\eqitem{Separación máxima de código}{}{S_{\mathrm{epct}} = \min\left(s_{\mathrm{max1ct}},\ s_{\mathrm{max2ct}}\right) = \mathbf{45}\,\mathrm{cm}}

\eqitem{Separación calculada}{}{s_{\mathrm{ctc}} = \dfrac{A_{\mathrm{basct}}}{A_{\mathrm{sctc}}} = \mathbf{70}\,\mathrm{cm}}

\begin{longtable}{@{}l p{0.52\linewidth} r l l@{}}
\toprule
\textbf{Dato} & \textbf{Descripción} & \textbf{Valor} & \textbf{Unidad} & \textbf{Ref.}\\
\midrule
$s_{\mathrm{ct,e}}$ & Separación escogida & \cellcolor{fondodato}45 & $\mathrm{cm}$ & {\scriptsize } \\
\bottomrule
\end{longtable}

\eqitem{Separación escogida $\le$ máxima}{}{s_{\mathrm{ct,e}} \le S_{\mathrm{epct}}\;\;\Rightarrow\;\;\text{\ok{CUMPLE}}}

\eqitem{Acero proporcionado}{}{A_{\mathrm{sctp}} = \dfrac{A_{\mathrm{basct}}}{s_{\mathrm{ct,e}}} = \mathbf{4.468}\,\mathrm{cm^{2}/m}}

\subsection{Acero de refuerzo superficial [STD A 8.17.2.1]}

\eqitem{Refuerzo superficial requerido}{}{A_{\mathrm{sk}} = 0.012 \cdot \dfrac{\mathrm{in}}{\mathrm{ft}} \cdot \left(d_{\mathrm{c}} - 30\,\mathrm{in}\right) = \mathbf{12.650}\,\mathrm{cm^{2}/m}}

\eqitem{Mitad del refuerzo principal por metro}{}{A_{\mathrm{skT}} = \dfrac{0.5 \cdot A_{\mathrm{sT}}}{\mathrm{m}} = \mathbf{28.510}\,\mathrm{cm^{2}/m}}

\eqitem{2 Ask < 0.5 AsT}{}{2 \cdot A_{\mathrm{sk}} < A_{\mathrm{skT}}\;\;\Rightarrow\;\;\text{\ok{CUMPLE}}}

\eqitem{Acero por lado}{}{A_{\mathrm{slk}} = A_{\mathrm{sk}} \cdot \mathit{HVC}_{\mathrm{in}} = \mathbf{26.565}\,\mathrm{cm^{2}}}

\begin{longtable}{@{}l p{0.52\linewidth} r l l@{}}
\toprule
\textbf{Dato} & \textbf{Descripción} & \textbf{Valor} & \textbf{Unidad} & \textbf{Ref.}\\
\midrule
$D_{\mathrm{refLongk}}$ & Diámetro del refuerzo & \cellcolor{fondodato}1.600 & $\mathrm{cm}$ & {\scriptsize } \\
\bottomrule
\end{longtable}

\eqitem{Área de refuerzo}{}{A_{\mathrm{sklong}} = 0.25 \cdot \pi \cdot {D_{\mathrm{refLongk}}}^{2} = \mathbf{2.011}\,\mathrm{cm^{2}}}

\begin{longtable}{@{}l p{0.52\linewidth} r l l@{}}
\toprule
\textbf{Dato} & \textbf{Descripción} & \textbf{Valor} & \textbf{Unidad} & \textbf{Ref.}\\
\midrule
$n_{\mathrm{var2}}$ & Número de varillas por lado & \cellcolor{fondodato}14 & $$ & {\scriptsize } \\
\bottomrule
\end{longtable}

\eqitem{Área total por lado}{}{A_{\mathrm{sctladok}} = A_{\mathrm{sklong}} \cdot n_{\mathrm{var2}} = \mathbf{28.149}\,\mathrm{cm^{2}}}

\eqitem{Cuantía proporcionada}{}{A_{\mathrm{sctck}} = \dfrac{A_{\mathrm{sctladok}}}{\mathit{HVC}_{\mathrm{in}}} = \mathbf{13.404}\,\mathrm{cm^{2}/m}}

\eqitem{Refuerzo superficial}{}{A_{\mathrm{sctck}} > A_{\mathrm{sk}}\;\;\Rightarrow\;\;\text{\ok{CUMPLE}}}

\eqitem{Separación máxima d/6}{}{s_{\mathrm{sk1}} = \dfrac{d_{\mathrm{c}}}{6} = \mathbf{33.783}\,\mathrm{cm}}

\begin{longtable}{@{}l p{0.52\linewidth} r l l@{}}
\toprule
\textbf{Dato} & \textbf{Descripción} & \textbf{Valor} & \textbf{Unidad} & \textbf{Ref.}\\
\midrule
$s_{\mathrm{sk2}}$ & Separación máxima 30 cm & \cellcolor{fondodato}30 & $\mathrm{cm}$ & {\scriptsize } \\
\bottomrule
\end{longtable}

\eqitem{Separación máxima de código}{}{S_{\mathrm{epssk}} = \min\left(s_{\mathrm{sk1}},\ s_{\mathrm{sk2}}\right) = \mathbf{30}\,\mathrm{cm}}

\eqitem{Separación calculada}{}{s_{\mathrm{ctc2}} = \dfrac{A_{\mathrm{sklong}}}{A_{\mathrm{sctck}}} = \mathbf{15}\,\mathrm{cm}}

\begin{longtable}{@{}l p{0.52\linewidth} r l l@{}}
\toprule
\textbf{Dato} & \textbf{Descripción} & \textbf{Valor} & \textbf{Unidad} & \textbf{Ref.}\\
\midrule
$s_{\mathrm{pe}}$ & Separación escogida & \cellcolor{fondodato}15 & $\mathrm{cm}$ & {\scriptsize } \\
\bottomrule
\end{longtable}

\eqitem{Acero proporcionado}{}{A_{\mathrm{skp}} = \dfrac{A_{\mathrm{sklong}}}{s_{\mathrm{pe}}} = \mathbf{13.404}\,\mathrm{cm^{2}/m}}

\eqitem{Acero superficial proporcionado $\ge$ requerido}{}{A_{\mathrm{sk}} < A_{\mathrm{skp}}\;\;\Rightarrow\;\;\text{\ok{CUMPLE}}}

\clearpage
\section{Fuerza de frenado BR}

\nota{La mayor entre: 1) 25 \% del peso de los ejes del camión o tándem de diseño; 2) 5 \% del camión o tándem de diseño más la carga de carril. [AASHTO LRFD A. 3.6.4]}

\subsection{Datos}

\begin{longtable}{@{}l p{0.52\linewidth} r l l@{}}
\toprule
\textbf{Dato} & \textbf{Descripción} & \textbf{Valor} & \textbf{Unidad} & \textbf{Ref.}\\
\midrule
$L$ & Longitud del puente & \cellcolor{fondodato}30 & $\mathrm{m}$ & {\scriptsize } \\
$N_{carril}$ & Número de carriles & \cellcolor{fondodato}2 & $$ & {\scriptsize } \\
$m$ & Factor de presencia múltiple & \cellcolor{fondodato}1 & $$ & {\scriptsize AASHTO LRFD A. 3.6.1.1.2} \\
\bottomrule
\end{longtable}

\subsection{Camión de diseño HS 20-44}

\begin{longtable}{@{}l p{0.52\linewidth} r l l@{}}
\toprule
\textbf{Dato} & \textbf{Descripción} & \textbf{Valor} & \textbf{Unidad} & \textbf{Ref.}\\
\midrule
$P_{\mathrm{1cd}}$ & Eje delantero & \cellcolor{fondodato}8 & $\mathrm{kip}$ & {\scriptsize } \\
$P_{\mathrm{2cd}}$ & Eje central & \cellcolor{fondodato}32 & $\mathrm{kip}$ & {\scriptsize } \\
$P_{\mathrm{3cd}}$ & Eje trasero & \cellcolor{fondodato}32 & $\mathrm{kip}$ & {\scriptsize } \\
\bottomrule
\end{longtable}

\eqitem{Peso total del camión}{}{P_{\mathrm{cd}} = P_{\mathrm{1cd}} + P_{\mathrm{2cd}} + P_{\mathrm{3cd}} = \mathbf{32\,658.65}\,\mathrm{kgf}}

\subsection{Tándem de diseño}

\begin{longtable}{@{}l p{0.52\linewidth} r l l@{}}
\toprule
\textbf{Dato} & \textbf{Descripción} & \textbf{Valor} & \textbf{Unidad} & \textbf{Ref.}\\
\midrule
$P_{\mathrm{1td}}$ & Eje 1 & \cellcolor{fondodato}25 & $\mathrm{kip}$ & {\scriptsize } \\
$P_{\mathrm{2td}}$ & Eje 2 & \cellcolor{fondodato}25 & $\mathrm{kip}$ & {\scriptsize } \\
\bottomrule
\end{longtable}

\eqitem{Peso total del tándem}{}{P_{\mathrm{td}} = P_{\mathrm{1td}} + P_{\mathrm{2td}} = \mathbf{22\,679.62}\,\mathrm{kgf}}

\subsection{Carga de carril de diseño}

\begin{longtable}{@{}l p{0.52\linewidth} r l l@{}}
\toprule
\textbf{Dato} & \textbf{Descripción} & \textbf{Valor} & \textbf{Unidad} & \textbf{Ref.}\\
\midrule
$\mathit{WLL}$ & Carga de carril & \cellcolor{fondodato}0.6400 & $\mathrm{kip/ft}$ & {\scriptsize } \\
\bottomrule
\end{longtable}

\subsection{Fuerza de frenado}

\eqitem{Opción 1: 25 \% del camión}{}{\mathit{BR}_{\mathrm{1}} = 0.25 \cdot P_{\mathrm{cd}} \cdot N_{carril} \cdot m = \mathbf{16\,329.33}\,\mathrm{kgf}}

\eqitem{Opción 1: 25 \% del tándem}{}{\mathit{BR}_{\mathrm{2}} = 0.25 \cdot P_{\mathrm{td}} \cdot N_{carril} \cdot m = \mathbf{11\,339.81}\,\mathrm{kgf}}

\eqitem{Opción 2: 5 \% (camión + carril)}{}{\mathit{BR}_{\mathrm{3}} = 0.05 \cdot N_{carril} \cdot m \cdot \left(P_{\mathrm{cd}} + \mathit{WLL} \cdot L\right) = \mathbf{6\,123.14}\,\mathrm{kgf}}

\eqitem{Opción 2: 5 \% (tándem + carril)}{}{\mathit{BR}_{\mathrm{4}} = 0.05 \cdot N_{carril} \cdot m \cdot \left(P_{\mathrm{td}} + \mathit{WLL} \cdot L\right) = \mathbf{5\,125.24}\,\mathrm{kgf}}

\eqitem{Fuerza de frenado}{}{\mathit{BR} = \max\left(\mathit{BR}_{\mathrm{1}},\ \mathit{BR}_{\mathrm{2}},\ \mathit{BR}_{\mathrm{3}},\ \mathit{BR}_{\mathrm{4}}\right) = \mathbf{16\,329.33}\,\mathrm{kgf}}

\clearpage
\section{Reacción producida por la carga viva sobre el estribo}

\nota{Reacción máxima en el apoyo A por camión HS 20, tándem y carga de carril (AASHTO LRFD).}

\subsection{Camión de diseño HS 20}

\begin{longtable}{@{}l p{0.52\linewidth} r l l@{}}
\toprule
\textbf{Dato} & \textbf{Descripción} & \textbf{Valor} & \textbf{Unidad} & \textbf{Ref.}\\
\midrule
$L$ & Longitud de la viga & \cellcolor{fondodato}30 & $\mathrm{m}$ & {\scriptsize } \\
$L_{1,cd}$ & Distancia entre ejes P1-P2 & \cellcolor{fondodato}4.270 & $\mathrm{m}$ & {\scriptsize } \\
$L_{2,cd}$ & Distancia entre ejes P2-P3 & \cellcolor{fondodato}4.270 & $\mathrm{m}$ & {\scriptsize } \\
$P_{\mathrm{1cd}}$ & Eje delantero & \cellcolor{fondodato}8 & $\mathrm{kip}$ & {\scriptsize } \\
$P_{\mathrm{2cd}}$ & Eje central & \cellcolor{fondodato}32 & $\mathrm{kip}$ & {\scriptsize } \\
$P_{\mathrm{3cd}}$ & Eje trasero & \cellcolor{fondodato}32 & $\mathrm{kip}$ & {\scriptsize } \\
\bottomrule
\end{longtable}

\nota{Ordenadas de la línea de influencia de la reacción: Ay = 1 - X/L}

\begin{longtable}{@{}l p{0.52\linewidth} r l l@{}}
\toprule
\textbf{Dato} & \textbf{Descripción} & \textbf{Valor} & \textbf{Unidad} & \textbf{Ref.}\\
\midrule
$x_{\mathrm{P3cd}}$ & Posición de P3 (sobre el apoyo) & \cellcolor{fondodato}0 & $\mathrm{m}$ & {\scriptsize } \\
\bottomrule
\end{longtable}

\eqitem{Ordenada P3}{}{o_{\mathrm{P3cd}} = 1 - \dfrac{x_{\mathrm{P3cd}}}{L} = \mathbf{1}}

\eqitem{Posición de P2}{}{x_{\mathrm{P2cd}} = L_{2,cd} = \mathbf{4.270}\,\mathrm{m}}

\eqitem{Ordenada P2}{}{o_{\mathrm{P2cd}} = 1 - \dfrac{x_{\mathrm{P2cd}}}{L} = \mathbf{0.8577}}

\eqitem{Posición de P1}{}{x_{\mathrm{P1cd}} = L_{2,cd} + L_{1,cd} = \mathbf{8.540}\,\mathrm{m}}

\eqitem{Ordenada P1}{}{o_{\mathrm{P1cd}} = 1 - \dfrac{x_{\mathrm{P1cd}}}{L} = \mathbf{0.7153}}

\eqitem{Reacción en el apoyo A (camión)}{}{\mathit{RA}_{\mathrm{cd}} = P_{\mathrm{1cd}} \cdot o_{\mathrm{P1cd}} + P_{\mathrm{2cd}} \cdot o_{\mathrm{P2cd}} + P_{\mathrm{3cd}} \cdot o_{\mathrm{P3cd}} = \mathbf{29\,559.71}\,\mathrm{kgf}}

\subsection{Tándem de diseño}

\begin{longtable}{@{}l p{0.52\linewidth} r l l@{}}
\toprule
\textbf{Dato} & \textbf{Descripción} & \textbf{Valor} & \textbf{Unidad} & \textbf{Ref.}\\
\midrule
$L_{1,td}$ & Distancia entre ejes & \cellcolor{fondodato}1.200 & $\mathrm{m}$ & {\scriptsize } \\
$P_{\mathrm{1td}}$ & Eje 1 & \cellcolor{fondodato}25 & $\mathrm{kip}$ & {\scriptsize } \\
$P_{\mathrm{2td}}$ & Eje 2 & \cellcolor{fondodato}25 & $\mathrm{kip}$ & {\scriptsize } \\
$x_{\mathrm{P2td}}$ & Posición de P2 (sobre el apoyo) & \cellcolor{fondodato}0 & $\mathrm{m}$ & {\scriptsize } \\
\bottomrule
\end{longtable}

\eqitem{Ordenada P2}{}{o_{\mathrm{P2td}} = 1 - \dfrac{x_{\mathrm{P2td}}}{L} = \mathbf{1}}

\eqitem{Posición de P1}{}{x_{\mathrm{P1td}} = L_{1,td} = \mathbf{1.200}\,\mathrm{m}}

\eqitem{Ordenada P1}{}{o_{\mathrm{P1td}} = 1 - \dfrac{x_{\mathrm{P1td}}}{L} = \mathbf{0.9600}}

\eqitem{Reacción en el apoyo A (tándem)}{}{\mathit{RA}_{\mathrm{td}} = P_{\mathrm{1td}} \cdot o_{\mathrm{P1td}} + P_{\mathrm{2td}} \cdot o_{\mathrm{P2td}} = \mathbf{22\,226.03}\,\mathrm{kgf}}

\subsection{Carga de carril de diseño}

\begin{longtable}{@{}l p{0.52\linewidth} r l l@{}}
\toprule
\textbf{Dato} & \textbf{Descripción} & \textbf{Valor} & \textbf{Unidad} & \textbf{Ref.}\\
\midrule
$\mathit{WLL}$ & Carga de carril & \cellcolor{fondodato}0.6400 & $\mathrm{kip/ft}$ & {\scriptsize } \\
\bottomrule
\end{longtable}

\eqitem{Reacción por carril}{}{\mathit{RA}_{\mathrm{cc}} = \mathit{WLL} \cdot L \cdot 0.5 = \mathbf{14\,286.37}\,\mathrm{kgf}}

\subsection{Combinación}

\eqitem{Reacción máxima (camión o tándem)}{}{\mathit{RA}_{\mathrm{max}} = \max\left(\mathit{RA}_{\mathrm{cd}},\ \mathit{RA}_{\mathrm{td}}\right) = \mathbf{29\,559.71}\,\mathrm{kgf}}

\begin{longtable}{@{}l p{0.52\linewidth} r l l@{}}
\toprule
\textbf{Dato} & \textbf{Descripción} & \textbf{Valor} & \textbf{Unidad} & \textbf{Ref.}\\
\midrule
$IM$ & Incremento por carga dinámica & \cellcolor{fondodato}33 & $\%$ & {\scriptsize } \\
\bottomrule
\end{longtable}

\eqitem{Combinación máxima por carril}{}{\mathit{RA}_{\mathrm{1carril}} = \left(1 + IM\right) \cdot \mathit{RA}_{\mathrm{max}} + \mathit{RA}_{\mathrm{cc}} = \mathbf{53\,600.78}\,\mathrm{kgf}}

\begin{longtable}{@{}l p{0.52\linewidth} r l l@{}}
\toprule
\textbf{Dato} & \textbf{Descripción} & \textbf{Valor} & \textbf{Unidad} & \textbf{Ref.}\\
\midrule
$N_{carriles}$ & Número de carriles & \cellcolor{fondodato}2 & $$ & {\scriptsize } \\
$m$ & Factor de presencia múltiple & \cellcolor{fondodato}1 & $$ & {\scriptsize AASHTO LRFD A. 3.6.1.1.2} \\
\bottomrule
\end{longtable}

\eqitem{Reacción total sobre el estribo}{}{\mathit{RA}_{\mathrm{2carril}} = N_{carriles} \cdot m \cdot \mathit{RA}_{\mathrm{1carril}} = \mathbf{107\,201.57}\,\mathrm{kgf}}

\clearpage
\section{Estribo sobre cimentación superficial - AASHTO LRFD}

\nota{Diseño de estribo en voladizo de un puente de tramo simple sobre suelo no rocoso. AASHTO LRFD 2017.}

\subsection{Datos}

\begin{longtable}{@{}l p{0.52\linewidth} r l l@{}}
\toprule
\textbf{Dato} & \textbf{Descripción} & \textbf{Valor} & \textbf{Unidad} & \textbf{Ref.}\\
\midrule
$f'_{c,est}$ & f'c del estribo & \cellcolor{fondodato}21 & $\mathrm{MPa}$ & {\scriptsize } \\
$f'_{c,sup}$ & f'c de la superestructura & \cellcolor{fondodato}28 & $\mathrm{MPa}$ & {\scriptsize } \\
$f_{\mathrm{y}}$ & Fluencia del acero de refuerzo & \cellcolor{fondodato}420 & $\mathrm{MPa}$ & {\scriptsize } \\
$\gamma_{H}$ & Peso específico del hormigón & \cellcolor{fondodato}2\,400 & $\mathrm{kgf/m^{3}}$ & {\scriptsize } \\
\bottomrule
\end{longtable}

\nota{Información del suelo: suelo no rocoso.}

\begin{longtable}{@{}l p{0.52\linewidth} r l l@{}}
\toprule
\textbf{Dato} & \textbf{Descripción} & \textbf{Valor} & \textbf{Unidad} & \textbf{Ref.}\\
\midrule
$\gamma_s$ & Peso específico del suelo & \cellcolor{fondodato}2\,000 & $\mathrm{kgf/m^{3}}$ & {\scriptsize } \\
$q_{\mathrm{n}}$ & Capacidad de carga nominal del suelo & \cellcolor{fondodato}1.200 & $\mathrm{MPa}$ & {\scriptsize } \\
$\phi'_f$ & Ángulo efectivo de fricción interna del suelo & \cellcolor{fondodato}30 & $^{\circ}$ & {\scriptsize } \\
$\mathit{PGA}$ & Coeficiente de aceleración pico del terreno & \cellcolor{fondodato}0.2000 & $$ & {\scriptsize } \\
\bottomrule
\end{longtable}

\subsection{Cargas provenientes de la superestructura (por metro de estribo)}

\begin{longtable}{@{}l p{0.52\linewidth} r l l@{}}
\toprule
\textbf{Dato} & \textbf{Descripción} & \textbf{Valor} & \textbf{Unidad} & \textbf{Ref.}\\
\midrule
$\mathit{DC}$ & Carga muerta de componentes & \cellcolor{fondodato}11\,000 & $\mathrm{kgf/m}$ & {\scriptsize } \\
$\mathit{DW}$ & Carga muerta de superficie de rodadura & \cellcolor{fondodato}1\,500 & $\mathrm{kgf/m}$ & {\scriptsize } \\
$LL_{LL+IM}$ & Carga viva + impacto & \cellcolor{fondodato}9\,300 & $\mathrm{kgf/m}$ & {\scriptsize } \\
$\mathit{BR}$ & Fuerza de frenado & \cellcolor{fondodato}1\,570 & $\mathrm{kgf/m}$ & {\scriptsize } \\
\bottomrule
\end{longtable}

\subsection{Cargas y momentos debido al peso de la subestructura}

\begin{longtable}{@{}l p{0.52\linewidth} r l l@{}}
\toprule
\textbf{Dato} & \textbf{Descripción} & \textbf{Valor} & \textbf{Unidad} & \textbf{Ref.}\\
\midrule
$L_{\mathrm{est}}$ & Longitud del estribo & \cellcolor{fondodato}10.500 & $\mathrm{m}$ & {\scriptsize } \\
$l_{1}$ & Dimensión l1 del estribo & \cellcolor{fondodato}1.100 & $\mathrm{m}$ & {\scriptsize } \\
$l_{2}$ & Dimensión l2 del estribo & \cellcolor{fondodato}3.400 & $\mathrm{m}$ & {\scriptsize } \\
$l_{3}$ & Dimensión l3 del estribo & \cellcolor{fondodato}1 & $\mathrm{m}$ & {\scriptsize } \\
$l_{4}$ & Dimensión l4 del estribo & \cellcolor{fondodato}1.500 & $\mathrm{m}$ & {\scriptsize } \\
$l_{5}$ & Dimensión l5 del estribo & \cellcolor{fondodato}0.3000 & $\mathrm{m}$ & {\scriptsize } \\
$l_{6}$ & Dimensión l6 del estribo & \cellcolor{fondodato}0.6000 & $\mathrm{m}$ & {\scriptsize } \\
$l_{8}$ & Dimensión l8 del estribo & \cellcolor{fondodato}1.200 & $\mathrm{m}$ & {\scriptsize } \\
$l_{9}$ & Dimensión l9 del estribo & \cellcolor{fondodato}2.800 & $\mathrm{m}$ & {\scriptsize } \\
$l_{10}$ & Dimensión l10 del estribo & \cellcolor{fondodato}1.100 & $\mathrm{m}$ & {\scriptsize } \\
\bottomrule
\end{longtable}

\eqitem{Altura de la pantalla bajo el espaldar}{}{l_{7} = l_{2} + l_{3} = \mathbf{4.400}\,\mathrm{m}}

\eqitem{Altura del relleno}{}{H = l_{1} + l_{2} + l_{3} + l_{4} = \mathbf{7}\,\mathrm{m}}

\eqitem{Altura del cimiento}{}{D = l_{1} = \mathbf{1.100}\,\mathrm{m}}

\eqitem{Base del estribo}{}{B = l_{8} + l_{5} + l_{6} + l_{9} = \mathbf{4.900}\,\mathrm{m}}

\nota{DC por metro: zapata, pantalla y espaldar se calculan por metro de ancho; topes y muros de ala se reparten en la longitud del estribo Lest. Brazos medidos desde la punta de la zapata.}

\begin{center}
{\small\textbf{Peso propio del estribo}}\\[3pt]
\adjustbox{max width=\linewidth}{\begin{tabular}{ccccccccc}
\toprule
\begin{tabular}[b]{@{}c@{}}\textbf{Sección}\end{tabular} & \begin{tabular}[b]{@{}c@{}}\textbf{Base}\\ {\scriptsize $\mathrm{m}$}\end{tabular} & \begin{tabular}[b]{@{}c@{}}\textbf{Altura}\\ {\scriptsize $\mathrm{m}$}\end{tabular} & \begin{tabular}[b]{@{}c@{}}\textbf{Ancho}\\ {\scriptsize $\mathrm{m}$}\end{tabular} & \begin{tabular}[b]{@{}c@{}}\textbf{Cantidad}\end{tabular} & \begin{tabular}[b]{@{}c@{}}\textbf{Volumen}\\ {\scriptsize $\mathrm{m^{3}}$}\end{tabular} & \begin{tabular}[b]{@{}c@{}}\textbf{DC}\\ {\scriptsize $\mathrm{tonf/m}$}\end{tabular} & \begin{tabular}[b]{@{}c@{}}\textbf{x}\\ {\scriptsize $\mathrm{m}$}\end{tabular} & \begin{tabular}[b]{@{}c@{}}\textbf{Mx}\\ {\scriptsize $\mathrm{tonf\cdot m/m}$}\end{tabular} \\
\midrule
zapata & $4.900$ & $1.100$ & \cellcolor{fondodato}$1$ & \cellcolor{fondodato}$1$ & $5.390$ & $12.936$ & $2.450$ & $31.693$ \\
pantalla & $0.9000$ & $4.400$ & \cellcolor{fondodato}$1$ & \cellcolor{fondodato}$1$ & $3.960$ & $9.504$ & $1.650$ & $15.682$ \\
espaldar & $0.3000$ & $1.500$ & \cellcolor{fondodato}$1$ & \cellcolor{fondodato}$1$ & $0.4500$ & $1.080$ & $1.950$ & $2.106$ \\
topes & $0.6000$ & $1.100$ & \cellcolor{fondodato}$0.8000$ & \cellcolor{fondodato}$2$ & $1.056$ & $0.2414$ & $1.500$ & $0.3621$ \\
muros ala & $2.800$ & $5.900$ & \cellcolor{fondodato}$0.5000$ & \cellcolor{fondodato}$2$ & $16.520$ & $3.776$ & $3.500$ & $13.216$ \\
\bottomrule
\end{tabular}}
\end{center}

\eqitem{Peso propio del estribo}{}{DC_{est} = \sum DC_i = \mathbf{27.537}\,\mathrm{tonf/m}}

\eqitem{Momento del peso propio}{}{Mx_{DC,est} = \sum DC_i\,x_i = \mathbf{63.059}\,\mathrm{tonf\cdot m/m}}

\eqitem{Brazo del peso propio}{}{x_{DC,est} = \dfrac{Mx_{DC,est}}{DC_{est}} = \mathbf{2.290}\,\mathrm{m}}

\subsubsection{Carga y momento debidos al peso propio del relleno}

\eqitem{Base del relleno}{}{Base_{EV} = l_{9} = \mathbf{2.800}\,\mathrm{m}}

\eqitem{Altura del relleno}{}{Alto_{EV} = l_{2} + l_{3} + l_{4} = \mathbf{5.900}\,\mathrm{m}}

\eqitem{Peso del relleno}{}{\mathit{EV} = Base_{EV} \cdot Alto_{EV} \cdot \gamma_s = \mathbf{33.040}\,\mathrm{tonf/m}}

\eqitem{Brazo del relleno}{}{x_{relleno} = B - \dfrac{l_{9}}{2} = \mathbf{3.500}\,\mathrm{m}}

\eqitem{Momento del relleno}{}{Mx_{EV} = \mathit{EV} \cdot x_{relleno} = \mathbf{115.64}\,\mathrm{tonf\cdot m/m}}

\subsection{Sobrecarga por carga viva sobre el relleno (LS) [A. 3.11.5.3, A. 3.11.6.4]}

\begin{longtable}{@{}l p{0.52\linewidth} r l l@{}}
\toprule
\textbf{Dato} & \textbf{Descripción} & \textbf{Valor} & \textbf{Unidad} & \textbf{Ref.}\\
\midrule
$\delta$ & Ángulo de fricción relleno-muro & \cellcolor{fondodato}0 & $^{\circ}$ & {\scriptsize } \\
$\beta$ & Ángulo del relleno con la horizontal & \cellcolor{fondodato}0 & $^{\circ}$ & {\scriptsize } \\
$\theta$ & Ángulo de la cara posterior del muro con la horizontal & \cellcolor{fondodato}90 & $^{\circ}$ & {\scriptsize } \\
\bottomrule
\end{longtable}

\eqitem{Coeficiente de empuje activo (Coulomb)}{A. 3.11.5.3-1}{k_a = \dfrac{{\left(\sin\left(\theta + \phi'_f\right)\right)}^{2}}{{\left(1 + \sqrt{\dfrac{\sin\left(\phi'_f + \delta\right) \cdot \sin\left(\phi'_f - \beta\right)}{\sin\left(\theta - \delta\right) \cdot \sin\left(\theta + \beta\right)}}\right)}^{2} \cdot {\left(\sin\left(\theta\right)\right)}^{2} \cdot \sin\left(\theta - \delta\right)} = \mathbf{0.3333}}

\begin{longtable}{@{}l p{0.52\linewidth} r l l@{}}
\toprule
\textbf{Dato} & \textbf{Descripción} & \textbf{Valor} & \textbf{Unidad} & \textbf{Ref.}\\
\midrule
$h_{\mathrm{eq}}$ & Altura equivalente de suelo por sobrecarga & \cellcolor{fondodato}600 & $\mathrm{mm}$ & {\scriptsize T. 3.11.6.4-1} \\
\bottomrule
\end{longtable}

\eqitem{Presión por sobrecarga}{}{\Delta p_{LS} = k_a \cdot \gamma_s \cdot h_{\mathrm{eq}} = \mathbf{0.4000}\,\mathrm{tonf/m^{2}}}

\eqitem{Componente vertical de la sobrecarga}{}{\mathit{LS}_{\mathrm{y}} = \Delta p_{LS} \cdot l_{9} = \mathbf{1.120}\,\mathrm{tonf/m}}

\eqitem{Momento estabilizante de la sobrecarga}{}{Mx_{LSy} = \mathit{LS}_{\mathrm{y}} \cdot x_{relleno} = \mathbf{3.920}\,\mathrm{tonf\cdot m/m}}

\subsubsection{Acciones verticales por metro de estribo}

\eqitem{Brazo de DC y DW de la superestructura}{}{x_{DC} = l_{8} + \dfrac{l_{6}}{2} = \mathbf{1.500}\,\mathrm{m}}

\eqitem{Momento de DC superestructura}{}{Mx_{DC} = \mathit{DC} \cdot x_{DC} = \mathbf{16.500}\,\mathrm{tonf\cdot m/m}}

\eqitem{Momento de DW}{}{Mx_{DW} = \mathit{DW} \cdot x_{DC} = \mathbf{2.250}\,\mathrm{tonf\cdot m/m}}

\eqitem{Momento de LL+IM}{}{Mx_{LL+IM} = LL_{LL+IM} \cdot x_{DC} = \mathbf{13.950}\,\mathrm{tonf\cdot m/m}}

\eqitem{Suma de fuerzas verticales}{}{\Sigma V = DC_{est} + \mathit{DC} + \mathit{DW} + LL_{LL+IM} + \mathit{EV} + \mathit{LS}_{\mathrm{y}} = \mathbf{83.497}\,\mathrm{tonf/m}}

\eqitem{Suma de momentos estabilizantes}{}{\Sigma MV = Mx_{DC,est} + Mx_{DC} + Mx_{DW} + Mx_{LL+IM} + Mx_{EV} + Mx_{LSy} = \mathbf{215.32}\,\mathrm{tonf\cdot m/m}}

\subsection{Empuje horizontal del suelo (EH) [A. 3.11.5.3]}

\eqitem{Coeficiente de presión activa (Rankine)}{}{K_A = \dfrac{1 - \sin\left(\phi'_f\right)}{1 + \sin\left(\phi'_f\right)} = \mathbf{0.3333}}

\eqitem{Presión activa en la base}{}{\Delta p_{EH} = K_A \cdot \gamma_s \cdot H = \mathbf{4.667}\,\mathrm{tonf/m^{2}}}

\eqitem{Empuje activo estático}{}{\mathit{EH} = \dfrac{\Delta p_{EH} \cdot H}{2} = \mathbf{16.333}\,\mathrm{tonf/m}}

\eqitem{Brazo}{}{y_{EH} = \dfrac{H}{3} = \mathbf{2.333}\,\mathrm{m}}

\eqitem{Momento de volteo}{}{My_{EH} = \mathit{EH} \cdot y_{EH} = \mathbf{38.111}\,\mathrm{tonf\cdot m/m}}

\eqitem{Empuje por sobrecarga viva}{}{\mathit{LS}_{\mathrm{x}} = \Delta p_{LS} \cdot H = \mathbf{2.800}\,\mathrm{tonf/m}}

\eqitem{Brazo}{}{y_{LSx} = \dfrac{H}{2} = \mathbf{3.500}\,\mathrm{m}}

\eqitem{Momento de volteo}{}{My_{LSx} = \mathit{LS}_{\mathrm{x}} \cdot y_{LSx} = \mathbf{9.800}\,\mathrm{tonf\cdot m/m}}

\subsection{Fuerzas sísmicas [A. 11.6.5]}

\nota{Efecto combinado de PAE (empuje dinámico del relleno, Ec. 11.6.5.3-2) y PIR (fuerza inercial del muro, Ec. 11.6.5.1-1). kh = 0.5 kh0 [A. 11.6.5.2.2].}

\begin{longtable}{@{}l p{0.52\linewidth} r l l@{}}
\toprule
\textbf{Dato} & \textbf{Descripción} & \textbf{Valor} & \textbf{Unidad} & \textbf{Ref.}\\
\midrule
$F_{pga}$ & Coeficiente de sitio para período corto & \cellcolor{fondodato}1.200 & $$ & {\scriptsize } \\
\bottomrule
\end{longtable}

\eqitem{Coeficiente de aceleración pico As}{}{A_s = F_{pga} \cdot \mathit{PGA} = \mathbf{0.2400}}

\eqitem{kh0}{}{k_{h0} = A_s = \mathbf{0.2400}}

\eqitem{Coeficiente sísmico horizontal}{}{k_h = 0.5 \cdot k_{h0} = \mathbf{0.1200}}

\begin{longtable}{@{}l p{0.52\linewidth} r l l@{}}
\toprule
\textbf{Dato} & \textbf{Descripción} & \textbf{Valor} & \textbf{Unidad} & \textbf{Ref.}\\
\midrule
$k_v$ & Coeficiente sísmico vertical & \cellcolor{fondodato}0 & $$ & {\scriptsize A. 11.6.5.2.1} \\
$\delta$ & Fricción suelo-muro (sismo) & \cellcolor{fondodato}0 & $^{\circ}$ & {\scriptsize } \\
$\beta$ & Inclinación del muro respecto a la vertical & \cellcolor{fondodato}0 & $^{\circ}$ & {\scriptsize } \\
$i$ & Inclinación del relleno & \cellcolor{fondodato}0 & $^{\circ}$ & {\scriptsize } \\
\bottomrule
\end{longtable}

\eqitem{Ángulo $\theta$MO}{}{\theta_{MO} = \arctan\left(\dfrac{k_h}{1 - k_v}\right) = \mathbf{6.843}^{\circ}}

\eqitem{$\phi$ $\ge$ i + $\theta$MO (Ec. 11.6.5.3-1)}{}{\phi'_f \ge i + \theta_{MO}\;\;\Rightarrow\;\;\text{\ok{CUMPLE}}}

\eqitem{Coeficiente de empuje sísmico (Mononobe-Okabe)}{}{K_{AE} = \dfrac{{\left(\cos\left(\phi'_f - \theta_{MO} - \beta\right)\right)}^{2}}{\cos\left(\theta_{MO}\right) \cdot {\left(\cos\left(\beta\right)\right)}^{2} \cdot \cos\left(\delta + \theta_{MO} + \beta\right)} \cdot {\left(1 + \sqrt{\dfrac{\sin\left(\phi'_f + \delta\right) \cdot \sin\left(\phi'_f - \theta_{MO} - i\right)}{\cos\left(\delta + \theta_{MO} + \beta\right) \cdot \cos\left(i - \beta\right)}}\right)}^{-2} = \mathbf{0.4107}}

\eqitem{Fuerza sísmica PAE}{}{\mathit{PAE} = 0.5 \cdot K_{AE} \cdot {H}^{2} \cdot \gamma_s = \mathbf{20.123}\,\mathrm{tonf/m}}

\eqitem{Incremento dinámico $\Delta$PAE}{}{\Delta P_{AE} = \mathit{PAE} - \mathit{EH} = \mathbf{3.790}\,\mathrm{tonf/m}}

\eqitem{Brazo de $\Delta$PAE}{}{Y_{\Delta PAE} = 0.5 \cdot H = \mathbf{3.500}\,\mathrm{m}}

\eqitem{Momento de $\Delta$PAE}{}{M_{\Delta PAE} = \Delta P_{AE} \cdot Y_{\Delta PAE} = \mathbf{13.265}\,\mathrm{tonf\cdot m/m}}

\eqitem{Peso del suelo detrás del muro}{}{W_s = \mathit{EV} = \mathbf{33.040}\,\mathrm{tonf/m}}

\eqitem{Peso del estribo}{}{W_w = DC_{est} = \mathbf{27.537}\,\mathrm{tonf/m}}

\eqitem{Fuerza inercial PIR}{}{P_{IR} = k_h \cdot \left(W_w + W_s\right) = \mathbf{7.269}\,\mathrm{tonf/m}}

\eqitem{Empuje activo sísmico}{}{Empuje_{act.sis} = \max\left(\Delta P_{AE},\ \mathit{EH}\right) = \mathbf{16.333}\,\mathrm{tonf/m}}

\eqitem{Combinación 1}{}{P_{seis1} = \Delta P_{AE} + 0.5 \cdot P_{IR} = \mathbf{7.425}\,\mathrm{tonf/m}}

\eqitem{Combinación 2}{}{P_{seis2} = 0.5 \cdot Empuje_{act.sis} + P_{IR} = \mathbf{15.436}\,\mathrm{tonf/m}}

\eqitem{Fuerza sísmica de diseño}{}{P_{seis} = \max\left(P_{seis1},\ P_{seis2}\right) = \mathbf{15.436}\,\mathrm{tonf/m}}

\eqitem{Punto de aplicación}{}{Y_{Pseis} = 0.4 \cdot H = \mathbf{2.800}\,\mathrm{m}}

\eqitem{Momento respecto al punto de volteo}{}{MP_{seis} = P_{seis} \cdot Y_{Pseis} = \mathbf{43.221}\,\mathrm{tonf\cdot m/m}}

\subsubsection{Fuerza sísmica proveniente de la superestructura, Hu [A. 14.6.3, C. 11.6.5.1]}

\begin{longtable}{@{}l p{0.52\linewidth} r l l@{}}
\toprule
\textbf{Dato} & \textbf{Descripción} & \textbf{Valor} & \textbf{Unidad} & \textbf{Ref.}\\
\midrule
$\mu$ & Coeficiente de fricción & \cellcolor{fondodato}0.2000 & $$ & {\scriptsize } \\
$\gamma_{DC}$ & Factor DC (Resistencia I) & \cellcolor{fondodato}1.250 & $$ & {\scriptsize } \\
$\gamma_{DW}$ & Factor DW (Resistencia I) & \cellcolor{fondodato}1.500 & $$ & {\scriptsize } \\
$\gamma_{EQ}$ & Factor LL (Evento Extremo I) & \cellcolor{fondodato}0.5000 & $$ & {\scriptsize } \\
\bottomrule
\end{longtable}

\eqitem{Fuerza de fricción por desplazamiento}{}{H_{\mathrm{u}} = \mu \cdot \left(\gamma_{DC} \cdot \mathit{DC} + \gamma_{DW} \cdot \mathit{DW} + \gamma_{EQ} \cdot LL_{LL+IM}\right) = \mathbf{4.130}\,\mathrm{tonf/m}}

\eqitem{Brazo de Hu}{}{Y_{Hu} = l_{1} + l_{7} = \mathbf{5.500}\,\mathrm{m}}

\eqitem{Momento de Hu}{}{MH_u = H_{\mathrm{u}} \cdot Y_{Hu} = \mathbf{22.715}\,\mathrm{tonf\cdot m/m}}

\eqitem{Suma de acciones sísmicas}{}{\Sigma H_{sis} = P_{seis} + H_{\mathrm{u}} = \mathbf{19.566}\,\mathrm{tonf/m}}

\eqitem{Suma de momentos sísmicos}{}{\Sigma MH_{sis} = MP_{seis} + MH_u = \mathbf{65.936}\,\mathrm{tonf\cdot m/m}}

\subsection{Verificación de la estabilidad y esfuerzos sobre el suelo [T. 3.4.1-1, T. 3.4.1-2]}

\eqitem{Brazo del frenado}{}{Y_{BR} = H + 1.8\,\mathrm{m} = \mathbf{8.800}\,\mathrm{m}}

\eqitem{Momento del frenado}{}{M_{BR} = \mathit{BR} \cdot Y_{BR} = \mathbf{13.816}\,\mathrm{tonf\cdot m/m}}

\eqitem{Suma de fuerzas horizontales}{}{\Sigma H = \mathit{EH} + P_{seis} + H_{\mathrm{u}} + \mathit{BR} + \mathit{LS}_{\mathrm{x}} = \mathbf{40.269}\,\mathrm{tonf/m}}

\eqitem{Suma de momentos de volteo}{}{\Sigma MH = My_{EH} + MP_{seis} + MH_u + M_{BR} + My_{LSx} = \mathbf{127.66}\,\mathrm{tonf\cdot m/m}}

\begin{center}
{\small\textbf{Factores de carga - fuerzas verticales y momentos estabilizantes}}\\[3pt]
\adjustbox{max width=\linewidth}{\begin{tabular}{ccccccc}
\toprule
\begin{tabular}[b]{@{}c@{}}\textbf{Estado límite}\end{tabular} & \begin{tabular}[b]{@{}c@{}}\textbf{DC super.}\end{tabular} & \begin{tabular}[b]{@{}c@{}}\textbf{DC est.}\end{tabular} & \begin{tabular}[b]{@{}c@{}}\textbf{DW}\end{tabular} & \begin{tabular}[b]{@{}c@{}}\textbf{EV}\end{tabular} & \begin{tabular}[b]{@{}c@{}}\textbf{LL+IM}\end{tabular} & \begin{tabular}[b]{@{}c@{}}\textbf{LSy}\end{tabular} \\
\midrule
Resistencia I - Máx. & \cellcolor{fondodato}$1.250$ & \cellcolor{fondodato}$1.250$ & \cellcolor{fondodato}$1.500$ & \cellcolor{fondodato}$1.350$ & \cellcolor{fondodato}$1.750$ & \cellcolor{fondodato}$1.750$ \\
Resistencia I - Mín. & \cellcolor{fondodato}$0.9000$ & \cellcolor{fondodato}$0.9000$ & \cellcolor{fondodato}$0.6500$ & \cellcolor{fondodato}$1$ & \cellcolor{fondodato}$0$ & \cellcolor{fondodato}$0$ \\
Evento Extremo I - Máx. & \cellcolor{fondodato}$1.250$ & \cellcolor{fondodato}$1.250$ & \cellcolor{fondodato}$1.500$ & \cellcolor{fondodato}$1.350$ & \cellcolor{fondodato}$0.5000$ & \cellcolor{fondodato}$0.5000$ \\
Evento Extremo I - Mín. & \cellcolor{fondodato}$0.9000$ & \cellcolor{fondodato}$0.9000$ & \cellcolor{fondodato}$0.6500$ & \cellcolor{fondodato}$1$ & \cellcolor{fondodato}$0$ & \cellcolor{fondodato}$0$ \\
Servicio I & \cellcolor{fondodato}$1$ & \cellcolor{fondodato}$1$ & \cellcolor{fondodato}$1$ & \cellcolor{fondodato}$1$ & \cellcolor{fondodato}$1$ & \cellcolor{fondodato}$1$ \\
\bottomrule
\end{tabular}}
\end{center}

\begin{center}
{\small\textbf{Factores de carga - fuerzas horizontales y momentos de volteo}}\\[3pt]
\adjustbox{max width=\linewidth}{\begin{tabular}{cccccc}
\toprule
\begin{tabular}[b]{@{}c@{}}\textbf{Estado límite}\end{tabular} & \begin{tabular}[b]{@{}c@{}}\textbf{EH}\end{tabular} & \begin{tabular}[b]{@{}c@{}}\textbf{LSx}\end{tabular} & \begin{tabular}[b]{@{}c@{}}\textbf{BR}\end{tabular} & \begin{tabular}[b]{@{}c@{}}\textbf{Pseis}\end{tabular} & \begin{tabular}[b]{@{}c@{}}\textbf{Hbu}\end{tabular} \\
\midrule
Resistencia I - Máx. & \cellcolor{fondodato}$1.500$ & \cellcolor{fondodato}$1.750$ & \cellcolor{fondodato}$1.750$ & \cellcolor{fondodato}$0$ & \cellcolor{fondodato}$0$ \\
Resistencia I - Mín. & \cellcolor{fondodato}$0.9000$ & \cellcolor{fondodato}$0$ & \cellcolor{fondodato}$0$ & \cellcolor{fondodato}$0$ & \cellcolor{fondodato}$0$ \\
Evento Extremo I - Máx. & \cellcolor{fondodato}$1.500$ & \cellcolor{fondodato}$0.5000$ & \cellcolor{fondodato}$0.5000$ & \cellcolor{fondodato}$1$ & \cellcolor{fondodato}$1$ \\
Evento Extremo I - Mín. & \cellcolor{fondodato}$0.9000$ & \cellcolor{fondodato}$0$ & \cellcolor{fondodato}$0$ & \cellcolor{fondodato}$1$ & \cellcolor{fondodato}$1$ \\
Servicio I & \cellcolor{fondodato}$1$ & \cellcolor{fondodato}$1$ & \cellcolor{fondodato}$1$ & \cellcolor{fondodato}$0$ & \cellcolor{fondodato}$0$ \\
\bottomrule
\end{tabular}}
\end{center}

\begin{center}
{\small\textbf{Combinaciones de carga (por metro de estribo)}}\\[3pt]
\adjustbox{max width=\linewidth}{\begin{tabular}{ccccc}
\toprule
\begin{tabular}[b]{@{}c@{}}\textbf{Estado límite}\end{tabular} & \begin{tabular}[b]{@{}c@{}}\textbf{$\Sigma$V}\\ {\scriptsize $\mathrm{tonf/m}$}\end{tabular} & \begin{tabular}[b]{@{}c@{}}\textbf{$\Sigma$MV}\\ {\scriptsize $\mathrm{tonf\cdot m/m}$}\end{tabular} & \begin{tabular}[b]{@{}c@{}}\textbf{$\Sigma$H}\\ {\scriptsize $\mathrm{tonf/m}$}\end{tabular} & \begin{tabular}[b]{@{}c@{}}\textbf{$\Sigma$MH}\\ {\scriptsize $\mathrm{tonf\cdot m/m}$}\end{tabular} \\
\midrule
Resistencia I - Máx. & $113.26$ & $290.21$ & $32.148$ & $98.495$ \\
Resistencia I - Mín. & $68.699$ & $188.71$ & $14.700$ & $34.300$ \\
Evento Extremo I - Máx. & $100.24$ & $267.87$ & $46.251$ & $134.91$ \\
Evento Extremo I - Mín. & $68.699$ & $188.71$ & $34.266$ & $100.24$ \\
Servicio I & $83.497$ & $215.32$ & $20.703$ & $61.727$ \\
\bottomrule
\end{tabular}}
\end{center}

\subsubsection{Verificación al volcamiento [A. 11.6.3, A. 11.6.5, A. 10.5.2]}

\nota{Xo = (MV - MH)/V ;  e = B/2 - Xo ;  se exige e < emax.}

\eqitem{Excentricidad máxima - Resistencia I (A. 11.6.3.3)}{}{e_{max,R} = \dfrac{B}{3} = \mathbf{1.633}\,\mathrm{m}}

\begin{longtable}{@{}l p{0.52\linewidth} r l l@{}}
\toprule
\textbf{Dato} & \textbf{Descripción} & \textbf{Valor} & \textbf{Unidad} & \textbf{Ref.}\\
\midrule
$\gamma_{EQ}$ & Factor de carga viva simultánea con sismo & \cellcolor{fondodato}0.5000 & $$ & {\scriptsize } \\
\bottomrule
\end{longtable}

\eqitem{Excentricidad máxima - Evento Extremo I (A. 11.6.5)}{}{e_{max,EV} = \dfrac{11 \cdot B}{30} = \mathbf{1.797}\,\mathrm{m}}

\begin{center}
{\small\textbf{Verificación al volcamiento}}\\[3pt]
\adjustbox{max width=\linewidth}{\begin{tabular}{cccccccc}
\toprule
\begin{tabular}[b]{@{}c@{}}\textbf{Estado}\end{tabular} & \begin{tabular}[b]{@{}c@{}}\textbf{$\Sigma$V}\\ {\scriptsize $\mathrm{tonf/m}$}\end{tabular} & \begin{tabular}[b]{@{}c@{}}\textbf{MV}\\ {\scriptsize $\mathrm{tonf\cdot m/m}$}\end{tabular} & \begin{tabular}[b]{@{}c@{}}\textbf{MH}\\ {\scriptsize $\mathrm{tonf\cdot m/m}$}\end{tabular} & \begin{tabular}[b]{@{}c@{}}\textbf{Xo}\\ {\scriptsize $\mathrm{m}$}\end{tabular} & \begin{tabular}[b]{@{}c@{}}\textbf{e}\\ {\scriptsize $\mathrm{m}$}\end{tabular} & \begin{tabular}[b]{@{}c@{}}\textbf{emax}\\ {\scriptsize $\mathrm{m}$}\end{tabular} & \begin{tabular}[b]{@{}c@{}}\textbf{Cumple}\end{tabular} \\
\midrule
Resistencia I - Máx. & $113.26$ & $290.21$ & $98.495$ & $1.693$ & $0.7573$ & $1.633$ & \ok{OK} \\
Resistencia I - Mín. & $68.699$ & $188.71$ & $34.300$ & $2.248$ & $0.2024$ & $1.633$ & \ok{OK} \\
Evento Extremo I - Máx. & $100.24$ & $267.87$ & $134.91$ & $1.326$ & $1.124$ & $1.797$ & \ok{OK} \\
Evento Extremo I - Mín. & $68.699$ & $188.71$ & $100.24$ & $1.288$ & $1.162$ & $1.797$ & \ok{OK} \\
\bottomrule
\end{tabular}}
\end{center}

\subsubsection{Esfuerzos sobre el suelo [A. 10.6.3, A. 10.5.5.2.2]}

\begin{longtable}{@{}l p{0.52\linewidth} r l l@{}}
\toprule
\textbf{Dato} & \textbf{Descripción} & \textbf{Valor} & \textbf{Unidad} & \textbf{Ref.}\\
\midrule
$\varphi_{b,r}$ & Factor de resistencia - Resistencia I & \cellcolor{fondodato}0.4500 & $$ & {\scriptsize } \\
\bottomrule
\end{longtable}

\eqitem{Capacidad mayorada - Resistencia I}{}{q_{\mathrm{R}} = \varphi_{b,r} \cdot q_{\mathrm{n}} = \mathbf{55.065}\,\mathrm{tonf/m^{2}}}

\begin{longtable}{@{}l p{0.52\linewidth} r l l@{}}
\toprule
\textbf{Dato} & \textbf{Descripción} & \textbf{Valor} & \textbf{Unidad} & \textbf{Ref.}\\
\midrule
$\varphi_{b,EV}$ & Factor de resistencia - Evento Extremo I & \cellcolor{fondodato}0.8000 & $$ & {\scriptsize } \\
\bottomrule
\end{longtable}

\eqitem{Capacidad mayorada - Evento Extremo I}{}{q_{R,EV} = \varphi_{b,EV} \cdot q_{\mathrm{n}} = \mathbf{97.893}\,\mathrm{tonf/m^{2}}}

\begin{longtable}{@{}l p{0.52\linewidth} r l l@{}}
\toprule
\textbf{Dato} & \textbf{Descripción} & \textbf{Valor} & \textbf{Unidad} & \textbf{Ref.}\\
\midrule
$\varphi_{b,S}$ & Factor de resistencia - Servicio I & \cellcolor{fondodato}0.6500 & $$ & {\scriptsize } \\
\bottomrule
\end{longtable}

\eqitem{Capacidad mayorada - Servicio I}{}{q_{R,S} = \varphi_{b,S} \cdot q_{\mathrm{n}} = \mathbf{79.538}\,\mathrm{tonf/m^{2}}}

\eqitem{Xo - Servicio I}{}{X_{\mathrm{o,S}} = \dfrac{\mathit{MV}_{\mathrm{S}} - \mathit{MH}_{\mathrm{S}}}{V_{\mathrm{S}}} = \mathbf{1.839}\,\mathrm{m}}

\eqitem{Excentricidad - Servicio I}{}{e_{\mathrm{S}} = \dfrac{B}{2} - X_{\mathrm{o,S}} = \mathbf{0.6105}\,\mathrm{m}}

\begin{center}
{\small\textbf{Esfuerzo último actuante sobre suelo no rocoso ($\sigma$v = V/(B - 2e))}}\\[3pt]
\adjustbox{max width=\linewidth}{\begin{tabular}{cccccc}
\toprule
\begin{tabular}[b]{@{}c@{}}\textbf{Estado}\end{tabular} & \begin{tabular}[b]{@{}c@{}}\textbf{$\Sigma$V}\\ {\scriptsize $\mathrm{tonf/m}$}\end{tabular} & \begin{tabular}[b]{@{}c@{}}\textbf{e}\\ {\scriptsize $\mathrm{m}$}\end{tabular} & \begin{tabular}[b]{@{}c@{}}\textbf{$\sigma$v}\\ {\scriptsize $\mathrm{tonf/m^{2}}$}\end{tabular} & \begin{tabular}[b]{@{}c@{}}\textbf{qR}\\ {\scriptsize $\mathrm{tonf/m^{2}}$}\end{tabular} & \begin{tabular}[b]{@{}c@{}}\textbf{Cumple}\end{tabular} \\
\midrule
Resistencia I - Máx. & $113.26$ & $0.7573$ & $33.456$ & $55.065$ & \ok{OK} \\
Resistencia I - Mín. & $68.699$ & $0.2024$ & $15.283$ & $55.065$ & \ok{OK} \\
Evento Extremo I - Máx. & $100.24$ & $1.124$ & $37.782$ & $97.893$ & \ok{OK} \\
Evento Extremo I - Mín. & $68.699$ & $1.162$ & $26.673$ & $97.893$ & \ok{OK} \\
Servicio I & $83.497$ & $0.6105$ & $22.696$ & $79.538$ & \ok{OK} \\
\bottomrule
\end{tabular}}
\end{center}

\subsubsection{Verificación al desplazamiento [A. 10.6.3.4, T. 10.5.5.2.2-1]}

\nota{RR = $\phi$$\tau$$\cdot$V$\cdot$tan $\delta$ + $\phi$ep$\cdot$Rep ; para hormigón vaciado contra el suelo tan $\delta$ = tan $\phi$f.}

\eqitem{tan $\phi$f}{}{\tan\phi'_f = \tan\left(\phi'_f\right) = \mathbf{0.5774}}

\begin{longtable}{@{}l p{0.52\linewidth} r l l@{}}
\toprule
\textbf{Dato} & \textbf{Descripción} & \textbf{Valor} & \textbf{Unidad} & \textbf{Ref.}\\
\midrule
$\phi_\tau$ & Factor de resistencia al corte suelo-cimentación & \cellcolor{fondodato}0.8000 & $$ & {\scriptsize } \\
$\phi_{ep}$ & Factor de resistencia pasiva & \cellcolor{fondodato}0.5000 & $$ & {\scriptsize } \\
\bottomrule
\end{longtable}

\eqitem{Coeficiente de presión pasiva}{}{K_P = \dfrac{1 + \sin\left(\phi'_f\right)}{1 - \sin\left(\phi'_f\right)} = \mathbf{3}}

\begin{longtable}{@{}l p{0.52\linewidth} r l l@{}}
\toprule
\textbf{Dato} & \textbf{Descripción} & \textbf{Valor} & \textbf{Unidad} & \textbf{Ref.}\\
\midrule
$H$ & Altura de suelo pasivo & \cellcolor{fondodato}0.9000 & $\mathrm{m}$ & {\scriptsize } \\
\bottomrule
\end{longtable}

\eqitem{Resistencia pasiva nominal}{}{R_{\mathrm{ep}} = 0.5 \cdot K_P \cdot \gamma_s \cdot {H}^{2} = \mathbf{2.430}\,\mathrm{tonf/m}}

\begin{center}
{\small\textbf{Verificación del deslizamiento del estribo}}\\[3pt]
\adjustbox{max width=\linewidth}{\begin{tabular}{ccccc}
\toprule
\begin{tabular}[b]{@{}c@{}}\textbf{Estado}\end{tabular} & \begin{tabular}[b]{@{}c@{}}\textbf{V}\\ {\scriptsize $\mathrm{tonf/m}$}\end{tabular} & \begin{tabular}[b]{@{}c@{}}\textbf{H actuante}\\ {\scriptsize $\mathrm{tonf/m}$}\end{tabular} & \begin{tabular}[b]{@{}c@{}}\textbf{RR resistente}\\ {\scriptsize $\mathrm{tonf/m}$}\end{tabular} & \begin{tabular}[b]{@{}c@{}}\textbf{Cumple}\end{tabular} \\
\midrule
Resistencia I - Máx. & $113.26$ & $32.148$ & $53.528$ & \ok{OK} \\
Resistencia I - Mín. & $68.699$ & $14.700$ & $32.946$ & \ok{OK} \\
Evento Extremo I - Máx. & $100.24$ & $46.251$ & $47.512$ & \ok{OK} \\
Evento Extremo I - Mín. & $68.699$ & $34.266$ & $32.946$ & \nook{NO OK} \\
\bottomrule
\end{tabular}}
\end{center}

\subsection{Diseño de la pantalla - armadura interior}

\eqitem{Altura de la pantalla}{}{H_{pantalla} = l_{2} + l_{3} + l_{4} = \mathbf{5.900}\,\mathrm{m}}

\eqitem{Coeficiente de empuje (redondeado, como en el PDF)}{}{k_a = \operatorname{red}\left(K_A\right) = \mathbf{0.3300}}

\eqitem{Presión en la base de la pantalla}{}{\Delta p = k_a \cdot H_{pantalla} \cdot \gamma_s = \mathbf{3.894}\,\mathrm{tonf/m^{2}}}

\begin{center}
{\small\textbf{Cargas horizontales sobre la pantalla (Resistencia I y Evento Extremo I)}}\\[3pt]
\adjustbox{max width=\linewidth}{\begin{tabular}{cccc}
\toprule
\begin{tabular}[b]{@{}c@{}}\textbf{Carga}\end{tabular} & \begin{tabular}[b]{@{}c@{}}\textbf{Fuerza}\\ {\scriptsize $\mathrm{tonf/m}$}\end{tabular} & \begin{tabular}[b]{@{}c@{}}\textbf{Brazo Y}\\ {\scriptsize $\mathrm{m}$}\end{tabular} & \begin{tabular}[b]{@{}c@{}}\textbf{Mx}\\ {\scriptsize $\mathrm{tonf\cdot m/m}$}\end{tabular} \\
\midrule
EH = $\Delta$p$\cdot$H/2 & $11.487$ & $1.967$ & $22.592$ \\
Pseis & $15.436$ & $2.360$ & $36.429$ \\
Hbu & $4.130$ & $4.400$ & $18.172$ \\
BR & $1.570$ & $7.700$ & $12.089$ \\
LSx = H$\cdot$$\Delta$p & $2.336$ & $2.950$ & $6.892$ \\
\bottomrule
\end{tabular}}
\end{center}

\subsubsection{Momento de diseño}

\begin{longtable}{@{}l p{0.52\linewidth} r l l@{}}
\toprule
\textbf{Dato} & \textbf{Descripción} & \textbf{Valor} & \textbf{Unidad} & \textbf{Ref.}\\
\midrule
$\gamma_{EH}$ & Factor de carga EH & \cellcolor{fondodato}1.500 & $$ & {\scriptsize } \\
$\gamma_{EQ}$ & Factor de carga EQ & \cellcolor{fondodato}1 & $$ & {\scriptsize } \\
$\gamma_{LS}$ & Factor de carga LS y BR & \cellcolor{fondodato}1.750 & $$ & {\scriptsize } \\
$\eta_D$ & Factor de ductilidad & \cellcolor{fondodato}1 & $$ & {\scriptsize } \\
$\eta_R$ & Factor de redundancia & \cellcolor{fondodato}1 & $$ & {\scriptsize } \\
$\eta_I$ & Factor de importancia & \cellcolor{fondodato}1 & $$ & {\scriptsize } \\
\bottomrule
\end{longtable}

\eqitem{Modificador de carga}{}{\eta = \eta_D \cdot \eta_R \cdot \eta_I = \mathbf{1}}

\eqitem{Evento Extremo I}{}{M_{u,EV} = \eta \cdot \left(\gamma_{EH} \cdot p_{\mathrm{nMEH}} + \gamma_{EQ} \cdot \left(p_{\mathrm{nMPs}} + p_{\mathrm{nMHu}}\right)\right) = \mathbf{88.488}\,\mathrm{tonf\cdot m/m}}

\eqitem{Resistencia I}{}{M_{u,R} = \eta \cdot \left(\gamma_{EH} \cdot p_{\mathrm{nMEH}} + \gamma_{LS} \cdot \left(p_{\mathrm{nMLS}} + p_{\mathrm{nMBR}}\right)\right) = \mathbf{67.105}\,\mathrm{tonf\cdot m/m}}

\eqitem{Momento último de la pantalla}{}{M_U = \max\left(M_{u,R},\ M_{u,EV}\right) = \mathbf{88.488}\,\mathrm{tonf\cdot m/m}}

\subsubsection{Estimación del refuerzo}

\begin{longtable}{@{}l p{0.52\linewidth} r l l@{}}
\toprule
\textbf{Dato} & \textbf{Descripción} & \textbf{Valor} & \textbf{Unidad} & \textbf{Ref.}\\
\midrule
$b$ & Ancho tributario & \cellcolor{fondodato}1 & $\mathrm{m}$ & {\scriptsize } \\
$d_{hierro}$ & Diámetro del acero de refuerzo & \cellcolor{fondodato}25 & $\mathrm{mm}$ & {\scriptsize } \\
\bottomrule
\end{longtable}

\eqitem{Área de una barra}{}{A_{hierro} = 0.25 \cdot \pi \cdot {d_{hierro}}^{2} = \mathbf{4.909}\,\mathrm{cm^{2}}}

\begin{longtable}{@{}l p{0.52\linewidth} r l l@{}}
\toprule
\textbf{Dato} & \textbf{Descripción} & \textbf{Valor} & \textbf{Unidad} & \textbf{Ref.}\\
\midrule
$rec$ & Recubrimiento & \cellcolor{fondodato}7.500 & $\mathrm{cm}$ & {\scriptsize } \\
\bottomrule
\end{longtable}

\eqitem{Espesor de la pantalla}{}{A_{est} = l_{5} + l_{6} = \mathbf{0.9000}\,\mathrm{m}}

\eqitem{Peralte efectivo}{}{d = A_{est} - \left(rec + 0.5 \cdot d_{hierro}\right) = \mathbf{0.8125}\,\mathrm{m}}

\eqitem{Factor $\beta$1}{}{\beta_{\mathrm{1}} = \begin{cases}0.85 & \text{si } \operatorname{red}\left(\dfrac{f'_{c,est}}{\mathrm{MPa}}\right) \le 28\\ \begin{cases}0.65 & \text{si } \operatorname{red}\left(\dfrac{f'_{c,est}}{\mathrm{MPa}}\right) > 56\\ 0.85 - \dfrac{0.05 \cdot \left(f'_{c,est} - 28\,\mathrm{MPa}\right)}{7\,\mathrm{MPa}} & \text{en otro caso}\end{cases} & \text{en otro caso}\end{cases} = \mathbf{0.8500}}

\begin{longtable}{@{}l p{0.52\linewidth} r l l@{}}
\toprule
\textbf{Dato} & \textbf{Descripción} & \textbf{Valor} & \textbf{Unidad} & \textbf{Ref.}\\
\midrule
$\phi$ & Factor de reducción de resistencia & \cellcolor{fondodato}0.9000 & $$ & {\scriptsize } \\
$j$ & Brazo relativo (expresión aproximada) & \cellcolor{fondodato}0.9000 & $$ & {\scriptsize } \\
\bottomrule
\end{longtable}

\eqitem{Acero requerido}{}{A_{\mathrm{sp}} = \dfrac{M_U}{\phi \cdot f_{\mathrm{y}} \cdot j \cdot d} = \mathbf{31.394}\,\mathrm{cm^{2}/m}}

\eqitem{Profundidad del bloque de compresión}{}{a_{\mathrm{p}} = \dfrac{A_{\mathrm{sp}} \cdot b \cdot f_{\mathrm{y}}}{0.85 \cdot f'_{c,est} \cdot b} = \mathbf{7.387}\,\mathrm{cm}}

\eqitem{Eje neutro}{}{c_{\mathrm{p}} = \dfrac{a_{\mathrm{p}}}{\beta_{\mathrm{1}}} = \mathbf{8.690}\,\mathrm{cm}}

\eqitem{Deformación del acero en tensión}{}{\varepsilon_t = \dfrac{0.003 \cdot \left(d - c_{\mathrm{p}}\right)}{c_{\mathrm{p}}} = \mathbf{0.0250}}

\eqitem{$\varepsilon$t $\ge$ 0.005 (tensión controlada) - pantalla}{}{\varepsilon_t \ge 0.005\;\;\Rightarrow\;\;\text{\ok{OK}}}

\begin{longtable}{@{}l p{0.52\linewidth} r l l@{}}
\toprule
\textbf{Dato} & \textbf{Descripción} & \textbf{Valor} & \textbf{Unidad} & \textbf{Ref.}\\
\midrule
$N$ & Número de barras por metro & \cellcolor{fondodato}7 & $$ & {\scriptsize } \\
\bottomrule
\end{longtable}

\eqitem{Acero total}{}{A_{\mathrm{sT}} = \dfrac{N \cdot A_{hierro}}{b} = \mathbf{34.361}\,\mathrm{cm^{2}/m}}

\eqitem{Acero proporcionado $\ge$ requerido - pantalla}{}{A_{\mathrm{sT}} \ge A_{\mathrm{sp}}\;\;\Rightarrow\;\;\text{\ok{OK}}}

\eqitem{Capacidad de momento}{}{M_{U,V} = \phi \cdot A_{\mathrm{sT}} \cdot f_{\mathrm{y}} \cdot \left(d - 0.5 \cdot a_{\mathrm{p}}\right) = \mathbf{102.72}\,\mathrm{tonf\cdot m/m}}

\eqitem{Capacidad $\ge$ momento último - pantalla}{}{M_{U,V} > M_U\;\;\Rightarrow\;\;\text{\ok{CUMPLE}}}

\eqitem{Separación}{}{S = \dfrac{A_{hierro}}{A_{\mathrm{sT}}} = \mathbf{0.1429}\,\mathrm{m}}

\begin{longtable}{@{}l p{0.52\linewidth} r l l@{}}
\toprule
\textbf{Dato} & \textbf{Descripción} & \textbf{Valor} & \textbf{Unidad} & \textbf{Ref.}\\
\midrule
$S_{asumido}$ & Separación adoptada (emplear $\varnothing$25 mm @ 14 cm) & \cellcolor{fondodato}14 & $\mathrm{cm}$ & {\scriptsize } \\
\bottomrule
\end{longtable}

\subsubsection{Comprobación del As mínimo [Ec. 5.7.3.3.2-1]}

\eqitem{Módulo de la sección}{}{S_{\mathrm{s}} = \dfrac{b \cdot {A_{est}}^{2}}{6} = \mathbf{0.1350}\,\mathrm{m^{3}}}

\eqitem{Módulo de ruptura}{}{f_{\mathrm{r}} = 0.63 \cdot \sqrt{\dfrac{f'_{c,est}}{\mathrm{MPa}}} \cdot \mathrm{MPa} = \mathbf{29.439}\,\mathrm{kgf/cm^{2}}}

\eqitem{1.2 Mcr}{}{\mathit{FM}_{\mathrm{cr}} = \dfrac{1.2 \cdot f_{\mathrm{r}} \cdot S_{\mathrm{s}}}{b} = \mathbf{47.692}\,\mathrm{tonf\cdot m/m}}

\eqitem{1.33 Mu}{}{\mathit{FM}_{\mathrm{u}} = 1.33 \cdot M_U = \mathbf{117.69}\,\mathrm{tonf\cdot m/m}}

\eqitem{Menor valor}{}{\mathit{MU}_{\mathrm{min}} = \min\left(\mathit{FM}_{\mathrm{cr}},\ \mathit{FM}_{\mathrm{u}}\right) = \mathbf{47.692}\,\mathrm{tonf\cdot m/m}}

\eqitem{As mínimo - pantalla}{}{M_{U,V} > \mathit{MU}_{\mathrm{min}}\;\;\Rightarrow\;\;\text{\ok{CUMPLE}}}

\subsubsection{Refuerzo por contracción y temperatura en la pantalla [A. 5.10.8]}

\eqitem{Acero de temperatura}{}{A_{s,temp} = \dfrac{0.75 \cdot A_{est} \cdot l_{7} \cdot \mathrm{MPa}}{2 \cdot \left(A_{est} + l_{7}\right) \cdot f_{\mathrm{y}}} = \mathbf{0.6671}\,\mathrm{mm^{2}/mm}}

\eqitem{0.233 $\le$ As $\le$ 1.27 mm$^{2}$/mm - pantalla}{}{A_{s,temp} \ge \dfrac{0.233 \cdot {\mathrm{mm}}^{2}}{\mathrm{mm}} \;\wedge\; A_{s,temp} \le \dfrac{1.27 \cdot {\mathrm{mm}}^{2}}{\mathrm{mm}}\;\;\Rightarrow\;\;\text{\ok{OK}}}

\begin{longtable}{@{}l p{0.52\linewidth} r l l@{}}
\toprule
\textbf{Dato} & \textbf{Descripción} & \textbf{Valor} & \textbf{Unidad} & \textbf{Ref.}\\
\midrule
$d_{hierro,t}$ & Diámetro del refuerzo & \cellcolor{fondodato}16 & $\mathrm{mm}$ & {\scriptsize } \\
\bottomrule
\end{longtable}

\eqitem{Área del refuerzo}{}{A_{hierro,t} = \dfrac{\pi \cdot {d_{hierro,t}}^{2}}{4} = \mathbf{2.011}\,\mathrm{cm^{2}}}

\eqitem{Separación requerida}{}{S_{a,p,t} = \dfrac{A_{hierro,t}}{A_{s,temp}} = \mathbf{0.3014}\,\mathrm{m}}

\eqitem{Separación máxima 3$\cdot$espesor}{}{S_{\mathrm{max1}} = 3 \cdot A_{est} = \mathbf{270}\,\mathrm{cm}}

\begin{longtable}{@{}l p{0.52\linewidth} r l l@{}}
\toprule
\textbf{Dato} & \textbf{Descripción} & \textbf{Valor} & \textbf{Unidad} & \textbf{Ref.}\\
\midrule
$S_{\mathrm{max2}}$ & Separación máxima 45 cm & \cellcolor{fondodato}45 & $\mathrm{cm}$ & {\scriptsize } \\
\bottomrule
\end{longtable}

\nota{Emplear $\varnothing$16 mm @ 25 cm.}

\subsection{Diseño de la armadura de la zapata - punta (Resistencia I)}

\eqitem{Esfuerzo máximo en el suelo}{}{\sigma_{R,max} = s_{\mathrm{ig,RMax}} = \mathbf{33.456}\,\mathrm{tonf/m^{2}}}

\eqitem{Momento en la punta}{}{M_{p,z} = \dfrac{\sigma_{R,max} \cdot {l_{8}}^{2}}{2} = \mathbf{24.088}\,\mathrm{tonf\cdot m/m}}

\eqitem{Módulo de la sección}{}{S_s = 0.1667 \cdot b \cdot {l_{1}}^{2} = \mathbf{0.2017}\,\mathrm{m^{3}}}

\eqitem{1.2 Mcr}{}{FMcr_{p,z} = \dfrac{1.2 \cdot f_{\mathrm{r}} \cdot S_s}{b} = \mathbf{71.258}\,\mathrm{tonf\cdot m/m}}

\eqitem{1.33 Mu}{}{FMu_{p,z} = 1.33 \cdot M_{p,z} = \mathbf{32.037}\,\mathrm{tonf\cdot m/m}}

\eqitem{Menor valor}{}{MU_{min,p,z} = \min\left(FMcr_{p,z},\ FMu_{p,z}\right) = \mathbf{32.037}\,\mathrm{tonf\cdot m/m}}

\eqitem{Mu > MUmin (si no, se diseña con MUmin)}{}{M_{p,z} > MU_{min,p,z}\;\;\Rightarrow\;\;\text{\info{NO CUMPLE}}}

\eqitem{Momento de diseño de la punta}{}{M_{u,p,z} = \max\left(M_{p,z},\ MU_{min,p,z}\right) = \mathbf{32.037}\,\mathrm{tonf\cdot m/m}}

\begin{longtable}{@{}l p{0.52\linewidth} r l l@{}}
\toprule
\textbf{Dato} & \textbf{Descripción} & \textbf{Valor} & \textbf{Unidad} & \textbf{Ref.}\\
\midrule
$d_{hierro,p,z}$ & Diámetro del refuerzo & \cellcolor{fondodato}20 & $\mathrm{mm}$ & {\scriptsize } \\
\bottomrule
\end{longtable}

\eqitem{Área de una barra}{}{A_{hierro,p,z} = 0.25 \cdot \pi \cdot {d_{hierro,p,z}}^{2} = \mathbf{3.142}\,\mathrm{cm^{2}}}

\eqitem{Peralte efectivo}{}{d_{p,e} = l_{1} - \left(rec + 0.5 \cdot d_{hierro,p,z}\right) = \mathbf{101.50}\,\mathrm{cm}}

\eqitem{Acero requerido}{}{A_{s,p,z} = \dfrac{M_{u,p,z}}{\phi \cdot f_{\mathrm{y}} \cdot j \cdot d_{p,e}} = \mathbf{9.099}\,\mathrm{cm^{2}/m}}

\eqitem{Bloque de compresión}{}{a_{p,z} = \dfrac{A_{s,p,z} \cdot b \cdot f_{\mathrm{y}}}{0.85 \cdot f'_{c,est} \cdot b} = \mathbf{2.141}\,\mathrm{cm}}

\eqitem{Eje neutro}{}{c_{p,z} = \dfrac{a_{p,z}}{\beta_{\mathrm{1}}} = \mathbf{2.519}\,\mathrm{cm}}

\eqitem{Deformación del acero}{}{\varepsilon_{t,p,z} = \dfrac{0.003 \cdot \left(d_{p,e} - c_{p,z}\right)}{c_{p,z}} = \mathbf{0.1179}}

\eqitem{$\varepsilon$t $\ge$ 0.005 - punta}{}{\varepsilon_{t,p,z} \ge 0.005\;\;\Rightarrow\;\;\text{\ok{OK}}}

\begin{longtable}{@{}l p{0.52\linewidth} r l l@{}}
\toprule
\textbf{Dato} & \textbf{Descripción} & \textbf{Valor} & \textbf{Unidad} & \textbf{Ref.}\\
\midrule
$N_{p,z}$ & Número de barras por metro & \cellcolor{fondodato}3 & $$ & {\scriptsize } \\
\bottomrule
\end{longtable}

\eqitem{Acero total}{}{AsT_{p,z} = \dfrac{N_{p,z} \cdot A_{hierro,p,z}}{b} = \mathbf{9.425}\,\mathrm{cm^{2}/m}}

\eqitem{Acero proporcionado $\ge$ requerido - punta}{}{AsT_{p,z} \ge A_{s,p,z}\;\;\Rightarrow\;\;\text{\ok{CUMPLE}}}

\eqitem{Capacidad de momento}{}{MU_{V,p,z} = \phi \cdot AsT_{p,z} \cdot f_{\mathrm{y}} \cdot \left(d_{p,e} - 0.5 \cdot a_{p,z}\right) = \mathbf{36.484}\,\mathrm{tonf\cdot m/m}}

\eqitem{Capacidad de momento - punta}{}{MU_{V,p,z} > MU_{min,p,z}\;\;\Rightarrow\;\;\text{\ok{CUMPLE}}}

\eqitem{Separación requerida}{}{S_{p,z} = \dfrac{A_{hierro,p,z}}{A_{s,p,z}} = \mathbf{0.3453}\,\mathrm{m}}

\begin{longtable}{@{}l p{0.52\linewidth} r l l@{}}
\toprule
\textbf{Dato} & \textbf{Descripción} & \textbf{Valor} & \textbf{Unidad} & \textbf{Ref.}\\
\midrule
$S_{asumido,p,z}$ & Separación adoptada (emplear $\varnothing$20 mm @ 35 cm) & \cellcolor{fondodato}35 & $\mathrm{cm}$ & {\scriptsize } \\
\bottomrule
\end{longtable}

\subsubsection{Verificación del esfuerzo cortante [A. 5.7.3.3, A. 5.7.3.4.1]}

\eqitem{Cortante en la unión zapata-pantalla}{}{V_{\mathrm{u}} = \sigma_{R,max} \cdot l_{8} \cdot b = \mathbf{40.147}\,\mathrm{tonf}}

\eqitem{Peralte de cortante}{}{d_{\mathrm{v}} = d_{p,e} = \mathbf{1.015}\,\mathrm{m}}

\eqitem{3 dv > l8}{}{3 \cdot d_{\mathrm{v}} > l_{8}\;\;\Rightarrow\;\;\text{\ok{OK}}}

\begin{longtable}{@{}l p{0.52\linewidth} r l l@{}}
\toprule
\textbf{Dato} & \textbf{Descripción} & \textbf{Valor} & \textbf{Unidad} & \textbf{Ref.}\\
\midrule
$\beta_v$ & Factor $\beta$ & \cellcolor{fondodato}2 & $$ & {\scriptsize } \\
$\theta_v$ & Ángulo $\theta$ & \cellcolor{fondodato}45 & $^{\circ}$ & {\scriptsize } \\
\bottomrule
\end{longtable}

\eqitem{Resistencia nominal del hormigón}{}{V_{\mathrm{c}} = 0.083 \cdot \beta_v \cdot \sqrt{\dfrac{f'_{c,est}}{\mathrm{MPa}}} \cdot \mathrm{MPa} \cdot b \cdot d_{\mathrm{v}} = \mathbf{78.734}\,\mathrm{tonf}}

\eqitem{Resistencia última del hormigón}{}{V_{u,c} = \phi \cdot V_{\mathrm{c}} = \mathbf{70.861}\,\mathrm{tonf}}

\eqitem{Cortante en la punta}{}{V_{u,c} > V_{\mathrm{u}}\;\;\Rightarrow\;\;\text{\ok{OK}}}

\subsection{Diseño del talón de la zapata}

\eqitem{Peso de la zapata}{}{DC_{zapata} = e_{\mathrm{DC,1}} = \mathbf{12.936}\,\mathrm{tonf/m}}

\eqitem{Esfuerzo mínimo en el suelo}{}{\sigma_{R,min} = s_{\mathrm{ig,RMin}} = \mathbf{15.283}\,\mathrm{tonf/m^{2}}}

\eqitem{Brazo del relleno}{}{b_{EV} = 0.5 \cdot l_{9} = \mathbf{1.400}\,\mathrm{m}}

\eqitem{Brazo del peso de la zapata}{}{b_{DC,z} = 0.5 \cdot l_{9} = \mathbf{1.400}\,\mathrm{m}}

\eqitem{Excentricidad}{}{e_{\sigma,min} = e_{\mathrm{RMin}} = \mathbf{0.2024}\,\mathrm{m}}

\eqitem{B - 2e}{}{B_{\mathrm{2e}} = B - 2 \cdot e_{\sigma,min} = \mathbf{4.495}\,\mathrm{m}}

\eqitem{Longitud cargada bajo el talón}{}{b_{\sigma} = B_{\mathrm{2e}} - \left(l_{8} + l_{6} + l_{5}\right) = \mathbf{2.395}\,\mathrm{m}}

\begin{longtable}{@{}l p{0.52\linewidth} r l l@{}}
\toprule
\textbf{Dato} & \textbf{Descripción} & \textbf{Valor} & \textbf{Unidad} & \textbf{Ref.}\\
\midrule
$\gamma_{EV}$ & Factor de carga EV & \cellcolor{fondodato}1.350 & $$ & {\scriptsize } \\
$\gamma_{DC,z}$ & Factor de carga DC & \cellcolor{fondodato}1.250 & $$ & {\scriptsize } \\
\bottomrule
\end{longtable}

\eqitem{Momento último en el talón}{}{M_{u,t,z} = \eta \cdot \left(\gamma_{EV} \cdot \mathit{EV} \cdot b_{EV} + \gamma_{DC,z} \cdot DC_{zapata} \cdot b_{DC,z} - \dfrac{\sigma_{R,min} \cdot {b_{\sigma}}^{2}}{2}\right) = \mathbf{41.247}\,\mathrm{tonf\cdot m/m}}

\eqitem{1.2 Mcr}{}{FMcr_{t,z} = \dfrac{1.2 \cdot f_{\mathrm{r}} \cdot S_s}{b} = \mathbf{71.258}\,\mathrm{tonf\cdot m/m}}

\eqitem{1.33 Mu}{}{FMu_{t,z} = 1.33 \cdot M_{u,t,z} = \mathbf{54.858}\,\mathrm{tonf\cdot m/m}}

\eqitem{Menor valor}{}{MU_{min,t,z} = \min\left(FMcr_{t,z},\ FMu_{t,z}\right) = \mathbf{54.858}\,\mathrm{tonf\cdot m/m}}

\eqitem{Mu > MUmin (si no, se diseña con MUmin)}{}{M_{u,t,z} > MU_{min,t,z}\;\;\Rightarrow\;\;\text{\info{NO CUMPLE}}}

\eqitem{Momento de diseño del talón}{}{M_{u,t,z,d} = \max\left(M_{u,t,z},\ MU_{min,t,z}\right) = \mathbf{54.858}\,\mathrm{tonf\cdot m/m}}

\begin{longtable}{@{}l p{0.52\linewidth} r l l@{}}
\toprule
\textbf{Dato} & \textbf{Descripción} & \textbf{Valor} & \textbf{Unidad} & \textbf{Ref.}\\
\midrule
$d_{hierro,t,z}$ & Diámetro del refuerzo & \cellcolor{fondodato}20 & $\mathrm{mm}$ & {\scriptsize } \\
\bottomrule
\end{longtable}

\eqitem{Área de una barra}{}{A_{hierro,t,z} = 0.25 \cdot \pi \cdot {d_{hierro,t,z}}^{2} = \mathbf{3.142}\,\mathrm{cm^{2}}}

\eqitem{Peralte efectivo}{}{d_{t,z} = l_{1} - \left(rec + 0.5 \cdot d_{hierro,t,z}\right) = \mathbf{1.015}\,\mathrm{m}}

\eqitem{Acero requerido}{}{A_{s,t,z} = \dfrac{M_{u,t,z,d}}{\phi \cdot f_{\mathrm{y}} \cdot j \cdot d_{t,z}} = \mathbf{15.580}\,\mathrm{cm^{2}/m}}

\eqitem{Bloque de compresión}{}{a_{t,z} = \dfrac{A_{s,t,z} \cdot b \cdot f_{\mathrm{y}}}{0.85 \cdot f'_{c,est} \cdot b} = \mathbf{3.666}\,\mathrm{cm}}

\eqitem{Eje neutro}{}{c_{t,z} = \dfrac{a_{t,z}}{\beta_{\mathrm{1}}} = \mathbf{4.313}\,\mathrm{cm}}

\eqitem{Deformación del acero}{}{\varepsilon_{t,t,z} = \dfrac{0.003 \cdot \left(d_{t,z} - c_{t,z}\right)}{c_{t,z}} = \mathbf{0.0676}}

\eqitem{$\varepsilon$t $\ge$ 0.005 - talón}{}{\varepsilon_{t,t,z} \ge 0.005\;\;\Rightarrow\;\;\text{\ok{OK}}}

\begin{longtable}{@{}l p{0.52\linewidth} r l l@{}}
\toprule
\textbf{Dato} & \textbf{Descripción} & \textbf{Valor} & \textbf{Unidad} & \textbf{Ref.}\\
\midrule
$N_{t,z}$ & Número de barras por metro & \cellcolor{fondodato}5 & $$ & {\scriptsize } \\
\bottomrule
\end{longtable}

\eqitem{Acero total}{}{AsT_{t,z} = \dfrac{N_{t,z} \cdot A_{hierro,t,z}}{b} = \mathbf{15.708}\,\mathrm{cm^{2}/m}}

\eqitem{Acero proporcionado $\ge$ requerido - talón}{}{AsT_{t,z} \ge A_{s,t,z}\;\;\Rightarrow\;\;\text{\ok{CUMPLE}}}

\eqitem{Capacidad de momento}{}{MU_{V,t,z} = \phi \cdot AsT_{t,z} \cdot f_{\mathrm{y}} \cdot \left(d_{t,z} - 0.5 \cdot a_{t,z}\right) = \mathbf{60.345}\,\mathrm{tonf\cdot m/m}}

\eqitem{Capacidad de momento - talón}{}{MU_{V,t,z} > MU_{min,t,z}\;\;\Rightarrow\;\;\text{\ok{CUMPLE}}}

\eqitem{Separación}{}{S_{t,z} = \dfrac{A_{hierro,t,z}}{AsT_{t,z}} = \mathbf{0.2000}\,\mathrm{m}}

\begin{longtable}{@{}l p{0.52\linewidth} r l l@{}}
\toprule
\textbf{Dato} & \textbf{Descripción} & \textbf{Valor} & \textbf{Unidad} & \textbf{Ref.}\\
\midrule
$S_{asumido,t,z}$ & Separación adoptada (emplear $\varnothing$20 mm @ 20 cm) & \cellcolor{fondodato}20 & $\mathrm{cm}$ & {\scriptsize } \\
\bottomrule
\end{longtable}

\subsubsection{Refuerzo por contracción y temperatura en la zapata [A. 5.10.8]}

\eqitem{Acero de temperatura}{}{A_{s,t,z} = \dfrac{0.75 \cdot B \cdot l_{1} \cdot \mathrm{MPa}}{2 \cdot \left(B + l_{1}\right) \cdot f_{\mathrm{y}}} = \mathbf{0.8021}\,\mathrm{mm^{2}/mm}}

\eqitem{0.233 $\le$ As $\le$ 1.27 mm$^{2}$/mm - zapata}{}{A_{s,t,z} \ge \dfrac{0.233 \cdot {\mathrm{mm}}^{2}}{\mathrm{mm}} \;\wedge\; A_{s,t,z} \le \dfrac{1.27 \cdot {\mathrm{mm}}^{2}}{\mathrm{mm}}\;\;\Rightarrow\;\;\text{\ok{OK}}}

\begin{longtable}{@{}l p{0.52\linewidth} r l l@{}}
\toprule
\textbf{Dato} & \textbf{Descripción} & \textbf{Valor} & \textbf{Unidad} & \textbf{Ref.}\\
\midrule
$d_{h,t,z}$ & Diámetro del refuerzo & \cellcolor{fondodato}16 & $\mathrm{mm}$ & {\scriptsize } \\
\bottomrule
\end{longtable}

\eqitem{Área del refuerzo}{}{A_{h,t,z} = 0.25 \cdot \pi \cdot {d_{h,t,z}}^{2} = \mathbf{2.011}\,\mathrm{cm^{2}}}

\eqitem{Separación requerida}{}{S_{a,p,t,z} = \dfrac{A_{h,t,z}}{A_{s,t,z}} = \mathbf{25.067}\,\mathrm{cm}}

\eqitem{Separación máxima 3$\cdot$espesor}{}{Smax_{1,t,z} = 3 \cdot l_{1} = \mathbf{330}\,\mathrm{cm}}

\begin{longtable}{@{}l p{0.52\linewidth} r l l@{}}
\toprule
\textbf{Dato} & \textbf{Descripción} & \textbf{Valor} & \textbf{Unidad} & \textbf{Ref.}\\
\midrule
$Smax_{2,t,z}$ & Separación máxima 45 cm & \cellcolor{fondodato}45 & $\mathrm{cm}$ & {\scriptsize } \\
\bottomrule
\end{longtable}

\nota{Emplear $\varnothing$16 mm @ 25 cm.}

\clearpage
\end{document}
